\documentclass[11pt,a4paper]{article}
\usepackage[margin=23mm,headheight=16pt]{geometry}
\usepackage{fontspec}
\setmainfont{NotoSerifCJK-Regular.ttc}[Path=/usr/share/fonts/opentype/noto/,FontIndex=2,BoldFont=NotoSerifCJK-Bold.ttc,ItalicFont=NotoSerifCJK-Regular.ttc,ItalicFeatures={FakeSlant=0.2}]
\setsansfont{NotoSansCJK-Regular.ttc}[Path=/usr/share/fonts/opentype/noto/,FontIndex=2,BoldFont=NotoSansCJK-Bold.ttc]
\XeTeXlinebreaklocale "zh"
\XeTeXlinebreakskip = 0pt plus 1pt
\usepackage{amsmath,amssymb,booktabs,longtable,array,tabularx,enumitem,graphicx,xcolor,tikz,fancyhdr}
\usetikzlibrary{arrows.meta,positioning,calc,shapes.geometric,fit,backgrounds,decorations.pathreplacing,patterns}
\usepackage{hyperref}
\definecolor{inkblue}{HTML}{173B58}
\hypersetup{colorlinks=true,linkcolor=inkblue,urlcolor=inkblue,pdftitle={高级图像处理：零基础详解},pdfauthor={}}
\setlength{\parindent}{2em}
\setlength{\parskip}{5pt plus 1pt}
\setlength{\emergencystretch}{3em}
\linespread{1.18}
\renewcommand{\contentsname}{目录}
\renewcommand{\figurename}{图}
\renewcommand{\tablename}{表}
\pagestyle{fancy}\fancyhf{}\fancyhead[L]{\small\sffamily 高级图像处理}\fancyhead[R]{\small\sffamily 零基础详解}\fancyfoot[C]{\thepage}

\setlist{itemsep=3pt,topsep=4pt,leftmargin=2em}
\renewcommand{\arraystretch}{1.25}

\begin{document}
\begin{titlepage}
\vspace*{20mm}
{\Huge\sffamily 高级图像处理\par}\vspace{8mm}
{\LARGE 零基础详解\par}\vspace{12mm}
{\large 从像素、灰度到滤波与频谱\par}
\vfill
\noindent 以提供的《小测复习讲义（黑白打印版）》为范围，补足理解所需的基础概念与逐步解题过程。原题编号保留，公式与图示重新编排；原稿中的已核实错误在对应讲解中更正。\par
\vspace{6mm}
\noindent 先理解图像中的变化，再把变化写成数字和公式。数学符号在使用前说明，计算例题写明单位、代入过程与结论条件。\par
\vspace{10mm}
2026年10月19日
\end{titlepage}
\setcounter{tocdepth}{2}\tableofcontents\clearpage
\section{开始之前：把图像问题读成数字问题}
\subsection{本书的学习顺序}
一张灰度照片可以看作一张数字表。表中的每个数表示一个小位置的亮暗；把这些位置按原来的排列显示出来，便成为图像。图像处理的大部分入门操作，都可以从这个事实出发理解：改变一个数，可以改变一个像素的亮暗；把周围几个数平均，可以减小局部起伏；比较相邻位置的数，可以发现边缘。

学习过程分为四步。第一步理解图像是怎样记录位置与亮暗的，以及人眼怎样观察这些变化。第二步研究对单个灰度值作变换，并统计整张图的灰度分布。第三步让相邻像素参与计算，理解平滑、去噪和锐化。第四步把图像理解为不同快慢变化的组合，进入频率域滤波。最后一种表示并没有创造另一张与原图无关的图，而是把同一份信息换一种方式描述。

本书沿用附件中的题号，例如E1、D4，使原题与详解可以相互定位。但知识讲解与对应例题放在一起，不要求先读完全部概念再集中做题。题目中的条件、计算结果和适用范围必须一起理解。原稿中的“必考”“高频考点”等标签不是本书对考试命题的保证。

\subsection{符号只是名称：先明确输入与输出}
字母常用来给一个数或一组数起名字。例如，用$r$表示处理前的灰度，用$s$表示处理后的灰度。若规则是“在原数上加20”，就写成
\[
s=r+20.
\]
当$r=30$时，$s=30+20=50$；当$r=100$时，$s=120$。字母没有增加新的计算，只是避免为每个可能的输入重新写一遍同样的话。

函数（Function）是从输入得到输出的规则。把上述规则叫作$T$，可以写成$s=T(r)$，其中$T(r)=r+20$。这里的括号表示“将$r$交给规则$T$处理”，不是相乘。图像处理中常见的$f(x,y)$也遵循这个习惯：$x,y$指明位置，$f(x,y)$给出该位置的灰度值。它不表示$f$乘$x$再乘$y$。

下标（Subscript）用于区分同类对象。$r_0,r_1,r_2$表示第0、第1、第2个灰度值；$n_k$表示第$k$类的数量。下标的作用像座位编号，它本身不是乘法。上标有时表示幂，例如$r^2=r\times r$，有时是标记；具体含义必须看定义。

\subsection{分数、平均与比例}
分数线表示除法。$\frac{a}{b}$表示$a\div b$，前提是$b\ne0$。如果5个像素的灰度分别是10、20、20、30、40，它们的总和是120，平均值为
\[
\frac{10+20+20+30+40}{5}=\frac{120}{5}=24.
\]
平均值回答的是“如果把总量均匀分给这5个位置，每个位置得到多少”。它不要求24恰好是原来5个数之一。

比例回答的是“某一部分占整体多少”。若100个像素中有20个灰度为0，灰度0的比例为$20/100=0.2=20\%$。后文把这种比例称为相对频数；在图像的经验统计模型中，也把它当作该灰度的概率（Probability）。这里所需的概率概念首先就是计数与除法，不必先学习复杂的随机过程。

\subsection{求和符号不是一种新运算}
希腊大写字母$\Sigma$表示把一串数相加。例如
\[
\sum_{i=1}^{4}a_i=a_1+a_2+a_3+a_4.
\]
符号下面的$i=1$给出起点，上面的4给出终点。$i$是依次取1、2、3、4的编号。若$(a_1,a_2,a_3,a_4)=(2,3,5,6)$，则结果为16。

二维图像既有行又有列，因此有时出现两层求和。例如一个$2\times2$小区域的4个数为
\[
\begin{bmatrix}1&2\\3&4\end{bmatrix}.
\]
先把每一行相加，得到3和7，再把行的结果相加，得到10。双重求和符号只是把这个“逐行、逐列遍历”的过程缩写起来。后文的卷积与相关公式，都是先给出实际小网格算例，再写这种缩写形式。

\subsection{平方、平方根与指数}
平方是同一个数乘自己：$3^2=3\times3=9$。平方根（Square Root）是反过来寻找非负的那个数：$\sqrt9=3$。例如直角三角形的两条直角边分别为3和4，斜边长度是$\sqrt{3^2+4^2}=5$。频谱图中某点离中心有多远，也用同一种距离计算。

指数（Exponent）表示乘法的重复或其连续扩展。$2^3=2\times2\times2=8$，$2^{-1}=1/2$，$r^{1/2}=\sqrt r$。伽马变换中的$r^\gamma$就是这种幂运算；判断它如何改变亮暗，可以先代入$r=0.25$：$0.25^{1/2}=0.5$，而$0.25^2=0.0625$。在0至1的灰度尺度上，前者变亮，后者变暗。

自然指数函数写作$\exp(x)=e^x$，其中$e\approx2.71828$。对数（Logarithm）是指数运算的逆运算：在$a>0$、$a\ne1$、$c>0$的条件下，若$a^b=c$，就有$\log_a c=b$。自然对数记为$\ln c$，以$e$为底。特别重要的一条规则是
\[
\ln(ab)=\ln a+\ln b,\qquad a>0,\ b>0.
\]
这解释了后面同态滤波为什么先取对数：它希望把“照明乘反射”变成“两项相加”。在使用这条公式前，必须确保要取对数的数为正。

\subsection{单位是计算的一部分}
一个数后面的单位说明它在量什么。毫米描述长度，像素描述离散采样点，毫米/像素描述每个采样点对应的实际长度。若10毫米宽的纸面由20个像素记录，则
\[
\frac{10\ \mathrm{mm}}{20\ \mathrm{pixel}}=0.5\ \mathrm{mm/pixel}.
\]
长度单位必须统一后再相加或比较。常用换算为
\[
1\ \mathrm{m}=1000\ \mathrm{mm},\qquad
1\ \mathrm{mm}=1000\ \mu\mathrm{m}.
\]
例如$4.52\,\mu\mathrm{m}=0.00452\,\mathrm{mm}$。如果把4.52直接当作毫米代入，相机视场会被放大1000倍；这种错误不是算法问题，而是量纲不一致。

分辨率计算常需要“向上取整”。若计算出至少333.33个像素，实际不能制造三分之一个像素，也不能取333而低于下限，所以至少取334。向上取整记为$\lceil333.33\rceil=334$。普通四舍五入与向上取整不是同一种操作，必须按题意选择。

\subsection{示意图、精确计算与现实测量}
示意图用于突出关系，例如两个像素是否相邻、一条变换曲线位于恒等线的哪一侧。它不一定按真实尺寸比例绘制。精确计算则必须明确输入数据、公式和假设。现实测量还受到噪声、模糊、几何误差等因素影响，因此“计算得到足够像素”不自动等于“实际一定能识别目标”。

后文会区分这三个层次。需要理解机制时使用小整数或小网格；求数值答案时给出具体单位和逐步代入；涉及设备选择时再说明理想模型的边界。这样可以先建立清楚的基本关系，再逐渐加入真实场景的复杂因素。

% 基础篇：原稿模块一至三及 E2、E4、E5、F2--F7、G1、G2。
\section{从照片到可以计算的数字图像}
\label{sec:foundations-image}

一张手机照片看上去有天空、道路和行人，计算机读到的却是一份按位置排列的数值。学习图像处理的第一步，是把这两种看法联系起来：数值记录某个位置的明暗或颜色，数值之间的关系记录物体的形状。处理图像，就是在保留任务所需信息的前提下，有目的地改变或解释这些数值。

本章从方格和数表开始，不要求预先学过信号处理。后续章节中反复出现的“逐像素”“邻域”“灰度”“分辨率”，都可以回到这里找到具体含义。

\subsection{一幅图像为什么可以写成一张数表}

把一张黑白照片盖上透明方格纸，每格用一个数代表明暗，就得到数字图像（digital image）的直观模型。方格中的一个记录单位叫像素（pixel）。像素不是被拍物体上的一个天然小方块，而是数字图像中的一个位置及其数值；把照片放大后看到的方块，是显示软件画出像素的一种方式。

对只有明暗、没有颜色的灰度图像（grayscale image），每个像素用一个灰度级（gray level）表示。常见的 8 位图像把灰度编号为 0 到 255：0 通常显示为黑，255 通常显示为白，中间的数显示为不同深浅的灰。这里的“通常”很重要：编号如何变成屏幕亮度，还要经过显示设备和显示设置；灰度 200 并不直接等于某个固定物理亮度。

\begin{figure}[htbp]
\centering
\begin{tikzpicture}[x=.78cm,y=.78cm,font=\small]
% Display white for 255 and black for 0, explicitly overpaint correct gray mapping.
\foreach \x/\y/\v in {0/2/0,1/2/60,2/2/180,3/2/255,0/1/20,1/1/90,2/1/200,3/1/240,0/0/50,1/0/120,2/0/220,3/0/255}{
  \pgfmathsetmacro{\shade}{100-\v/2.55}
  \fill[black!\shade] (\x,\y) rectangle ++(1,1);
}
\draw[step=1,gray] (0,0) grid (4,3);
\draw[-{Latex}] (0,3.5)--(4,3.5) node[right]{列 $j$};
\draw[-{Latex}] (-.45,3)--(-.45,0) node[below]{行 $i$};
\foreach \x in {0,...,3} \node at (\x+.5,3.25) {\x};
\foreach \y/\v in {2/0,1/1,0/2} \node at (-.2,\y+.5) {\v};
\draw[line width=1.2pt] (1,1) rectangle (2,2);
\draw[-{Latex}] (4.4,1.5)--(5.3,1.5);
\node[anchor=west] at (5.6,1.5) {$\displaystyle f=\begin{bmatrix}0&60&180&255\\20&90&200&240\\50&120&220&255\end{bmatrix}$};
\node at (2,-.65) {3 行、4 列，共 12 个像素};
\node[anchor=west] at (5.6,-.05) {$f(1,1)=90$：第 2 行、第 2 列};
\end{tikzpicture}
\caption{小灰度图与数表的一一对应。这里从 0 开始编号，因此编号 1 指第二行或第二列。}
\label{fig:foundation-grid}
\end{figure}

这种数表在数学中叫矩阵（matrix）。记图像为 $f(i,j)$，意思是“到第 $i$ 行、第 $j$ 列去取一个灰度数”。字母 $f$ 是图像的名字，括号里的 $i,j$ 是地址，不是乘法。若图像有 $M$ 行、$N$ 列，并从 0 编号，则
\[
  i=0,1,\ldots,M-1,\qquad j=0,1,\ldots,N-1.
\]
例如 $M=3,N=4$，共有 $MN=3\times4=12$ 个像素。教材也常写连续坐标 $f(x,y)$，其中 $x$ 向右、$y$ 向上；图像数组常把行号向下增加。两者都可以使用，但计算位置或距离前必须先说明约定。本基础篇的离散数表统一使用“行 $i$、列 $j$”。

\subsection{二值 灰度和彩色三种记录方式}

二值图像（binary image）只使用两个值，通常记为 0 和 1。例如，把车牌文字记为 1、背景记为 0，就得到“哪里是文字”的地图。这里的 1 表示被选中的目标，不一定必须显示成白色。判断像素是不是目标，与目标显示成什么颜色，是两个不同问题。

彩色图像（color image）通常为每个像素记录三个通道（channel），例如红绿蓝颜色模型（red--green--blue color model, RGB）中的 $R,G,B$。像素 $(255,0,0)$ 可显示为红色，$(0,255,0)$ 可显示为绿色，$(80,80,80)$ 可显示为灰色。一个三通道 RGB 图像可以看成三张大小相同的灰度数表叠在一起。三通道不是“一个像素只能有三种颜色”，而是用三个强度组合表达大量颜色。

记录一个二选一结果需要一个比特（bit）。每增加一个比特，可用组合数就翻倍。因此，$b$ 比特可以表示
\[
  L=2^b
\]
个不同编号。式中 $b$ 是每个灰度像素的位数，$L$ 是可用灰度级数。例如 $b=1$ 时 $L=2$，$b=8$ 时 $L=256$，$b=10$ 时 $L=1024$。8 比特又叫 1 字节（byte）。一张 $1000\times1000$、每像素 8 比特的单通道图像，未压缩像素数据为
\[
  1000\times1000\times8=8\,000\,000\ \text{bit}
  =1\,000\,000\ \text{byte}.
\]
同样尺寸、每个通道 8 比特的 RGB 图像，像素数据是它的三倍。实际文件还可能包含拍摄信息并经过压缩，不能仅凭文件大小反推像素数。

\subsection{图像处理要完成什么任务}

想看清一张昏暗的门禁照片时，输入是图像，输出通常仍是图像：把暗部提亮、减弱噪点、让边缘更清楚。这属于低级图像处理（low-level image processing）常见的任务。图像改善（image improvement）强调让人更容易观察和解释，也可以为后面的机器分析准备更合适的输入。低级并不表示技术简单或效果不重要，而是指离像素数据比较近的处理层次。

当任务变成“这张照片里有几个人”或“这个号码是什么”时，输出已经是数量、位置或类别。图像分析（image analysis）和高级图像处理（high-level image processing）更关注结构、特征与语义。其中，图像分割（image segmentation）把像素分成有意义的区域，特征提取（feature extraction）把形状、颜色或纹理概括成可比较的量，识别（recognition）进一步判断对象属于哪一类。低级处理常为这些步骤服务。本讲义原课程的重点是低级图像处理，尤其是图像增强。

这些层次之间没有一条适用于所有教材的硬边界。例如，压缩表示（compressed representation）既能为存储和传输节省空间，也能服务后续分析，不能仅因为它出现在分析流程里就认定为“只属于高级处理”。理解具体任务的输入、输出和目的，比死记分类更可靠。

\begin{figure}[htbp]
\centering
\begin{tikzpicture}[font=\small,node distance=.52cm]
\node[draw,align=center,text width=2.5cm,minimum height=1.05cm] (a) {图像获取\\把场景记录成图像};
\node[draw,align=center,text width=2.5cm,minimum height=1.05cm,right=of a] (b) {图像增强与复原\\改善可见性与质量};
\node[draw,align=center,text width=2.5cm,minimum height=1.05cm,right=of b] (c) {分割与特征提取\\找到目标并描述};
\node[draw,align=center,text width=2.5cm,minimum height=1.05cm,right=of c] (d) {识别与决策\\回答实际问题};
\draw[-{Latex}] (a)--(b);\draw[-{Latex}] (b)--(c);\draw[-{Latex}] (c)--(d);
\node[align=center,below=.55cm of b,text width=10cm,xshift=1.6cm] {彩色处理、几何校正、压缩和存储可按任务插入流程；\\实际系统不必把所有步骤依次执行一遍。};
\end{tikzpicture}
\caption{从图像到任务答案的一条典型路线。流程由用途决定。}
\end{figure}

图像处理属于信号处理（signal processing）的一部分。信号可以理解为“随位置或时间变化、携带信息的量”：录音随时间变，照片随二维位置变。它也可以作为工程工具，例如用图像判断工件是否缺角。传统课程从电路分析、信号与系统、随机信号处理逐步进入图像处理；本书先从可见的像素操作建立直觉，再补上需要的数学语言。

计算机图形学（computer graphics）主要研究如何根据模型、几何和材质生成图像，例如根据三维模型合成汽车广告；图像处理则常从已有图像出发。计算机视觉（computer vision）关心机器如何通过图像理解场景，图像处理常是它的组成步骤。模式识别（pattern recognition）研究如何依据特征判别类别，既可处理图像，也可处理语音等数据。人工智能（artificial intelligence）与图像处理、图像识别相交；“图像处理只做像素而人工智能只做识别”不是可靠的分界。

\subsection{相机以外也可以产生图像}

图像获取（image acquisition）把来自场景的测量结果组织成二维或更高维数据。普通摄像机记录可见光；鱼眼相机（fisheye camera）用很宽的视角观察环境，画面也因此有明显的几何变形；全向相机（omnidirectional camera）用于观察周围多个方向。红外相机（infrared camera）记录红外波段的信号。雷达（radar）利用电磁波回波，X 射线成像（X-ray imaging）利用穿透后的测量，计算机断层成像（computed tomography, CT）从多方向测量重建内部结构。测距相机（range camera）可记录距离，深度图（depth map）的像素值代表远近，不一定代表明暗。

因此，“图像里的亮点”要结合数据含义解释。普通灰度照片里的高值通常显示得较亮；深度图中的高值可能只是更远的地方。增强或识别之前，必须知道像素记录的是什么量。

原课程的应用案例可以沿着“观测对象、希望得到的答案”来理解。门禁与客流统计要知道谁进出、多少人经过；A--B 门防尾随系统要判断两道门之间是否多出一个人；遗留异物检测要发现原来没有而持续出现的物体；车牌识别要把字符图像转成文字；机场飞机起降跟踪要估计飞机位置随时间怎样变化；无人机航拍可用于交通参数统计和竣工图测绘。它们对分辨率、光照、时间采样和几何校正的要求不同，因此不能只靠一种“万能增强”。

\subsection*{原题 F2 低级与高级图像处理}
\textbf{题目。}低级图像处理与高级图像处理的区别？本课程的重点是哪一类？

\textbf{完整解答。}先比较输出。低级处理常把输入图像变成另一幅图像，目的是改善可见性或数据质量，如去噪、平滑、锐化和对比度调整；高级处理更关注从图像中获取结构或语义，如分割对象、描述特征、识别类别。再比较服务对象：原课程把前者重点描述为面向人的观察，把后者重点描述为面向机器的自主感知；实际系统中，增强也经常服务机器识别。最后指出课程范围：本课程重点是低级图像处理。作答时给出“去噪后仍是一张图”与“识别后得到车牌号”两个例子，就能说明区别，而不只是重复名词。

\section{人眼为什么会改变对同一张图的判断}
\label{sec:foundations-eye}

从室外走进较暗的房间，最初可能看不清家具，停留一段时间后细节逐渐出现。家具没有突然变亮，变化的是视觉系统的工作状态。图像增强常是为人观看服务的，因此仅知道像素数值还不够，还要理解人眼怎样接收光，以及周围环境怎样影响明暗判断。

\subsection{光怎样到达视网膜}

角膜（cornea）是眼球前方透明、参与折射的结构；巩膜（sclera）是外部的坚韧保护层。虹膜（iris）中央的开口是瞳孔（pupil），其大小改变进入眼内的光量。晶状体（lens）进一步帮助光线聚焦。眼球后部的视网膜（retina）包含把光信号转成神经信号的感光细胞（photoreceptor）。脉络膜（choroid）位于巩膜与视网膜之间，参与供血并吸收杂散光。视神经（optic nerve）离开眼球的位置没有通常的感光细胞，形成盲点（blind spot）。

\begin{figure}[htbp]
\centering
\begin{tikzpicture}[x=1cm,y=1cm,font=\small]
\draw[thick] (0,0) ellipse (3 and 1.75);
\draw[thick] (-2.75,.65) .. controls (-3.55,.2) and (-3.55,-.2) .. (-2.75,-.65);
\draw[very thick] (2.3,1.08) arc (34:-34:2.78 and 1.92);
\draw[thick] (-1.85,0) ellipse (.3 and .85);
\draw[very thick] (-2.3,.95)--(-2.3,.35);
\draw[very thick] (-2.3,-.95)--(-2.3,-.35);
\draw[-{Latex},thin] (-4.4,.6)--(-2.2,.3)--(-1.6,.25)--(2.7,-.55);
\draw[-{Latex},thin] (-4.4,-.6)--(-2.2,-.3)--(-1.6,-.25)--(2.7,.55);
\draw[thick] (2.6,-.8)--(3.5,-1.2);\draw[thick] (2.8,-1.05)--(3.45,-1.45);
\node[anchor=east] at (-3.7,1.2) {角膜};\draw (-3.55,1.1)--(-3.05,.45);
\node at (-2.5,1.55) {虹膜};\draw (-2.5,1.35)--(-2.3,.75);
\node at (-.95,1.5) {晶状体};\draw (-1,1.25)--(-1.7,.65);
\node at (.8,2.05) {巩膜};\draw (.8,1.85)--(.8,1.68);
\node[anchor=west] at (3.25,.75) {视网膜};\draw (3.2,.65)--(2.78,.3);
\node[anchor=west] at (3.55,-1.3) {视神经};
\node[align=center] at (.3,-2.35) {示意图只表示位置与光路关系，未按真实解剖比例绘制};
\end{tikzpicture}
\caption{眼睛将进入的光线聚焦到视网膜上，再由神经系统处理。}
\end{figure}

晶状体与普通固定玻璃镜片不同，它可以改变形状。睫状体（ciliary body）及相关结构调节晶状体的形状，这种为不同距离重新聚焦的过程叫调节（accommodation）。一般而言，看近处时晶状体较凸，看远处时较扁。近视（myopia）与远视（hyperopia）涉及眼的屈光能力、眼轴长度和成像位置，并不只是“晶状体这一次变厚或变薄”这一动作的另一种叫法。

\subsection{视锥与视杆各自适合怎样的工作}

视锥细胞（cone）在较明亮环境中发挥主要作用，支持颜色感知和精细空间辨别。视网膜中央的中央凹（fovea）具有很高的视锥密度，所以把字放在视线正中通常更容易看清。视锥主导的状态叫明视觉（photopic vision）。

视杆细胞（rod）在较暗环境中更敏感，但不能提供与视锥相同的颜色信息和精细空间分辨能力。视杆分布于较广的视网膜区域，中心凹的中心缺乏视杆；视杆主导的状态叫暗视觉（scotopic vision）。从明到暗的过渡环境中，两类细胞可共同起作用，称为中间视觉（mesopic vision）。“暗光下视锥不工作”适合作为极暗条件下的简化说法，不应理解为光线稍暗，视锥就全部关闭。

\begin{center}
\begin{tabularx}{\linewidth}{@{}p{2.1cm}XX@{}}
\toprule
比较项目 & 视锥细胞 & 视杆细胞\\
\midrule
主要适用环境 & 较明亮，主导明视觉 & 昏暗，主导暗视觉\\
颜色信息 & 不同类型视锥共同支持颜色辨别 & 不提供正常的颜色辨别\\
细节与分布 & 中央凹密集，精细辨别能力强 & 分布较广，暗光敏感但精细辨别较弱\\
对图像设计的启发 & 良好照明有助于看清文字和细小边缘 & 很暗环境下不宜仅用颜色编码关键差别\\
\bottomrule
\end{tabularx}
\end{center}

这一区别解释了为什么夜间只用红绿两种颜色标记信息可能不够可靠，也解释了为什么“能看见某个亮点”不等于“能分辨亮点内部的细节”。但不能由此推出“亮度感觉主要总是由视杆完成”。白天、正常室内照明和常见屏幕观看条件下，视锥对亮度与细节判断都十分重要。

\subsection{亮度适应与可察觉的变化}

物理强度描述光的大小；感知亮度（perceived brightness）描述人觉得有多亮。两者相关，但不是简单的一对一数字复制。视觉系统会根据环境调整灵敏度，这叫亮度适应（brightness adaptation）。当前适应的背景水平称为亮度适应级（brightness adaptation level）。

原课程用约 $10^{-6}$ 至 $10^4\ \mathrm{mL}$ 描述视觉系统跨不同适应状态可处理的很宽亮度范围。这里 $\mathrm{mL}$ 是毫朗伯（millilambert），不是体积单位毫升。这个约十个数量级的范围是经典教学示意，不能当作任何观察者、任何实验下都精确相同的端点。关键结论是：\textbf{眼睛不能在同一适应状态下同时清楚分辨这个巨大范围内的所有亮暗层次。}就像相机一次曝光不一定同时保留太阳附近和阴影深处的细节，人眼的瞬时辨别也受到当前状态限制。

现在考虑一个具体实验：背景亮度为 $I$，在背景上增加一块很小的亮度差 $\Delta I$。符号 $\Delta$ 读作“变化量”，所以亮块的亮度是 $I+\Delta I$。差别太小时，人可能无法稳定辨认；增加到某个水平后，就更容易察觉。用约定的实验判据测得的可察觉增量叫差别阈限（difference threshold）。原题把回答“看见了”的比例达到 $50\%$ 时的增量记为 $\Delta I_{50}$。下标 50 是实验判据的标记，不是把 $\Delta I$ 乘 50。

韦伯比（Weber ratio）定义为
\[
  W=\frac{\Delta I_{50}}{I},\qquad I>0.
\]
$I$ 与 $\Delta I_{50}$ 必须用相同的物理量、相同单位；分子分母单位相消，所以 $W$ 没有单位。$W=0.02$ 表示增量达到背景的 $2\%$ 时达到所用察觉判据。比值越小，说明在该实验条件下能够察觉的\emph{相对变化}越小，亮度分辨能力越好。

\begin{figure}[htbp]
\centering
\begin{tikzpicture}[font=\small,x=1cm,y=1cm]
\fill[gray!30] (0,0) rectangle (4.3,1.8);
\fill[gray!27] (1.6,.5) rectangle (2.7,1.3);
\draw (0,0) rectangle (4.3,1.8);
\node at (2.15,2.2) {背景 $I$ 与亮块 $I+\Delta I$};
\node[align=center] at (2.15,-.5) {比较的是差值相对于背景的比例\\图中灰度仅为示意，不模拟实验阈限};
\draw[-{Latex}] (5.1,0)--(11.1,0) node[right] {$W$};
\foreach \x/\lab in {5.1/0,7.1/{0.02},9.1/{0.04}} {\draw (\x,.08)--(\x,-.08);\node[below] at (\x,-.08) {$\lab$};}
\draw[thick] (7.1,.15)--(7.1,1.15);\node[above,align=center] at (7.1,1.15) {$2/100$\\相对阈限 $2\%$};
\draw[thick] (9.1,.15)--(9.1,1.15);\node[above,align=center] at (9.1,1.15) {$16/400$\\相对阈限 $4\%$};
\draw[-{Latex}] (10.2,-1)--(6,-1) node[midway,below] {越小，能察觉的相对变化越小};
\end{tikzpicture}
\caption{韦伯比必须用相对量比较。绝对增量变大，并不自动代表相对分辨能力变差。}
\end{figure}

在一定范围和固定实验条件下，差别阈限可能近似与背景强度成比例，即 $\Delta I_{50}\approx kI$。这叫韦伯定律（Weber's law）的近似形式；$k$ 是该条件下近似不变的比例。如果背景变为原来的 4 倍，阈限也恰好变成 4 倍，韦伯比不变。因此，“更亮背景需要更大的绝对增量”与“更亮背景下相对辨别更差”不能混为一谈。在非常暗的环境以及不同适应、观察时间、目标大小等条件下，这个简单近似也可能不成立。

\subsection*{原题 E4 韦伯比的计算与解释}
\textbf{题目。}（1）写出韦伯比的定义并说明其含义；（2）某背景亮度 $I=100$（单位任意但前后一致），在题设 $50\%$ 察觉判据下的增量 $\Delta I_{50}=2$，求韦伯比；（3）换成更亮背景 $I=400$、$\Delta I_{50}=16$，比较两次的亮度分辨能力。

\textbf{第一步，写定义。}$W=\Delta I_{50}/I$，表示达到该察觉判据所需的相对亮度变化。不是直接比较亮块总亮度 $I+\Delta I_{50}$。

\textbf{第二步，代入第一组。}
\[
 W_1=\frac{2}{100}=0.02=2\%.
\]
背景每 100 个强度单位增加 2 个单位，就达到题设判据。

\textbf{第三步，代入第二组。}
\[
 W_2=\frac{16}{400}=0.04=4\%.
\]
背景每 100 个强度单位需要增加相当于 4 个单位，才能达到同样判据。

\textbf{比较结论。}$W_1<W_2$，所以在这两组指定实验条件下，第一组的\emph{相对亮度分辨能力}较好。这个结论只比较题中两次测量，不能推出“环境越暗越容易分辨亮度”，也不能证明“暗适应总比明适应分辨能力强”。若第二组给出的是 $\Delta I_{50}=8$，则 $8/400=0.02$，两组相对阈限就相同。

\subsection{周围的像素会改变中心的视觉效果}

把完全相同的灰色小方块分别放到黑色和白色背景上，黑背景上的小方块往往显得更亮。这叫同时对比（simultaneous contrast）：某处的感知亮度不仅取决于它本身，还取决于周围区域。检查下图中心小块的数字可知，两块都用相同填充绘制；视觉上的不同不是计算机改了中间数值。

\begin{figure}[htbp]
\centering
\begin{tikzpicture}[font=\small,x=1cm,y=1cm]
\fill[black!85] (0,0) rectangle (4.3,2.2);
\fill[black!8] (5.1,0) rectangle (9.4,2.2);
\draw (0,0) rectangle (4.3,2.2);\draw (5.1,0) rectangle (9.4,2.2);
\fill[black!45] (1.55,.5) rectangle (2.75,1.7);
\fill[black!45] (6.65,.5) rectangle (7.85,1.7);
\node at (2.15,-.38) {暗背景};\node at (7.25,-.38) {亮背景};
\node at (4.7,-.9) {两个中心方块的实际填充灰度完全相同};
\end{tikzpicture}
\caption{同时对比的示意图。打印条件可能改变错觉强弱，但中心数值相同这一点不变。}
\end{figure}

马赫带（Mach bands）是另一种例子。在某些明暗过渡的交界附近，人会觉得存在额外的亮边或暗边，而物理亮度并没有那样的局部峰谷。视觉系统的侧抑制（lateral inhibition）等机制有助于解释这种边缘强化现象：一个位置的神经反应会受到邻近位置影响。视错觉（optical illusion）则还包括对长度、方向和形状的误判。它们共同提醒我们，图像数值的客观变化与观看者的主观判断不能简单画等号。

对图像增强来说，这意味着均匀分布的灰度数值不一定最好看，更亮也不一定更清楚。判断结果时要说明观察任务：看文字是否清晰，找裂缝是否容易，还是颜色是否自然。与此同时，主观“看上去舒服”不能代替测量精度或识别准确率。

\subsection*{原题 F3 视锥与视杆及其设计意义}
\textbf{题目。}视锥细胞与视杆细胞有什么差别？这对图像设计有什么意义？

\textbf{完整解答。}视锥在较明亮环境中支持颜色与精细细节辨别，中央凹密度高；视杆在较暗环境中更敏感、分布较广，但颜色辨别和精细空间分辨较弱。设计供人观看的图像时，应兼顾观察照明、文字与背景的对比，以及颜色之外的形状或明暗标记。评价增强效果，还要考虑亮度适应与同时对比等视觉现象。

不能把这个结论延伸成“中值滤波不会损失视觉信息”。中值滤波可能去掉小目标、细线和纹理，其效果必须结合目标与噪声检查；细胞的空间汇聚也不是中值滤波无损的证明。视锥、视杆的区别给出的是视觉设计依据，而不是任意图像算法都有效的保证。

\section{从连续场景到有限像素及像素之间的关系}
\label{sec:foundations-digitization}

拍照时，真实世界并不会预先站在一张整数网格上。空间位置可以连续变化，接收到的光强也可以连续变化。计算机却只能保存有限个位置和有限精度的数值。数字化（digitization）因此包含两个彼此独立的选择：在哪里测量，以及把测量结果记成哪些数。

\subsection{采样决定在哪里量 量化决定能记多细}

先用测温作类比。沿走廊每隔一米放一个温度计，是在选择测量位置；把 $23.37\,^{\circ}\mathrm{C}$ 记成 $23\,^{\circ}\mathrm{C}$ 或 $23.4\,^{\circ}\mathrm{C}$，是在选择数值精细程度。增加温度计个数，不会自动增加每支温度计的小数位；保留更多小数，也不会补回没有放温度计的地方。

采样（sampling）是在空间上离散化：从连续位置选出一组测量位置。对二维图像，可以每隔 $\Delta x$ 选一列、每隔 $\Delta y$ 选一行。$\Delta x$ 和 $\Delta y$ 是采样间隔（sampling interval），单位可以是毫米、微米，或经过成像几何换算后的物体表面距离。符号 $\Delta x$ 表示相邻采样位置之间的横向距离，不是灰度变化。

量化（quantization）是在幅度上离散化：把测得的连续强度归入有限个灰度级。假设一个设备把归一化强度从 0 到 1 分成四档，并用 0、1、2、3 编码，则强度 $0.10$ 和 $0.20$ 可能都进入第 0 档，强度 $0.68$ 进入第 2 档。归一化（normalization）在这里是把原范围按比例变成便于描述的 0 到 1，不是引入新的物理亮度单位。

\begin{figure}[htbp]
\centering
\begin{tikzpicture}[font=\small,x=1cm,y=1cm]
\begin{scope}
\draw[-{Latex}] (0,0)--(5.4,0) node[right]{位置};
\draw[-{Latex}] (0,0)--(0,2.4) node[above]{强度};
\draw[thick,domain=0:5,samples=120] plot (\x,{1.2+.75*sin(95*\x)});
\foreach \x in {.4,1.2,2,2.8,3.6,4.4}{
\pgfmathsetmacro{\yy}{1.2+.75*sin(95*\x)}
\draw[dashed,gray] (\x,0)--(\x,\yy);\fill (\x,\yy) circle (1.8pt);
}
\draw[<->] (.4,-.35)--(1.2,-.35) node[midway,below] {$\Delta x$};
\node at (2.5,2.85) {采样：只取选定位置的值};
\end{scope}
\begin{scope}[xshift=7.5cm]
\draw[-{Latex}] (0,0)--(4.5,0) node[right]{输入};
\draw[-{Latex}] (0,0)--(0,2.5) node[above]{输出级};
\foreach \y/\lab in {.25/0,.9/1,1.55/2,2.2/3}{\draw[dotted,gray] (0,\y)--(4,\y);\node[left] at (0,\y){\lab};}
\draw[very thick] (0,.25)--(1,.25)--(1,.9)--(2,.9)--(2,1.55)--(3,1.55)--(3,2.2)--(4,2.2);
\foreach \x/\lab in {0/0,1/{.25},2/{.50},3/{.75},4/1}{\draw (\x,0)--(\x,-.06);\node[below] at (\x,-.06){$\lab$};}
\node at (2,2.85) {量化：一段输入归为同一级};
\end{scope}
\end{tikzpicture}
\caption{采样改变位置的密度；量化改变允许输出的数值。两者不能互相替代。}
\end{figure}

上图右侧是一种四档量化规则，规定 $[0,0.25)$ 对应 0，$[0.25,0.50)$ 对应 1，依此类推，端点 1 归最后一档。方括号表示包含端点，圆括号表示不包含端点。不同设备也可以使用“舍入到最近代表值”的规则；做数值题时应说明采用哪一种，不能把分档编号与该档的物理代表值混淆。

若以固定步长 $q$ 对连续值做最近值舍入，在不超过量化器范围的条件下，量化误差（quantization error）的绝对值至多为 $q/2$。例如允许值为 $0,0.1,0.2,\ldots$，输入 $0.26$ 舍入为 $0.3$，误差为 $0.04$，没有超过 $0.05$。这里“误差”是存下来的代表值与原值之差，不是某种额外随机噪声必然服从的分布。

实际传感器的像素会对一个小面积、一个曝光时段内的光作响应，并不真是无限小点上的瞬时测量。把它当作网格点采样，是便于学习和计算的模型。像素有限面积、镜头模糊与曝光运动都会影响最终细节，后面讨论分辨率时不能忽略这些因素。

更完整的光场可抽象写成 $I=E(x,y,z,\lambda,t)$。这句话只表示光的测量值可能依赖空间位置 $(x,y,z)$、波长 $\lambda$ 和时间 $t$；$E$ 是描述这种依赖关系的函数名称，不要求在这里进行任何微积分运算。固定观察平面、某个波段和某段拍摄时间后，才得到常用的二维图像 $f(x,y)$。

\subsection{空间分辨率和灰度分辨率不能混为一谈}

空间分辨率（spatial resolution）关注空间上多小的结构还能被区分。数字网格的行列数是相关因素，但必须连同覆盖范围理解：同样 1000 列，覆盖 10 毫米与覆盖 10 千米的实际细节能力完全不同。在同样视场内，更多有效像素通常意味着更细的空间采样；是否真正分辨出来，还受镜头、对焦、运动模糊、噪声和目标对比度影响。

灰度分辨率（gray-level resolution）关注可区分的强度档数，常用每像素位数 $b$ 或级数 $L=2^b$ 描述。空间采样变稀，斜边容易呈锯齿，小目标容易消失；灰度级太少，一段平滑的明暗变化可能变成一条条可见的带状层次，称为假轮廓（false contour）或色带（banding）。

\begin{figure}[htbp]
\centering
\begin{tikzpicture}[font=\small,x=1cm,y=1cm]
\foreach \x in {0,...,3}{\foreach \y in {0,...,3}{
\pgfmathtruncatemacro{\shade}{15+20*\x}
\fill[black!\shade] (.55*\x,.55*\y) rectangle ++(.55,.55);
}}
\draw[step=.55,white] (0,0) grid (2.2,2.2);\draw (0,0) rectangle (2.2,2.2);
\node[align=center] at (1.1,-.6) {位置稀疏\\小细节难以记录};
\foreach \x in {0,...,15}{\foreach \y in {0,...,15}{
\pgfmathtruncatemacro{\shade}{15+5*\x}
\fill[black!\shade] (4+.1375*\x,.1375*\y) rectangle ++(.1375,.1375);
}}
\draw (4,0) rectangle (6.2,2.2);
\node[align=center] at (5.1,-.6) {位置更密\\灰度也较丰富};
\foreach \x in {0,...,3}{\pgfmathtruncatemacro{\shade}{15+20*\x}\fill[black!\shade] (8+.55*\x,0) rectangle ++(.55,2.2);}
\draw (8,0) rectangle (10.2,2.2);
\node[align=center] at (9.1,-.6) {灰度只有四档\\渐变出现分带};
\end{tikzpicture}
\caption{采样不足与灰度档数不足产生不同的损失。网格线与条带为示意，不能据此量化真实设备的分辨率。}
\end{figure}

\subsection*{原题 F4 采样和量化}
\textbf{题目。}什么是采样？什么是量化？图像数字化包含哪两方面？

\textbf{完整解答。}采样把连续空间坐标变成有限网格位置，决定在哪里取值；量化把连续的强度变成有限个灰度级，决定每个位置的测量值怎样编码。数字化包含这两个方面。例如，每隔 $0.1\,\mathrm{mm}$ 测一次是采样选择，把每次测量存成 0 到 255 的整数是量化选择。前者的间隔单位是空间长度，后者的档数是数量、位数单位是 bit。写出这一具体例子，就能避免把“增加像素”与“增加灰度级”当作同一件事。

\subsection*{原题 F5 两种分辨率}
\textbf{题目。}空间分辨率与灰度分辨率的区别？分别由什么决定？

\textbf{完整解答。}空间分辨率描述区分空间细节的能力；在固定视场和其他成像条件下，像素采样越密，能够记录的空间变化通常越细。灰度分辨率描述数值上可用的强度级数，$b$ 比特对应 $2^b$ 个等级。前者不足常表现为锯齿、块状和细节丢失，后者不足常表现为渐变分带、假轮廓。必须补充条件：图像宽高只能描述网格规模，不能脱离物方覆盖范围、镜头与成像质量单独保证实际分辨率；把模糊图扩大成更多像素也不等于获得更多真实细节。

\subsection{周期和频率描述条纹有多密}

考虑一排重复的亮暗条纹。从一条亮纹的中心走到下一条相同亮纹的中心，是一个完整空间周期（spatial period），记为 $T$。若每隔 $0.2\,\mathrm{mm}$ 重复一次，则 $T=0.2\,\mathrm{mm}$。空间频率（spatial frequency）是单位长度内重复的次数，记为 $f$：
\[
  f=\frac{1}{T}.
\]
所以 $T=0.2\,\mathrm{mm}$ 时，$f=5\ \mathrm{cycles/mm}$，意思是每毫米有 5 个完整周期。周期越短，频率越高。“高频”不是说像素值很大，而是说灰度随位置变化得快。

采样频率（sampling frequency）同样是单位长度内取样的次数，记为 $f_s$。若相邻采样点间隔为 $\Delta x$，则
\[
  f_s=\frac{1}{\Delta x}.
\]
例如 $\Delta x=0.05\,\mathrm{mm}$，每毫米取 $1/0.05=20$ 次样，故 $f_s=20\ \mathrm{samples/mm}$。\emph{cycles/mm} 数的是信号的重复周期，\emph{samples/mm} 数的是测量点；两者有相同的“每长度”维度，但不是同一个对象。

\subsection{采样太稀为什么会产生假的条纹}

每隔很长距离才测一次条纹，就可能错过中间的明暗变化。更严重的是，稀疏测得的一串值可能同时符合两种不同条纹：原来的快速变化，以及另一种较慢变化。数字系统无法仅凭这串值判断哪种是真的，这叫混叠（aliasing）。相机拍摄细密织物、显示器或栅栏时出现的粗条纹，常与这种现象有关。

奈奎斯特采样定理（Nyquist sampling theorem）讨论一个理想条件：信号的频率成分存在最高上限 $f_{\max}$，超过它的成分完全不存在，这叫带限（band-limited）。在理想采样与理想重建等条件下，为避免最高频率正好落在含糊的边界，采用
\[
  f_s>2f_{\max},\qquad
  \Delta x<\frac{1}{2f_{\max}}.
\]
其中 $f_{\max}$ 是信号中最高的空间频率，$f_s$ 是采样频率，$\Delta x$ 是采样间隔。把间隔放在分母会使不等号方向反过来：要求频率更高，等价于要求间隔更小。

一些教材把临界条件写成 $f_s\geq2f_{\max}$，方便记忆“每周期至少两个样本”。但是等号是理想边界，不应作为对任意相位都可靠的工程设计点。相位（phase）在这里先理解为条纹相对采样网格平移了多少：同一排条纹平移半格，亮暗峰位置就会落到不同采样位置。

\begin{figure}[htbp]
\centering
\begin{tikzpicture}[font=\small,x=1.65cm,y=1cm]
\draw[-{Latex}] (0,0)--(6.2,0) node[right]{位置 $x$};
\draw[-{Latex}] (0,-1.3)--(0,1.5) node[above]{强度变化};
\draw[thick,domain=0:6,samples=180] plot (\x,{sin(180*\x)});
\foreach \x in {0,...,6}{\fill (\x,0) circle (2pt);\draw[dotted] (\x,-1.2)--(\x,1.1);}
\draw[very thick,dashed] (0,0)--(6,0);
\draw[<->] (0,-1.6)--(2,-1.6) node[midway,below]{一个周期 $T$};
\node[align=center] at (4,1.65) {每隔 $T/2$ 采样一次\\恰好全落在零点};
\node[align=center] at (4,-1.65) {采样结果与一条恒定线相同\\真实的振荡已无法由样本判定};
\end{tikzpicture}
\caption{临界采样的相位风险。曲线表示相对某背景的亮度变化，允许出现负的变化量，并非负的实际光强。}
\end{figure}

上图的真实曲线在采样点之间仍然起伏，但所有记录值都相同。增加一个足够大的背景亮度，就得到完全非负的真实强度，问题仍然存在。因此这不是“实际图像不会有负值”可以排除的数学小把戏，而是临界位置确实可能漏掉变化。

真实图像中的尖锐边缘并不严格带限。实际系统通常先用光学或数字低通来削弱过细变化，再采样；这叫抗混叠（anti-aliasing）。降低图像尺寸时也一样，不能只随手每隔几行几列丢掉像素，否则新的采样网格可能把细纹理误认为粗纹理。先适当平滑会牺牲一部分细节，却能减轻虚假结构。

\subsection*{原题 E5 从最细周期计算采样要求}
\textbf{题目。}某图像最细结构的空间周期为 $0.2\,\mathrm{mm}$，要在数字化时不发生混叠，采样间隔最大为多少？采样频率至少为多少？

\textbf{第一步，明确模型。}按题目的理想周期模型，把最细周期 $T_{\min}=0.2\,\mathrm{mm}$ 对应的频率作为最高频率。若所谓“结构”是尖锐条纹，还包含更高频成分，则需要先说明带限或预滤波假设，不能把条纹宽度机械代入这个周期公式。

\textbf{第二步，周期变频率。}
\[
 f_{\max}=\frac{1}{T_{\min}}
 =\frac{1}{0.2\,\mathrm{mm}}
 =5\ \mathrm{cycles/mm}.
\]
\textbf{第三步，求临界值。}
\[
 2f_{\max}=10\ \mathrm{samples/mm},\qquad
 \frac{1}{2f_{\max}}=0.1\,\mathrm{mm}.
\]
因此，原讲义习惯写的答案是“至少 10 次/mm、间隔不大于 $0.1\,\mathrm{mm}$”。更严谨地说，$10$ 次/mm 与 $0.1\,\mathrm{mm}$ 是\textbf{临界界限}：若要求对一般相位可靠，应取 $f_s>10$ 次/mm、$\Delta x<0.1\,\mathrm{mm}$，而不把等号当作安全点。

\textbf{第四步，给出有裕量的例子。}取 $\Delta x=0.05\,\mathrm{mm}$，则 $f_s=20$ 次/mm，一个最细周期内有 $0.2/0.05=4$ 个采样间隔。取 $0.04\,\mathrm{mm}$ 则对应 25 次/mm。它们比临界条件更稳妥，但“选四点或五点”仍不是任何镜头、噪声和目标条件下的质量保证，真实系统还需成像测试。

\subsection{视场把传感器尺寸换成物体上的尺寸}

相机参数中的像元尺寸（pixel pitch）描述传感器上相邻采样单元的间距，例如 $4.52\,\mu\mathrm{m}$。拍到的物体往往比传感器大得多，因此传感器上一个像元不等于物体上同样大小的一块区域。物方像素尺寸（object-space pixel size），在遥感中常称地面采样距离（ground sample distance, GSD），才描述相邻像素中心在被拍平面上相隔多远。

视场（field of view, FOV）在本节指相机在某个物体平面上覆盖的实际宽或高，不是镜头的视角角度。焦距（focal length）$f$ 描述成像几何的一项镜头参数；物距（object distance）$Z$ 描述物体平面离相机的距离。对近似正对相机的平面、较远物体，并忽略畸变，可以用相似三角形得到
\[
 \frac{W}{S_x}\approx\frac{Z}{f},
 \qquad W\approx S_x\frac{Z}{f}.
\]
式中 $W$ 是横向物方视场，$S_x$ 是传感器有效宽度；同理 $H\approx S_yZ/f$。相似三角形（similar triangles）只是说两只三角形形状相同，对应边按相同比例放大。物距远大于焦距时，像距可以近似用焦距代替；近距离精密测量时应使用更准确的成像模型并做标定。

\begin{figure}[htbp]
\centering
\begin{tikzpicture}[font=\small,x=1cm,y=1cm]
\draw[thick] (0,-1.8)--(0,1.8);
\draw[thick] (8,-.65)--(8,.65);
\fill (6,0) circle (2pt);
\draw (0,1.8)--(6,0)--(8,-.6);
\draw (0,-1.8)--(6,0)--(8,.6);
\draw[dashed] (-.5,0)--(8.6,0);
\draw[<->] (-.5,-1.8)--(-.5,1.8) node[midway,left] {$W$};
\draw[<->] (8.5,-.6)--(8.5,.6) node[midway,right] {$S_x$};
\draw[<->] (0,-2.25)--(6,-2.25) node[midway,below] {物距 $Z$};
\draw[<->] (6,-1.2)--(8,-1.2) node[midway,below] {约为 $f$};
\node at (0,2.15) {物体平面};\node at (6,.5) {成像中心};\node at (8,1.05) {传感器};
\node[align=center] at (4,-3.1) {图为放大的几何示意；实际物距通常远大于焦距};
\end{tikzpicture}
\caption{相似三角形把传感器上的长度换算为物体平面上的长度。}
\label{fig:foundation-fov}
\end{figure}

若横向有 $N_x$ 个像素、像元间距为 $p$，则常用有效宽度近似为 $S_x=N_xp$。这个表达把每个像素看作宽度 $p$ 的覆盖单元；若只讨论第一个到最后一个像素\emph{中心}的距离，则是 $(N_x-1)p$，两种口径应按需求区分。在本题的视场估算中使用前一种常见工程口径。代入可得
\[
 W\approx N_xp\frac{Z}{f},\qquad
 d\approx\frac{W}{N_x}=p\frac{Z}{f}.
\]
$d$ 是一个像素在物方对应的长度。例如 $d=0.2\,\mathrm{mm/pixel}$，表示每像素覆盖的标称物方宽度约为 $0.2\,\mathrm{mm}$；10 个像素约对应 2 毫米。这里的 pixel 是计数说明，不是另一个物理长度单位。

公式也展示出直观关系。在其他条件不变时，拍摄更远的物体会覆盖更大的视场，物方像素也更大；增加焦距会缩小视场，物方采样变细。只增加传感器像素数而保持同样像元尺寸、焦距和物距，主要扩大覆盖范围，并不自动减小 $d$。固定视场下增加像素数，才直接把 $d=W/N_x$ 变小。

\subsection*{原题 E2 分辨率的工程换算}
\textbf{题目。}某相机像素数 $7952\times5304$，像素尺寸 $4.52\,\mu\mathrm{m}$，物距 $300\,\mathrm{m}$，镜头焦距 $36.18\,\mathrm{mm}$。求：（1）传感器尺寸；（2）视场，横向乘纵向；（3）图像中一个像素代表多大的实际尺寸。

\textbf{第一步，整理单位。}
\[
 1\,\mathrm{mm}=1000\,\mu\mathrm{m},\quad
 1\,\mathrm{m}=1000\,\mathrm{mm},\quad
 p=4.52\,\mu\mathrm{m}=0.00452\,\mathrm{mm}.
\]
将像元尺寸先换成毫米，可以直接与焦距的毫米相消。像素数是计数，不需要再换长度单位。

\textbf{第二步，求传感器有效尺寸。}
\begin{align*}
 S_x&=7952\times0.00452\,\mathrm{mm}=35.94304\,\mathrm{mm},\\
 S_y&=5304\times0.00452\,\mathrm{mm}=23.97408\,\mathrm{mm}.
\end{align*}
因此传感器约为 $35.94\,\mathrm{mm}\times23.97\,\mathrm{mm}$。保留中间计算的足够位数，可以减少反复四舍五入造成的差异。

\textbf{第三步，求物方视场。}在上述远距离、正视平面的近似下，
\begin{align*}
 W&=35.94304\,\mathrm{mm}\times
       \frac{300\,\mathrm{m}}{36.18\,\mathrm{mm}}
   \approx298.04\,\mathrm{m},\\
 H&=23.97408\,\mathrm{mm}\times
       \frac{300\,\mathrm{m}}{36.18\,\mathrm{mm}}
   \approx198.79\,\mathrm{m}.
\end{align*}
分子与分母中的毫米相消，最后留下米；若算出来仍是“毫米乘米”，就说明单位约分没有完成。

\textbf{第四步，求单像素的物方长度。}
\begin{align*}
 d&=0.00452\,\mathrm{mm}\times
       \frac{300\,\mathrm{m}}{36.18\,\mathrm{mm}}\\
  &\approx0.0374793\,\mathrm{m/pixel}
   =3.74793\,\mathrm{cm/pixel}
   \approx3.75\,\mathrm{cm/pixel}.
\end{align*}
也可以用 $298.04/7952$ 检查横向结果，用 $198.79/5304$ 检查纵向结果。方形像元和同一成像几何应给出相同的物方尺度。

\textbf{结果解释。}该参数组合的视场约为 $298.04\,\mathrm{m}\times198.79\,\mathrm{m}$，标称采样尺度约为 $3.75\,\mathrm{cm/pixel}$。原稿中的 $4.52\,\mathrm{cm/pixel}$ 与题设不相容，不能因为量级接近就当作同一个结果。物方长度明显小于一个像素的毫米级裂缝，通常不能靠这组参数可靠分辨其宽度；增加软件显示倍率不会把缺失测量补回来。

% 独立内容单元；对应原模块 4.1--4.3、4.6。
% 仅依赖主文档提供的标准数学、表格与 TikZ 宏包。
\section{灰度变换与直方图 从一个像素学到整幅图}

本节解决一个很具体的问题：照片里的东西明明存在，为什么看不清，怎样让它更容易看清？我们先学会给单个像素换一个数，再学会统计整幅图中的数，最后根据统计结果自动换数。对应原讲义的模块 4.1--4.3 和 4.6；原题 B1--B3、C1--C2、E1、F1、G3 都在本节保留题号并逐步解释。

\subsection{增强与复原 先分清想解决什么问题}

\textbf{图像增强（image enhancement）}是为了某个用途，改变图像的显示或特征。例如把较暗的车牌提亮，使字符容易辨认；把边缘突出，使零件轮廓容易测量。这里的“好”与用途有关：照片看起来自然，和机器容易读出字符，可能需要不同的处理。

\textbf{图像复原（image restoration）}是根据图像如何被破坏的模型，估计未被破坏时的图像。例如知道相机在曝光时水平晃动，就尝试根据运动模糊模型减轻拖影；知道噪声具有某种统计规律，就据此估计较干净的图像。

\textbf{先验知识（prior knowledge）}就是处理前已经掌握的信息，如运动方向、噪声特点或镜头模糊程度。“复原使用模型”不代表一定能恢复真实原图。数据可能已经丢失，模型也可能不准确。增强也可以有客观评价指标，例如字符识别正确率，因此不能把二者简单分成“完全主观”和“完全客观”。

想象同一张模糊的黑白照片：增大文字与背景的灰度差，是增强；估计造成拖影的过程再尝试反向补偿，是复原。两者也可以串联使用。

\subsubsection*{原题 F1 图像增强与图像复原有什么区别}
\label{original:F1}

\textbf{第一步，比较目标。}增强问“怎样让它更适合当前用途”；复原问“怎样从退化结果估计原来的图像”。

\textbf{第二步，比较依据。}增强可以根据观察效果、应用指标来选择变换；复原通常依赖退化模型、噪声模型等先验知识。

\textbf{第三步，举出例子。}低对比度照片的灰度拉伸属于增强；根据已知运动方向和程度进行去模糊属于复原。去噪可能属于增强流程，也可能建立在复原模型上，不能仅凭“去噪”这个名字判断。

\textbf{可写在答卷上的答案：}图像增强以特定应用为目标，改善可视效果或后续分析效果；图像复原依据退化过程的先验或模型，估计接近未退化状态的图像。二者的目标和依据不同，但可能使用相似工具，也都不保证恢复已经丢失的信息。

\textbf{处理途径也要分清。}直接在像素网格上处理称为\textbf{空间域处理（spatial-domain processing）}：只根据当前像素灰度换数，叫点处理；结合周围一小片像素计算，属于邻域或掩模处理，例如平滑与锐化。另一条途径是\textbf{频率域处理（frequency-domain processing）}：先把图像变换为不同变化频率的分量，再决定哪些保留、哪些减弱，最后变回图像。后面的滤波章节会解释后两类；本节先把点处理学明白。

\subsection{点运算的共同语言 输入一个数 输出一个数}

\textbf{灰度（gray level）}在本节中是表示像素明暗的数。对常见的 8 位无符号灰度图，0 表示黑，255 表示白，中间的数表示不同程度的灰。\textbf{灰度级数}记作 $L$，所以 8 位图有 $L=256$ 级，最大灰度为 $L-1=255$。注意“256 级”包含 0，因此最大数不是 256。

\textbf{点运算（point operation）}是根据某个像素自己的灰度，计算它的新灰度。把输入灰度记为 $r$，输出灰度记为 $s$，变换规则记为 $T$，就得到
\[
s=T(r).
\]
括号的意思是“把 $r$ 交给规则 $T$ 处理”，不是乘法。若规则是 $T(r)=r+20$，输入 $r=40$，便得到 $s=60$。像素仍在原位置，只是灰度变了。

画变换曲线时，横轴放输入 $r$，纵轴放输出 $s$。找一个输入，竖直向上碰到曲线，再水平向左读纵轴，就是输出。曲线在 $s=r$ 的对角线上方，该输入就变亮；在下方，该输入就变暗。

\textbf{全局点变换（global point transform）}使用全图共用的一条规则：两个位置的输入都等于 40，输出就相同。\textbf{局部处理（local processing）}会利用像素周围的信息；即使输入相同，不同位置也可能得到不同输出。HE 的规则由整幅图统计得到，但规则一旦确定，对所有像素使用同一张灰度对照表。

\subsubsection*{计算结果超出范围时怎么办}

计算可以产生小数，也可以得到负数或超过 255 的数。若最终保存为 8 位图，通常先\textbf{截断到合法范围（clipping）}，再\textbf{取整（rounding）}。例如 $260\to255$，$-7\to0$，$52.7\to53$。本节的 $\operatorname{round}(x)$ 表示把 $x$ 四舍五入到最近整数；题目若给出其他取整规则，要按题目执行。

截断会丢失差别：原本算得 260 与 300 的两个像素都变成 255。\textbf{饱和（saturation）}指许多不同的真实亮度被记录成同一个最大值或最小值。若一片白色区域的像素已经全部是 255，仅靠灰度变换无法知道它原来有哪些不同细节。

\subsection{逐个理解常见灰度变换}

\subsubsection*{线性变换 明暗平移与对比度缩放}

图像处理中常把 $s=ar+b$ 称为\textbf{线性灰度变换（linear gray-level transform）}；严格数学名称是\textbf{仿射变换（affine transform）}，当 $b\ne0$ 时不满足线性算子的定义。其中 $a$ 是乘在输入上的系数，$b$ 是最后加上的偏移量。先乘 $a$，再加 $b$，最后检查范围。

若 $a=1,b=30$，灰度 40 和 50 变成 70 和 80。两者都变亮，但差值仍是 10。若 $a=2,b=0$，它们变成 80 和 100，差值从 10 变成 20。第二种操作拉大了这两个灰度的区别。\textbf{对比度（contrast）}与明暗差异有关，不等于平均亮度。

这里的结论以没有截断为前提。把 140 和 160 都乘 2 后保存为 8 位，两个结果都截断成 255，差异反而消失。因此“系数越大，对比度越好”不成立。

\subsubsection*{反色 把灰度次序翻过来}

\textbf{图像反色或负片变换（image negative）}把黑变白、白变黑，中间灰度反向对应。公式是
\[
s=(L-1)-r.
\]
$L-1$ 是最大灰度，减号表示从最大值中减去输入。例如 8 位图中，$r=30$ 时 $s=255-30=225$；$r=200$ 时 $s=55$。输入越大，输出越小，所以曲线是一条下降直线。

反色改变黑白呈现方式，有时能让原先不容易观察的细节更醒目。但在未取整的情况下，两灰度的差值大小不变，仅次序颠倒；它没有凭空增加细节，也不等于一般意义上的“提高对比度”。

\subsubsection*{对数变换 压缩跨度很大的数}

\textbf{对数变换（logarithmic transform）}让输入越大时，输出增长得越慢。先不记公式，观察这个规律：从 1 增加到 10 是乘 10，从 10 增加到 100 也是乘 10；以 10 为底的对数把它们变成 $0,1,2$，乘法尺度被压成相近的加法尺度。

本节用自然对数 $\ln$，其底数是常数 $e\approx2.718$。常见灰度公式为 $s=c\ln(1+r)$：先加 1，再取对数，最后乘缩放系数 $c$。加 1 是为了让黑色 $r=0$ 可以安全计算，且 $\ln(1)=0$。

若输入范围为 $[0,L-1]$，又希望最大输入仍映射到最大输出，就令
\[
c=\frac{L-1}{\ln L},\qquad
s=\frac{L-1}{\ln L}\ln(1+r).
\]
方括号 $[a,b]$ 表示从 $a$ 到 $b$、含两端的范围。把 $r=L-1$ 代入，$1+r=L$，于是 $\ln L$ 与分母相除为 1，输出确实是 $L-1$。

用只有 8 级的图说明：$L=8$，所以 $c=7/\ln8\approx3.366$。输入 1 对应 $3.366\ln2\approx2.333$，输入 3 对应 $3.366\ln4\approx4.667$，输入 7 对应 7。较暗的输入被明显抬高，而亮部数值被挤得更近。

常见应用是显示\textbf{动态范围（dynamic range）}很大的数据，如傅里叶频谱幅度。若大部分数较小，少数数极大，直接显示可能只看见几个亮点；对数压缩后，较小数值也更容易观察。这里改善的是显示，不是证明微弱细节都是真实信号。

\subsubsection*{幂律与伽马变换 先归一化再判断明暗}

\textbf{幂律变换（power-law transform）}使用指数来弯曲灰度曲线；在图像处理中常称\textbf{伽马变换（gamma transform）}。指数记为希腊字母 $\gamma$，读作 gamma。例如 $u^2$ 表示 $u\times u$，$u^{1/2}$ 表示 $u$ 的平方根。

为了避免单位和范围造成误判，先把输入除以最大灰度：$u=r/(L-1)$，使 $u$ 落在 0 到 1。这个步骤叫\textbf{归一化（normalization）}。然后算 $v=u^\gamma$，最后乘回最大灰度：
\[
s=(L-1)\left(\frac{r}{L-1}\right)^\gamma,\qquad \gamma>0.
\]
$r$ 是原灰度，$u$ 是归一化输入，$v$ 是归一化输出，$s=(L-1)v$。式子从内到外读：先做除法，再做乘方，最后做乘法。

看三个容易手算的归一化输入：
\begin{center}
\begin{tabular}{ccccc}
\toprule
输入 $u$ & $u^{1/2}$ & $u^1$ & $u^2$ & 判断 \\
\midrule
$0.25$ & $0.50$ & $0.25$ & $0.0625$ & 平方根提亮，平方压暗 \\
$0.50$ & $0.7071$ & $0.50$ & $0.25$ & 规律相同 \\
$1$ & $1$ & $1$ & $1$ & 白端保持不变 \\
\bottomrule
\end{tabular}
\end{center}
例如 8 位图 $r=64,\gamma=1/2$：先算 $64/255\approx0.251$，取平方根约为 $0.501$，乘 255 得约 $127.75$，取整为 128。

所以在上述规范化公式下，$0<\gamma<1$ 提亮中间灰度，$\gamma=1$ 不变，$\gamma>1$ 压暗中间灰度。黑端和白端保持不变。\textbf{不能直接把 $s=r^\gamma$ 与这一结论混用。}例如直接对原始灰度算 $64^{1/2}=8$，当然不是提亮；它缺少了范围归一化与输出缩放。

有的资料写 $s=cr^\gamma$。若想保持两个端点，需要选 $c=(L-1)^{1-\gamma}$，这才与上面的规范化公式一致。另有软件把控制参数写成指数的倒数，因此使用软件时必须先看实际公式，不能只看“gamma 大于 1”这几个字。

伽马变换也会改变某些灰度区间的间距，所以可能改善局部灰度对比。但全图统一伽马使用同一条曲线，不能根据空间位置分别照顾暗角与亮区。\textbf{伽马校正（gamma correction）}还可用于补偿采集或显示设备的非线性响应；它与“把照片提亮”有关，但用途并不完全相同。

\subsubsection*{原题 C2 画反色 对数与两种幂律曲线}
\label{original:C2}

\textbf{先画基准。}以下使用归一化坐标 $u=r/(L-1)$、$v=s/(L-1)$，因此两轴均从 0 到 1。换回题目要求的原灰度，只要把轴端 1 改标为 $L-1$，形状不变。虚线 $v=u$ 是不作改变的参照线。

\begin{center}
\begin{tikzpicture}[x=3.9cm,y=3.9cm,>=Stealth,font=\small]
\begin{scope}
\draw[->] (0,0)--(1.12,0) node[right] {$u$};
\draw[->] (0,0)--(0,1.12) node[above] {$v$};
\draw[dashed] (0,0)--(1,1);
\draw[very thick] (0,1)--(1,0);
\node[below left] at (0,0) {0};\node[below] at (1,0) {1};\node[left] at (0,1) {1};
\node[align=center] at (.5,-.24) {反色\quad $v=1-u$};
\end{scope}
\begin{scope}[xshift=7.3cm]
\draw[->] (0,0)--(1.12,0) node[right] {$u$};
\draw[->] (0,0)--(0,1.12) node[above] {$v$};
\draw[dashed] (0,0)--(1,1);
\draw[very thick,domain=0:1,samples=60] plot (\x,{ln(1+7*\x)/ln(8)});
\node[below left] at (0,0) {0};\node[below] at (1,0) {1};\node[left] at (0,1) {1};
\node[align=center] at (.5,-.24) {对数示例\quad $v=\ln(1+7u)/\ln8$};
\end{scope}
\begin{scope}[yshift=-6.1cm]
\draw[->] (0,0)--(1.12,0) node[right] {$u$};
\draw[->] (0,0)--(0,1.12) node[above] {$v$};
\draw[dashed] (0,0)--(1,1);
\draw[very thick,domain=0:1,samples=70] plot (\x,{sqrt(\x)});
\node[below left] at (0,0) {0};\node[below] at (1,0) {1};\node[left] at (0,1) {1};
\node[align=center] at (.5,-.24) {幂律\quad $v=u^{1/2}$\quad 提亮};
\end{scope}
\begin{scope}[xshift=7.3cm,yshift=-6.1cm]
\draw[->] (0,0)--(1.12,0) node[right] {$u$};
\draw[->] (0,0)--(0,1.12) node[above] {$v$};
\draw[dashed] (0,0)--(1,1);
\draw[very thick,domain=0:1,samples=60] plot (\x,{\x*\x});
\node[below left] at (0,0) {0};\node[below] at (1,0) {1};\node[left] at (0,1) {1};
\node[align=center] at (.5,-.24) {幂律\quad $v=u^2$\quad 压暗};
\end{scope}
\end{tikzpicture}
\end{center}

\textbf{逐条作图：}反色连结左上角与右下角；对数从原点到右上角，前段较陡、后段较平；$\gamma<1$ 的规范化幂律位于对角线上方；$\gamma>1$ 的位于下方。用“上方/下方、越来越陡/越来越平”描述最不容易混淆。若使用数学凹凸术语，对数与 $0<\gamma<1$ 的幂律在内部是凹函数，$\gamma>1$ 的幂律是凸函数。

\subsubsection*{分段线性变换 给不同区间不同的待遇}

\textbf{分段线性变换（piecewise-linear transform）}把输入范围切成几段，每段使用一条直线。\textbf{对比度拉伸（contrast stretching）}是其中一个常见用途：把原来挤在较窄范围里的灰度拉开。

设希望输入 $a$ 映到黑色 0，输入 $b$ 映到白色 $L-1$，并且 $a<b$。中间灰度先减去 $a$，把起点移到 0；再除以 $b-a$，换成 0 到 1 之间的相对位置；最后乘 $L-1$：
\[
s=\begin{cases}
0,&r\le a,\\[2pt]
(L-1)\dfrac{r-a}{b-a},&a<r<b,\\[5pt]
L-1,&r\ge b.
\end{cases}
\]
大括号表示“按条件选一行”，不是三行全部相加。例：$L=256,a=80,b=160$，输入 120 位于两端正中，因此 $s=255(120-80)/(160-80)=127.5$，取整为 128。输入 100 变为约 64，输入 140 变为约 191，原来的 40 级间距变成约 127 级。

\begin{center}
\begin{tikzpicture}[x=.045cm,y=.018cm,>=Stealth,font=\small]
\draw[->] (0,0)--(275,0) node[right] {$r$};
\draw[->] (0,0)--(0,290) node[above] {$s$};
\draw[dashed] (0,0)--(255,255);
\draw[very thick] (0,0)--(80,0)--(160,255)--(255,255);
\draw[dotted] (160,0)--(160,255)--(0,255);
\node[below] at (80,0) {80};\node[below] at (160,0) {160};\node[below] at (255,0) {255};
\node[left] at (0,255) {255};\node[below left] at (0,0) {0};
\node[align=center,fill=white] at (211,130) {斜率较大的区段\\拉开灰度差};
\end{tikzpicture}
\end{center}

这条曲线牺牲了 $a$ 以下和 $b$ 以上的差别。若两端也有重要细节，可以设置多个控制点，如 $(0,0),(r_1,s_1),(r_2,s_2),(L-1,L-1)$，用直线连接，给中间重点区间较大斜率，给两端较小斜率。这里 $(r_1,s_1)$ 表示“输入 $r_1$ 对应输出 $s_1$”。

\subsubsection*{灰度级切片 突出一个数值区间}

\textbf{灰度级切片（gray-level slicing）}把落在指定区间中的像素突出显示。设目标范围为 $a\le r\le b$，想让这些像素显示为白色。有两种常见做法：
\begin{enumerate}[leftmargin=*]
\item \textbf{其余归零：}范围内输出 $L-1$，范围外输出 0。结果只有黑白，便于显示目标分布。
\item \textbf{保留背景：}范围内输出 $L-1$，范围外仍输出原值 $r$。目标突出，同时保留环境背景。
\end{enumerate}
例如取区间 $[100,140]$，输入 $90,120,170$，第一种输出 $0,255,0$，第二种输出 $90,255,170$。选择的是\textbf{灰度范围}，不是自动识别“哪个物体”：如果衣服和墙面恰好具有相同灰度，它们会一起被选中。

\subsubsection*{位平面切片 拆开一个数的二进制组成}

\textbf{位平面切片（bit-plane slicing）}按二进制位拆分图像，与“挑选一个灰度区间”完全不同。\textbf{二进制（binary）}只有 0 和 1；各位的权重依次是 $1,2,4,8,16,32,64,128$。

例如十进制灰度 173 可以拆成 $128+32+8+4+1$，其 8 位二进制写为 $10101101$。从右到左把位编号为 $k=0,1,\ldots,7$，该位的值记为 $b_k$，只会等于 0 或 1。于是
\[
r=\sum_{k=0}^{7}b_k2^k
=b_0\!\times1+b_1\!\times2+\cdots+b_7\!\times128.
\]
符号 $\sum$ 表示“把后面这类项按指定编号逐项相加”；下方 $k=0$ 是起点，上方 7 是终点。$2^k$ 是第 $k$ 位的权重。

把全图每个像素的同一位取出来，就得到一个位平面；显示时通常把位值 0 显示为黑，把位值 1 显示为白。若想用四则运算提取第 $k$ 位：先除以 $2^k$，向下取整，再除以 2 取余数。例如提取 173 的第 5 位：$173/32=5.40625$，向下取整得 5，5 除以 2 的余数为 1。

只保留 173 的最高四位，相当于保留 $128+32=160$，舍掉了 13。高位通常决定较大的明暗结构，低位影响较细的数值变化；\textbf{低位不一定全是噪声}，也可能含有真实微弱细节。是否丢弃，要看任务允许多大误差。

\subsection{直方图 从数像素开始}

\textbf{灰度直方图（gray-level histogram）}回答“每种灰度有多少像素”。它不记录像素在哪里。横轴是灰度，纵轴是像素数，每个灰度对应的计数格也称\textbf{组或箱（bin）}。

先看一幅只有 10 个像素、允许灰度为 $0,1,2,3$ 的小图，其灰度依次为
\[
0,\ 0,\ 1,\ 1,\ 1,\ 1,\ 1,\ 2,\ 3,\ 3.
\]
数一数：0 出现 2 次，1 出现 5 次，2 出现 1 次，3 出现 2 次。因此直方图计数是 $2,5,1,2$。把灰度为 $k$ 的个数记为 $n_k$，把总像素数记为 $n$，本例有 $n=10$。

\textbf{归一化直方图（normalized histogram）}把每个计数除以总数。这个比值可以看作：从图中等可能地随机选一个像素，它恰好为该灰度的\textbf{概率（probability）}。无需先学概率论，记住“这一类有多少，除以总共有多少”即可：
\[
p_k=\frac{n_k}{n}.
\]
$p_k$ 是灰度 $k$ 所占的比例。例如灰度 1 的比例是 $5/10=0.5=50\%$。所有比例相加必须等于 1，因为每个像素恰好属于一个灰度。

\textbf{累积分布函数（cumulative distribution function，CDF）}进一步回答“灰度不大于 $k$ 的像素，占全部的多少”。所谓累计，就是从最左边一直加到当前位置。把累计比例记为 $C_k$，累计个数记为 $N_k$：
\[
N_k=n_0+n_1+\cdots+n_k,\qquad
C_k=\frac{N_k}{n}=p_0+p_1+\cdots+p_k.
\]
例如“不大于 1”包含灰度 0 和 1，两类共 $2+5=7$ 个，所以 $C_1=7/10=0.7$，不是 $5/10$。

\begin{center}
\begin{tabular}{ccccc}
\toprule
灰度 $k$ & 计数 $n_k$ & 概率 $p_k$ & 累计个数 $N_k$ & 累计比例 $C_k$\\
\midrule
0 & 2 & $2/10=0.2$ & 2 & $2/10=0.2$\\
1 & 5 & $5/10=0.5$ & $2+5=7$ & $7/10=0.7$\\
2 & 1 & $1/10=0.1$ & $7+1=8$ & $8/10=0.8$\\
3 & 2 & $2/10=0.2$ & $8+2=10$ & $10/10=1$\\
\bottomrule
\end{tabular}
\end{center}

\begin{center}
\begin{tikzpicture}[x=1.1cm,y=.55cm,>=Stealth,font=\small]
\begin{scope}
\draw[->] (-.5,0)--(4,0) node[right] {$k$};
\draw[->] (-.5,0)--(-.5,6.2) node[above] {计数 $n_k$};
\foreach \k/\v in {0/2,1/5,2/1,3/2}{
\draw[fill=black!15] (\k-.28,0) rectangle (\k+.28,\v);
\node[above] at (\k,\v) {\v};\node[below] at (\k,0) {\k};}
\end{scope}
\begin{scope}[xshift=8cm,y=3cm]
\draw[->] (-.5,0)--(4,0) node[right] {$k$};
\draw[->] (-.5,0)--(-.5,1.15) node[above] {累计比例 $C_k$};
\draw[dotted] (-.5,1)--(3.5,1);\node[left] at (-.5,1) {1};
\draw[very thick] (0,.2)--(1,.2)--(1,.7)--(2,.7)--(2,.8)--(3,.8)--(3,1)--(3.5,1);
\foreach \k/\v in {0/.2,1/.7,2/.8,3/1}{\fill (\k,\v) circle (1.7pt);\node[above left] at (\k,\v) {\v};\node[below] at (\k,0) {\k};}
\end{scope}
\end{tikzpicture}
\end{center}

左图每根柱子单独计数；右图不断把新柱子的比例加进来。因此 CDF 只能上升或保持不变，不能下降，最右端必为 1。某灰度没有像素，该处累计值就不增加。\textbf{高柱子对应累计曲线较大的跳跃}，这是后面均衡化会放大某些小灰度差的原因。

\subsubsection*{直方图能告诉什么 不能告诉什么}

直方图偏左，说明低灰度像素占比大，图像整体常显得偏暗；偏右则常偏亮。集中在窄范围，通常说明可利用的灰度跨度较小；分布很宽，说明存在跨度较大的灰度。但少数黑点和白点就能把范围撑满，不能据此断言每个局部都清楚。

如果把同一批像素打乱位置，直方图不会改变，但一张人像可以被打乱成杂乱斑点。因此直方图\textbf{没有空间结构信息}，不能单独判断清晰度、目标位置、噪声种类，也不能证明图中恰好只有两个物体。

\subsection{直方图均衡 为什么累计比例能变成新灰度}

\textbf{直方图均衡化（histogram equalization，HE）}用输入灰度的累计比例来决定输出灰度。直觉上，它把“我排在全部像素的什么位置”变成“我应该落在输出灰度范围的什么位置”。

假设累计比例为 0.7，表示有 70\% 的像素不比当前灰度亮；让这个灰度落在输出范围约 70\% 的位置，就是均衡映射的基本想法。对于 $L$ 个输出灰度级，本节采用的离散公式为
\[
s_k=\operatorname{round}\!\bigl((L-1)C_k\bigr).
\]
$s_k$ 表示原灰度 $k$ 应换成的新灰度。$C_k$ 是 0 到 1 的累计比例，乘 $L-1$ 把它扩大到合法灰度范围，再取整得到可保存的整数。

刚才的 10 像素例子有 $L=4$，因此乘数为 3：
\begin{center}
\begin{tabular}{ccccc}
\toprule
输入 $k$ & $C_k$ & $3C_k$ & 取整输出 $s_k$ & 这些像素搬到哪里\\
\midrule
0 & 0.2 & 0.6 & 1 & 2 个全部移到 1\\
1 & 0.7 & 2.1 & 2 & 5 个全部移到 2\\
2 & 0.8 & 2.4 & 2 & 1 个也移到 2\\
3 & 1 & 3 & 3 & 2 个保持在 3\\
\bottomrule
\end{tabular}
\end{center}
输出灰度 $0,1,2,3$ 的计数为 $0,2,6,2$。灰度 2 收到原灰度 1 的 5 个像素和原灰度 2 的 1 个像素，所以为 $5+1=6$。\textbf{均衡化不是把像素重新放到空间位置，而是把它们的灰度重新编号。}

为什么名字叫“均衡”，结果却不平均？在灰度连续、分布满足相应条件的理想模型中，用连续 CDF 变换会产生均匀分布；实际图像有有限个整数灰度，同值像素又被同一个规则整体映射，不能照搬这个理想结论。上面的例子反而出现了更高的柱子，正好提醒我们：\textbf{HE 不保证每根柱子更接近平均高度，也不保证视觉效果一定改善。}

\textbf{软件结果不同时先查公式。}有的实现还会减去第一个非零累计值，再缩放到全范围；取整规则也可能不同。本节原题 E1 明确采用 $\operatorname{round}((L-1)C_k)$，计算中不要擅自换成另一种版本。常数图像等特殊输入还需要单独约定，不能把所有实现都当成一模一样。

\subsubsection*{进阶阅读 连续公式中的积分到底在累计什么}

原讲义还给出连续公式。若暂时没有学过积分，可以先把它看作“柱子无限变窄时，累计柱子面积”的写法，不需要先会求复杂积分。

把连续输入记为 $r$，并先归一化到 $[0,1]$。\textbf{概率密度函数（probability density function，PDF）}记为 $p_R(r)$，是描述灰度分布的密度曲线。离散图像的 $p_k$ 是某一个灰度的概率；连续情形中，\textbf{一个区间下方的面积}才是落在该区间的概率，曲线高度本身不是单点概率。曲线总面积为 1。

从 0 累计到当前灰度 $r$ 的面积就是 CDF：
\[
s=T(r)=\int_0^r p_R(w)\,\mathrm{d}w.
\]
弯曲的符号 $\int$ 表示积分，也就是连续累计；下方 0 是开始位置，上方 $r$ 是结束位置；$w$ 是累计过程中走过的临时灰度；$p_R(w)$ 是那个位置的密度；$\mathrm{d}w$ 表示极小的一段灰度宽度。可把 $p_R(w)\mathrm{d}w$ 想成一根极细小柱的面积，再把所有小面积加起来。输出 $s$ 在 0 到 1 之间。

例如输入在 $[0,1]$ 上本来就均匀，密度为 1；从 0 到 $r$ 的面积是“高 1 乘宽 $r$”，于是 $T(r)=r$，无需改变。若满足连续分布条件，CDF 变换后的 $s$ 在 $[0,1]$ 上均匀，写成 $p_S(s)=1$；$p_S$ 表示输出的密度。若输出范围改为 $[0,L-1]$，需在 CDF 外乘 $L-1$，均匀密度也相应为 $1/(L-1)$，不能仍写成 1。

实际数字图像要把积分换成累计和，再量化成整数；这一步破坏了连续结论所需的条件，所以不能要求实际输出逐级等高。

\subsubsection*{原题 E1 完整计算 64 乘 64 图像的均衡映射}
\label{original:E1}

\textbf{题目。}一幅 $64\times64$ 的图像只有 8 个允许的灰度级 $0,1,\ldots,7$，其像素计数如下，求均衡后的映射关系与新直方图。
\begin{center}
\begin{tabular}{c|rrrrrrrr}
\toprule
灰度 & 0&1&2&3&4&5&6&7\\
计数 &790&1023&850&656&329&245&122&81\\
\bottomrule
\end{tabular}
\end{center}

\textbf{第 1 步，确认两个不同的数量。}总像素数 $n=64\times64=4096$，用于求比例；灰度级数 $L=8$，所以输出最大值 $L-1=7$，用于缩放。4096 是“有多少个像素”，8 是“每个像素允许取几种数”，不能互换。

先验算计数总和：$790+1023+850+656+329+245+122+81=4096$，与图像尺寸一致。

\textbf{第 2 步，从左到右求累计个数。}第 0 项为 790；第 1 项是 $790+1023=1813$；第 2 项再加 850 得 2663。每一步都只需在前一累计数上加当前计数。最后必须得到 4096。

\textbf{第 3 步，用精确计数算比例和输出。}比如输入灰度 2：累计个数 2663，累计比例为 $2663/4096\approx0.650146$，乘 7 得 $4.551025$，四舍五入后输出 5。表中的小数用于展示，真正决定取整的是精确分数 $7N_k/4096$。

\begin{center}
\small
\begin{tabular}{rrrrrrr}
\toprule
$k$ & $n_k$ & $N_k$ & $p_k=n_k/4096$ & $C_k=N_k/4096$ & $7C_k$ & $s_k$\\
\midrule
0&790&790&0.192871&0.192871&1.350098&1\\
1&1023&1813&0.249756&0.442627&3.098389&3\\
2&850&2663&0.207520&0.650146&4.551025&5\\
3&656&3319&0.160156&0.810303&5.672119&6\\
4&329&3648&0.080322&0.890625&6.234375&6\\
5&245&3893&0.059814&0.950439&6.653076&7\\
6&122&4015&0.029785&0.980225&6.861572&7\\
7&81&4096&0.019775&1.000000&7.000000&7\\
\bottomrule
\end{tabular}
\end{center}

不要先把概率都保留两位小数再累计。粗略四舍五入可能把某一结果推过半整数分界，改变最终映射。本题恰好得到与原讲义相同的整数映射，不代表提前舍入总是安全。

\textbf{第 4 步，写成灰度对照表。}
\[
0\to1,\quad 1\to3,\quad 2\to5,\quad
3\to6,\quad 4\to6,\quad
5\to7,\quad 6\to7,\quad 7\to7.
\]
箭头表示“每个输入灰度换成哪个输出灰度”。例如原图左上角若是灰度 4，就把该位置写成 6；不需要移动那个像素。

\textbf{第 5 步，按输出灰度合并计数。}输出 6 接收输入 3 和 4，所以计数为 $656+329=985$。输出 7 接收输入 5、6、7，所以计数为 $245+122+81=448$。没人映射到 0、2、4，它们的计数就是 0。

\begin{center}
\begin{tabular}{c|rrrrrrrr}
\toprule
输出灰度 &0&1&2&3&4&5&6&7\\
新计数 &0&790&0&1023&0&850&985&448\\
\bottomrule
\end{tabular}
\end{center}
核查：$790+1023+850+985+448=4096$，总像素没有增减。若总数不对，通常是合并时漏了一组，或把灰度值误当作像素数来加。

\begin{center}
\begin{tikzpicture}[x=.66cm,y=.0034cm,>=Stealth,font=\small]
\begin{scope}
\draw[->] (-.6,0)--(8,0) node[right] {灰度};
\draw[->] (-.6,0)--(-.6,1210) node[above] {计数};
\node[above] at (3.5,1200) {E1 输入直方图};
\foreach \k/\v in {0/790,1/1023,2/850,3/656,4/329,5/245,6/122,7/81}{\draw[fill=black!15] (\k-.32,0) rectangle (\k+.32,\v);\node[below] at (\k,0) {\k};}
\draw[dotted] (-.6,1023)--(7.5,1023);\node[left] at (-.6,1023) {1023};
\end{scope}
\begin{scope}[xshift=7.8cm]
\draw[->] (-.6,0)--(8,0) node[right] {灰度};
\draw[->] (-.6,0)--(-.6,1210) node[above] {计数};
\node[above] at (3.5,1200) {E1 输出直方图};
\foreach \k/\v in {0/0,1/790,2/0,3/1023,4/0,5/850,6/985,7/448}{\draw[fill=black!15] (\k-.32,0) rectangle (\k+.32,\v);\node[below] at (\k,0) {\k};}
\draw[dotted] (-.6,1023)--(7.5,1023);
\end{scope}
\end{tikzpicture}
\end{center}

\textbf{第 6 步，解释结果而非只报数字。}输出占用的灰度级从 8 个变为 5 个，某些级合并、某些级空缺。若严格均匀，每一级应有 $4096/8=512$ 个像素；实际有 0、790、1023 等，显然不均匀。输入原本就有像素覆盖 0 到 7，所以也不能声称这道题“把原有范围变得更宽”。它实现了规定的累计映射，至于是否更适合观看或识别，还要检查图像和任务。

\subsubsection*{原题 B3 为什么均衡后不是完全均匀}
\label{original:B3}

\textbf{第一步，区分理想模型与数字图像。}连续 CDF 均衡的结论有连续分布等条件；数字图像只有有限个灰度级，需取整。

\textbf{第二步，抓住同值像素必须一起变换。}E1 中 1023 个输入灰度 1 的像素都被映射成 3，不能由这张单值灰度表把它们拆成若干份分配到不同灰度。所以只要输入某个计数很大，该组就不会被平均拆散。

\textbf{第三步，说明取整造成合并。}输入 3 与 4 的未取整输出分别约为 5.672 和 6.234，取整后都变为 6，于是计数相加。出现空级和合并级是正常结果。

\textbf{完整结论：}标准离散 HE 使用确定性的灰度映射，像素同值成组搬移，且输出要量化到整数级，因此一般不会得到完全均匀的直方图。它也不保证按某个均匀性指标一定更接近均匀；不能把“名字叫均衡”当作结果保证。

\subsection{直方图匹配 让图像接近指定的灰度分布}

\textbf{直方图规定化或匹配（histogram specification or matching）}是在给定目标分布后，让输出的灰度分布尽量接近该目标。均衡化使用固定的累计映射思路；匹配允许目标偏暗、偏亮、单峰或其他指定形状。

先理解“相同累计位置”：若源灰度 1 的累计比例为 0.7，就去目标分布中找累计比例接近 0.7 的灰度。它们都处在各自分布大约 70\% 的位置，把两者对应起来即可。

源灰度记为 $r$，它的 CDF 记为 $T(r)$；目标灰度记为 $z$，目标 CDF 记为 $G(z)$。先求 $s=T(r)$，再寻找满足 $G(z)\approx s$ 的目标灰度。在可逆的理想连续情形，可简记为
\[
z=G^{-1}(T(r)).
\]
这里 $G^{-1}$ 表示“沿着函数反查输入”，\textbf{不是 $1/G$}。数字图像的 CDF 常有台阶和平台，通常没有普通意义上的唯一逆函数，需要明确近似匹配规则。

\textbf{一个可以手算的离散例子。}仍用前面 10 像素小图，源累计比例是 $0.2,0.7,0.8,1$。希望目标灰度 $0,1,2,3$ 的概率为 $0.1,0.2,0.3,0.4$，目标累计比例便为 $0.1,0.3,0.6,1$。

本例采用“选累计比例差的绝对值最小者，若并列则选较低灰度”。绝对值 $|a-b|$ 表示两个数相差多远，例如 $|0.7-0.6|=0.1$。
\begin{center}
\begin{tabular}{cccl}
\toprule
源灰度 & 源累计值 & 所选目标灰度 & 比较过程\\
\midrule
0&0.2&0&到 0.1 和 0.3 都差 0.1，并列选低灰度\\
1&0.7&2&到目标累计值 0.6 最近，差 0.1\\
2&0.8&2&到 0.6 和 1 都差 0.2，并列选低灰度\\
3&1&3&目标累计值恰好为 1\\
\bottomrule
\end{tabular}
\end{center}
得到映射 $0\to0,1\to2,2\to2,3\to3$，输出计数为 $2,0,6,2$，并不等于目标计数 $1,2,3,4$。原因仍是同灰度像素成组变换，且匹配位置有限。另一种常用规则是选择“第一个累计值不小于源累计值”的灰度；它可能得到不同答案，答题时应说明规则。

\textbf{学到这里就够用的认识：}匹配把两个 CDF 对应起来，不要求先真正保存一幅均衡图。理论上它与一维分布的单调传输有关，但一般离散图像的单值灰度表未必能精确实现任意目标分布。对入门计算，不必先学习最优传输理论。

\subsection{读直方图的原题 先描述 再选方法}

\subsubsection*{原题 B1 比较四幅图像的明暗和对比度}
\label{original:B1}

\textbf{题目。}比较题中 a、b、c、d 四个直方图对应图像的明暗与对比度，指出哪些需要处理，以及可选的方法。

\textbf{图源说明。}原讲义保留了题目和文字答案，但没有可核对的原始四幅题图。下面是为讲解重新绘制的\textbf{教学示意图}，字母只代表本页四种情况，不能用它断言缺失原图的精确形状或编号关系。

\begin{center}
\begin{tikzpicture}[x=.018cm,y=.47cm,>=Stealth,font=\small]
\begin{scope}
\draw[->] (0,0)--(275,0) node[right] {$r$};\draw[->] (0,0)--(0,5.2) node[above] {计数};
\foreach \k/\v in {20/1,35/2.7,50/4.3,65/3.5,80/1.5}{\draw[fill=black!15] (\k-7,0) rectangle (\k+7,\v);}
\node[below] at (0,0) {0};\node[below] at (255,0) {255};\node[above] at (135,5.1) {a 暗部集中};
\end{scope}
\begin{scope}[xshift=7.8cm]
\draw[->] (0,0)--(275,0) node[right] {$r$};\draw[->] (0,0)--(0,5.2) node[above] {计数};
\foreach \k/\v in {180/1,195/2.8,210/4.4,225/3.5,240/1.8}{\draw[fill=black!15] (\k-7,0) rectangle (\k+7,\v);}
\node[below] at (0,0) {0};\node[below] at (255,0) {255};\node[above] at (135,5.1) {b 亮部集中};
\end{scope}
\begin{scope}[yshift=-4.5cm]
\draw[->] (0,0)--(275,0) node[right] {$r$};\draw[->] (0,0)--(0,5.2) node[above] {计数};
\foreach \k/\v in {100/1.5,115/3.8,130/4.5,145/2.8,160/1}{\draw[fill=black!15] (\k-7,0) rectangle (\k+7,\v);}
\node[below] at (0,0) {0};\node[below] at (255,0) {255};\node[above] at (135,5.1) {c 中部窄范围};
\end{scope}
\begin{scope}[xshift=7.8cm,yshift=-4.5cm]
\draw[->] (0,0)--(275,0) node[right] {$r$};\draw[->] (0,0)--(0,5.2) node[above] {计数};
\foreach \k/\v in {15/1.3,40/2.1,65/2.5,90/2,115/2.8,140/3,165/2.4,190/2.6,215/2.1,240/1.4}{\draw[fill=black!15] (\k-9,0) rectangle (\k+9,\v);}
\node[below] at (0,0) {0};\node[below] at (255,0) {255};\node[above] at (135,5.1) {d 较宽范围};
\end{scope}
\end{tikzpicture}
\end{center}

\textbf{第一步，找横向位置。}a 的主要像素靠近黑端，偏暗；b 靠近白端，偏亮；c 在中部，明暗适中；d 涵盖暗、中、亮范围。

\textbf{第二步，看主要分布的宽度。}a、b、c 都较窄，通常整体对比度不足。d 跨度较宽，有较大的全局灰度范围，但这还不能保证每个局部都清楚。

\textbf{第三步，把方法与问题配对。}a 可以拉伸暗部范围，若还需提亮可试规范化 $\gamma<1$；b 可试 $\gamma>1$ 调暗并温和拉开未饱和的亮部差异；c 可用分段线性拉伸；d 不应仅因为有增强算法就强行处理，应先看实际图像和任务是否已经满意。

\textbf{答题边界。}对本页示意，a、b、c 更值得检查和处理。不能仅凭 d 的宽直方图选择“中值去噪”，因为直方图没有告诉我们它是否有椒盐噪声。缺失原始题图时，原题 a--d 的逐项结论仍需与原图核对。

\subsubsection*{原题 B2 夜间道路的双峰直方图}
\label{original:B2}

\textbf{题目。}夜间道路监控图像的直方图呈“两个高耸的峰，中间接近平坦”，说明什么？怎样增强？

\textbf{先解释数据。}两个峰说明至少有两组灰度范围贡献了较多像素。结合夜间道路场景，一个峰可能来自暗路面或天空，另一个可能来自车灯、反光牌等亮区域；这是根据场景作出的合理解释，不能只凭峰的个数确定物体种类。若题目的“平坦”指接近横轴的低谷，说明中间灰度较少；若只是高度不变，则还要看平台的具体高度。\textbf{不能据此认定细节已丢失}。

\textbf{再明确目标。}若只想检测亮灯位置，可以考虑分割亮区域；若想看清暗处行人，重点应是暗部可见性；若需要保持监控画面自然，则还要控制明暗关系和噪声。

\textbf{然后给出方案。}先检查暗区噪声与亮区饱和情况。若照明不均且暗区有可用细节，可试温和 CLAHE；若整体只是偏暗，可试规范化 $\gamma<1$ 或分段曲线，并检查车灯附近是否失去层次。需要去噪时，应按噪声特点选择方法，不能默认所有夜景都用中值滤波。

\textbf{最后说明风险。}全局 HE、CLAHE 和伽马都可能把原有噪声变得更明显。CLAHE 的限制机制降低过度增强风险，却不是“不放大噪声”的保证。完全饱和的车灯内部，也不能靠这些灰度变换恢复真实纹理。

\subsubsection*{原题 C1 根据暗人物与亮投影屏幕画直方图}
\label{original:C1}

\textbf{题目。}左侧是较暗的人物与背景，右侧是一块较亮的投影屏幕。画出直方图示意，并提出增强建议。

\textbf{图源说明。}原始题图未提供，因此无法知道峰的精确灰度、宽度、高度或面积。以下只根据文字场景新画示意，\textbf{不是原题图像的真实统计结果}。原文中的具体峰位数字缺乏可核对图像，本节不把它们当作事实。

\begin{center}
\begin{tikzpicture}[x=.042cm,y=.62cm,>=Stealth,font=\small]
\draw[->] (0,0)--(272,0) node[right] {灰度 $r$};
\draw[->] (0,0)--(0,5.2) node[above] {像素数};
\draw[thick,fill=black!12] (0,0)--(6,.2)--(12,1.4)--(20,2.9)--(30,3.6)--(40,3.1)--(52,1.6)--(66,.4)--(90,.15)--(130,.12)--(163,.15)--(182,.5)--(192,2.5)--(199,4.2)--(205,3.2)--(216,.5)--(230,.15)--(255,0)--cycle;
\node[align=center] at (60,4.5) {较暗人物与背景\\低灰度的一组像素};
\node[align=center] at (205,5) {较亮屏幕\\高灰度的一组像素};
\node[below] at (0,0) {0};\node[below] at (255,0) {255};
\end{tikzpicture}
\end{center}

\textbf{作图第 1 步：}先画灰度横轴和计数纵轴。横轴左黑右白；纵轴越高，说明该灰度的像素越多。

\textbf{作图第 2 步：}在较暗一侧画一组隆起，在较亮一侧画另一组隆起。屏幕若亮度比较接近，可能形成较窄峰；人物与背景的材质和照明不一，暗峰可能更宽。这些是条件性的解释，而非精确测量。

\textbf{作图第 3 步：}不要把“图像左侧的人物”理解成直方图必然左侧。人物的峰在左边，是因为\textbf{人物较暗}，而不是因为人物位于画面左边。直方图不保存空间位置。

\textbf{方法建议：}若目的是看清人物，可在暗部使用有保护亮部的分段曲线，或尝试温和 CLAHE；若还需读取屏幕文字，应同时检查亮区是否饱和，以及增强后文字笔画是否变差。对不同区域分开处理再平滑过渡也是一种方案。全局 HE 可以作比较基线，但不能在没有测试的情况下保证它或 CLAHE 一定更好。

\subsection{HE 的改进方法 先看它试图解决哪个问题}

学这一组名称时，不必背成“后一代全面替代前一代”的时间线。它们处理的问题并不完全相同：有的关心局部照明，有的减少平均灰度变化，有的约束过度增强，还有的借助光照模型或学习模型。下面只讲机制与边界，不用未经核对的年份帮助记忆。

\subsubsection*{全局 HE 常见的三个局限}

\begin{enumerate}[leftmargin=*]
\item \textbf{小区域可能被大背景主导。}若全图大部分是天空，按全图统计得到的曲线未必适合只占少量像素的车牌。
\item \textbf{噪声可能变得明显。}在平坦暗区，灰度 30、31、32 的差别可能主要是噪声；若它们映射到 40、100、170，噪声就被拉大。
\item \textbf{整体明暗可能改变。}HE 没有强制“输出平均灰度等于输入平均灰度”。想保留夜景感觉或照片自然感时，这种变化可能不合适。
\end{enumerate}

\subsubsection*{AHE 让每个局部按自己的统计来处理}

\textbf{自适应直方图均衡化（adaptive histogram equalization，AHE）}用像素周围的局部直方图来确定映射。最直接的理解是：移动一个窗口，统计窗口内的灰度，为当前位置计算均衡输出。高效实现也可以划分若干小块，再结合相邻区域的映射。

与全局 HE 相比，同为灰度 80 的像素，在暗角和亮处可以得到不同输出，因为周围的灰度分布不同。它因此能适应空间上的照明差异。

但局部几乎全是灰度 30 时，极少量的 29 和 31 可能只是噪声；AHE 仍可能把这个狭窄分布拉开。窗口越小，局部统计越容易被少量像素影响；窗口过大，又逐渐失去对小范围变化的适应能力。

\subsubsection*{CLAHE 给局部增强加上限制}

\textbf{限制对比度的自适应直方图均衡化（contrast-limited adaptive histogram equalization，CLAHE）}在局部均衡之前限制过高的直方图柱子，降低由尖峰引起的过强灰度拉伸。常见分块实现包含四步：
\begin{enumerate}[leftmargin=*]
\item \textbf{分块并计数。}把图像划分为小块，每块统计自己的直方图。
\item \textbf{裁剪并再分配。}设置\textbf{裁剪限值（clip limit）}，超过限值的计数先截下来，再按实现规则分配给各灰度级。截下的是用于构建映射的统计量，不是删除那些像素。
\item \textbf{累计并建映射。}对修改后的直方图求 CDF，为每个小块建立灰度对照表。
\item \textbf{在相邻映射间插值。}处理一个像素时，结合周围小块映射的结果，使过渡平滑。常用\textbf{双线性插值（bilinear interpolation）}，本质上按水平和竖直距离加权平均。
\end{enumerate}

\begin{center}
\begin{tikzpicture}[>=Stealth,font=\small]
\draw (0,0) rectangle (4.4,3.6);\draw (2.2,0)--(2.2,3.6);\draw (0,1.8)--(4.4,1.8);
\node[align=center] at (1.1,2.7) {块 A\\局部映射 $T_A$};
\node[align=center] at (3.3,2.7) {块 B\\局部映射 $T_B$};
\node[align=center] at (1.1,.9) {块 C\\局部映射 $T_C$};
\node[align=center] at (3.3,.9) {块 D\\局部映射 $T_D$};
\fill (2.25,1.85) circle (2pt);\node[above right,fill=white,inner sep=1pt] at (2.25,1.85) {$P$};
\node[align=center] at (2.2,-.5) {像素 $P$ 结合相邻块的映射\\不是给每块强行贴一张不同亮度的图};
\draw[->,thick] (4.8,1.8)--(5.6,1.8);
\draw[->] (6,0)--(11.2,0) node[right] {灰度};\draw[->] (6,0)--(6,3.7) node[above] {计数};
\foreach \x/\h in {6.5/.6,7.2/1.1,7.9/3.3,8.6/2.9,9.3/1.3,10/.7}{\draw[fill=black!12] (\x-.25,0) rectangle (\x+.25,\h);}
\draw[dashed,thick] (6,2)--(10.8,2);\node[right,align=left] at (10.8,2) {裁剪\\限值};
\draw[pattern=north east lines] (7.65,2) rectangle (8.15,3.3);
\draw[pattern=north east lines] (8.35,2) rectangle (8.85,2.9);
\node[align=center] at (8.6,-.5) {斜线部分重新分配\\再用修改后的计数求 CDF};
\end{tikzpicture}
\end{center}

\textbf{一遍裁剪的算术例子。}局部四个灰度的计数是 $[1,1,6,0]$，共 8 个。若限值为 3，先得到 $[1,1,3,0]$，截下 3 个计数。为演示整数再分配，假设把它们分别加给第 0、1、3 级，得到 $[2,2,3,1]$，总数仍为 8。实际程序可能按更细的规则分配，并在需要时再次处理超限，不能把这一个演示当成所有软件的唯一实现。

这张新表的累计比例是 $[2/8,4/8,7/8,8/8]$，再乘输出最大灰度并取整，就得到该块映射。\textbf{裁剪的是用于求映射的直方图，不意味着输出直方图每根柱子都低于限值。}

\textbf{插值也可以手算。}某像素用四个邻块映射分别会得到 80、100、120、140；若位于四个块中心构成区域的正中央，每个结果权重相同，输出便是 $(80+100+120+140)/4=110$。靠近某个块时，那个块的结果权重更大。图中仅示意四块与像素的关系，完整实现还需规定图像外边缘的处理方式。

\textbf{限制是什么。}CLAHE 可抑制 AHE 的过度噪声放大，并借助插值减轻块状接缝，但不等于去噪器，也不保证无噪声放大、无伪影或识别率上升。块尺寸和限值都要测试。不同软件对限值的单位、归一化方式不同，不能直接照搬另一个软件的参数数值。

\subsubsection*{BBHE 与 DSIHE 按灰度分组以减少亮度漂移}

先定义\textbf{平均灰度（mean gray level）}：把所有像素灰度相加，再除以像素总数，记为 $\mu$。已知直方图时，不必逐个像素加，可用
\[
\mu=\frac{0n_0+1n_1+\cdots+(L-1)n_{L-1}}{n}.
\]
例如前面的 10 像素图平均灰度为 $(0\times2+1\times5+2\times1+3\times2)/10=13/10=1.3$。

\textbf{亮度保持双直方图均衡化（brightness-preserving bi-histogram equalization，BBHE）}以输入均值为依据选分割灰度，把直方图分成较暗和较亮两组，分别在各自子范围内均衡。名称中的“亮度保持”表达设计目标，不是任意数字图像上严格相等的数学保证。

为看清机制，取分割阈值 $t=1$，低组为灰度 0、1，高组为 2、3。低组共有 $2+5=7$ 个像素，在低组内部重新算概率：$2/7,5/7$，不是仍除以 10。低组累计值是 $2/7,1$，映射到输出范围 $[0,1]$。高组共 3 个，累计值是 $1/3,1$，映射到 $[2,3]$。

若采用本节的四舍五入规则，低组输出为 $\operatorname{round}(2/7)=0$ 和 1；高组输出为 $\operatorname{round}(2+1/3)=2$ 和 3。这个特定例子恰好不变，但其他输入未必如此。

\textbf{一个直接反例。}设允许灰度为 0 到 7，输入四个像素为 $0,0,4,4$，均值为 2。按 $[0,2]$ 与 $[3,7]$ 分组，各组只有一个出现的灰度，其组内 CDF 为 1；用未减最小 CDF 的子区间均衡，0 映到 2，4 映到 7，输出为 $2,2,7,7$，均值变为 4.5。这个反例足以否定“BBHE 对所有离散图像精确保均值”。实现对单一灰度子组另有约定时结果可能不同，也更说明结论必须带条件。

\textbf{原讲义中那个均值条件从哪里来。}原文给出 $P(r\le m)=1-m/255$。这不是一般离散图像公式，但可以在特定理想模型中推导：设灰度连续地落在 $[0,255]$，输入均值为 $m$，按 $m$ 分成两组；低组比例为 $q=P(r\le m)$，高组比例便是 $1-q$。再假定处理后低组在 $[0,m]$ 完全均匀、高组在 $[m,255]$ 完全均匀，且两组均非空。

均匀区间的平均值是两端的平均，因此两组输出均值分别为 $m/2$ 和 $(m+255)/2$。全图均值要按组内像素比例加权：
\[
\mu_{\rm out}
=q\frac m2+(1-q)\frac{m+255}{2}
=\frac{m+255-255q}{2}.
\]
若要求它等于输入均值 $m$，两边乘 2 得 $m+255-255q=2m$，移项得 $255q=255-m$，再除以 255，才得到 $q=1-m/255$。所以该条件仅对上述\textbf{连续、子范围完全均匀}的理想模型成立。实际 BBHE 的整数阈值、上下子范围端点、取整和非均匀离散输出，都使这个推导不能直接当成通用定理。

\textbf{等面积双子图像直方图均衡化（dualistic sub-image histogram equalization，DSIHE）}改用中位数附近的灰度分割。\textbf{中位数（median）}是把像素值从小到大排列后处于中间的位置；从 CDF 看，就是累计比例到达约 0.5 的地方。分割后同样分别均衡，目标是让两组像素数尽量接近。

“按灰度分组”与 AHE 的“按空间分块”不同：BBHE 的暗组像素可能散布在整幅图中。另一个细节是：若某一灰度占比很大，不拆开同值像素就可能无法把像素数恰好二等分，不能把 DSIHE 的两组说成一定严格等大。

\subsubsection*{RMSHE 与 MMBEBHE 进一步控制平均亮度变化}

\textbf{递归均值分割直方图均衡化（recursive mean-separate histogram equalization，RMSHE）}把 BBHE 的分割继续做下去。先按全图均值分两组；然后在每组内部重新求均值，再各分两组。这个“对得到的子问题重复同样操作”的过程叫\textbf{递归（recursion）}。

若每组都能有效分裂，分 1 层最多有 2 组，分 2 层最多 4 组，分 $d$ 层最多 $2^d$ 组，$d$ 是分裂层数。最后每组只在自己的灰度范围内均衡。分组细化可限制灰度移动幅度，但也可能使增强变弱；数字图像还有空组、单值组和取整问题，不能宣称递归层数越多视觉效果必然越好。

\textbf{最小平均亮度误差双直方图均衡化（minimum mean brightness error bi-histogram equalization，MMBEBHE）}尝试用使亮度变化较小的分割阈值。先定义\textbf{绝对平均亮度误差（absolute mean brightness error，AMBE）}：输入均值记为 $\mu_{\rm in}$，输出均值记为 $\mu_{\rm out}$，则
\[
\operatorname{AMBE}=|\mu_{\rm out}-\mu_{\rm in}|.
\]
竖线表示绝对值。例如输入均值 80、输出均值 95，误差为 $|95-80|=15$；输出均值 70 时误差为 10，不是负 10。

MMBEBHE 的核心是比较候选分割阈值对应的亮度误差，选择较优阈值，再分别均衡。理解时可以想象“尝试几种切法，比较输出均值，选择误差最小的切法”；实际算法可用高效计算避免生成所有完整输出图。\textbf{最小只针对所比较的阈值、误差定义和算法模型，不代表图像各方面最好，也不保证误差为零。}一幅完全不处理的图 AMBE 就为零，所以这个指标不能单独衡量增强质量。

\subsubsection*{WTHE 先修改直方图的权重}

\textbf{加权阈值直方图均衡化（weighted thresholded histogram equalization，WTHE）}先对直方图概率设置阈值并作加权，再用修改后的分布计算映射。它试图避免少数很高的柱子过度控制映射，也控制低频数灰度对增强的影响。

这里的“加权”是改变各灰度在构建映射时的贡献，不是把图像某些像素变重。修改后通常还要重新归一化，使用于累计的概率总和为 1。它与 CLAHE 都可能限制直方图的影响，但\textbf{WTHE 本身不等于空间分块}；具体权重和阈值需看所采用的版本。入门时先记“先改分布，再求 CDF”，不必背没有说明参数含义的公式。

\subsubsection*{Retinex 与 LIME 从光照和物体反射来理解暗图}

\textbf{视网膜皮层理论及其图像增强方法（Retinex）}常用一个简化模型理解观察到的图像：同一个物体，照得亮就更亮，照得暗就更暗；图像的明暗同时受照明和物体表面反射特性影响。

把位置记为 $(x,y)$，观察强度记为 $I(x,y)$，照明分量记为 $\ell(x,y)$，反射分量记为 $R(x,y)$，便有模型
\[
I(x,y)=\ell(x,y)R(x,y).
\]
这里的 $\ell$ 是小写字母 ell，表示照明，不是前文灰度级数 $L$。相乘表示照明与反射共同决定强度；例如反射分量为 0.5，照明从 0.2 变为 0.8，观察强度便从 0.1 变为 0.4。

Retinex 类方法估计或消除部分照明影响，以改善可见性和颜色表现。常见方法借助局部平滑估计背景照明，并在对数域中处理乘积关系。但单幅图像通常不能唯一确定照明和反射：$0.2=0.4\times0.5=0.5\times0.4$。因此必须引入平滑、边缘或颜色等假设，且可能出现光晕、噪声增强或颜色失真。不能把“物理分解”误解成自动获得真实光照。

\textbf{基于照明图估计的低照度增强（low-light image enhancement via illumination map estimation，LIME）}以照明图估计为重点。它可从彩色图每个像素的红、绿、蓝三个通道中的最大值建立初始照明估计，再利用空间结构对照明图进行细化，根据估计结果增强暗处。所谓\textbf{通道（channel）}，就是分别记录某种颜色分量的一张数值图。

直观上，同样的观察灰度若处于明显较暗的照明区域，就需要不同于亮区的调整。LIME 的光照估计与空间结构有关，并不是一条全图共用的普通伽马曲线。除以很小的照明估计时，原有噪声也可能被放大，因此照明正则化、数值保护及必要的去噪都很重要。

\subsubsection*{Zero-DCE 让模型预测局部曲线参数}

\textbf{零参考深度曲线估计（zero-reference deep curve estimation，Zero-DCE）}用神经网络预测增强曲线的参数。\textbf{神经网络（neural network）}在这里可先理解为一个通过训练学会“根据图像给出参数”的函数，不需要先掌握其全部内部结构。

其基础增强曲线可写成
\[
v=u+\alpha u(1-u).
\]
$u$ 是 0 到 1 的归一化输入，$v$ 是输出，$\alpha$ 是控制弯曲程度的参数。先算 $1-u$，再与 $u$ 相乘，乘以 $\alpha$，最后加回 $u$。当 $u=0.2,\alpha=0.5$ 时，$v=0.2+0.5\times0.2\times0.8=0.28$，于是变亮。若 $\alpha=-0.5$，则得到 0.12，变暗。当 $-1\le\alpha\le1$ 时，曲线在 $0\le u\le1$ 内保持输出范围并且单调不减，即输入较大时输出不会较小；多次重复曲线调整可产生更灵活的变换。

关键区别是：参数可以随位置、通道及迭代步骤变化，由网络结合图像内容预测。因此最终曲线作用于像素，不代表整个算法只看该像素自己的输入；\textbf{它不是一条固定的全局 $T(r)$}。

“零参考”表示训练不需要正常曝光参考图作为监督目标，并不表示“不需要训练数据”“不需要训练”或“任何暗图都能无损恢复”。模型仍需要图像和约束增强效果的损失函数；噪声、极暗细节和饱和信息仍有困难。

\subsubsection*{用一张对照表整理方法}
\begin{center}
\small
\begin{tabularx}{\linewidth}{@{}p{2.4cm}p{4.4cm}X@{}}
\toprule
方法 & 主要改变什么 & 先记住的限制\\
\midrule
HE & 全图共用累计映射 & 可改均值、放大噪声；离散输出不必均匀\\
AHE & 根据局部统计选映射 & 平坦局部易过度增强\\
CLAHE & 局部直方图限幅，再平滑结合映射 & 仍需调参，不能保证去噪\\
BBHE / DSIHE & 按均值 / 中位数分灰度组 & 不同于空间分块，不保证精确保均值\\
RMSHE & 反复按各组均值细分 & 细分更多可能减弱增强\\
MMBEBHE & 按亮度误差选分割阈值 & 均值误差小不等于整体质量高\\
WTHE & 修改直方图权重与阈值 & 具体参数影响增强和细节\\
Retinex、LIME & 估计或调整照明的影响 & 分解依赖假设，可能有光晕与噪声\\
Zero-DCE & 网络预测空间变化的曲线参数 & 仍有训练、数据与泛化要求\\
\bottomrule
\end{tabularx}
\end{center}

\subsection{原题 G3 低对比度车牌在 HE CLAHE 和伽马中怎么选}
\label{original:G3}

\textbf{题目。}对一张低对比度车牌图像，要在 HE、CLAHE、伽马变换中选一种，服务于后续字符识别。你选哪个，为什么？

\textbf{第 1 步，确认目前不知道什么。}“低对比度”本身不足以唯一决定方法，还需知道它是整体偏暗、局部照明不均、噪声重，还是已经模糊或饱和。题目不给图像时，要把选择的前提写出来，不能假装知道车牌实际照明。

\textbf{第 2 步，给出有条件的主方案。}若车牌受到明显不均匀照明，可把温和 CLAHE 作为优先测试方案。它根据局部统计增强字符与背景的差异，并限制过强拉伸；相邻块映射的插值有助于减轻分块接缝。

\textbf{第 3 步，与其他候选公平比较。}若主要问题是整体偏暗且照明均匀，规范化 $\gamma<1$ 可能已经足够，且参数简单。HE 可作为自动全局对照方案；它不利用位置，可能被大背景主导，但也不能断言一定失败。伽马能改变某些灰度区间的对比度，只是不具有 CLAHE 的空间自适应性。

\textbf{第 4 步，写出真正的验收方式。}尽量在车牌区域比较原图、HE、CLAHE 和伽马结果；保持后续识别器和其他处理条件一致，用有正确答案的样例检查字符或整牌识别正确率，同时查看笔画断裂、粘连、噪声与饱和。不能用“看起来更亮”代替“识别更准确”。

\textbf{可写在答卷上的答案：}若低对比度主要来自局部照明不均，我优先测试 CLAHE，因为它具有局部适应性和限制过度增强的机制；采用适当块尺寸与裁剪限值，并以识别结果验证。若只是均匀的整体偏暗，伽马变换也可能更合适。最终方法应由图像条件和字符识别效果决定，CLAHE 不保证始终最优。

\subsection{本节自查 把公式变成能解释的步骤}

不看前文，试着回答以下问题，再用后面的提示检查。
\begin{enumerate}[leftmargin=*]
\item 8 位图中 $L$ 为什么是 256，而最大灰度是 255？\par
\textbf{检查：}共有 256 个整数，从 0 开始数到 255。
\item 为何 $\gamma=1/2$ 有时算得 8，有时算得 128？\par
\textbf{检查：}对 64 直接开方得 8；对 $64/255$ 开方再乘 255，才是本节保持范围的规范化伽马，约得 128。
\item 灰度 2 的概率与累计概率有何不同？\par
\textbf{检查：}前者只数灰度 2，后者把 0、1、2 都计入，再除以总数。
\item E1 中输出 6 为什么有 985 个像素？\par
\textbf{检查：}输入 3 与 4 都映射成 6，所以把其计数 $656+329$ 相加。
\item 两张图直方图相同，为什么仍可能完全不像？\par
\textbf{检查：}直方图没有保存像素位置和排列结构。
\item CLAHE 比 AHE 多一个限制机制，为什么仍要检查噪声？\par
\textbf{检查：}限制增强幅度只是减少过度放大风险，不等于估计并去掉噪声。
\end{enumerate}

% 本文件独立定义所有所需局部写法；由主稿载入。
\section{代数运算：从多张图得到更可靠的图}
\label{sec:algebra-beginner}
本节展开原讲义模块4.4。先不考虑复杂公式：一幅灰度图像就是一张数值表，每个位置存着一个表示明暗的数。所谓\textbf{图像代数运算（image arithmetic）}，就是把两张或多张图\textbf{对应位置上的数}相加、相减、乘系数。例如，同一位置两次读到98和102，把它们平均便得到100。运算对象是数值；能否得到有意义的图像，取决于这些位置是否真的对应同一个物体点。

\subsection{先学会读一张小数值表}
\textbf{矩阵（matrix）}在这里可以先理解为整齐排列的数值网格，并不要求先学线性代数。横着的一排叫\textbf{行（row）}，竖着的一列叫\textbf{列（column）}。本章统一用$r$表示行号，向下增加；用$c$表示列号，向右增加。$f(r,c)$表示图$f$的第$r$行、第$c$列那个数。这样约定是为了避免把“图上的水平边缘”和“数值沿哪个方向变化”混为一谈。
\begin{center}
\begin{tikzpicture}[x=.78cm,y=.78cm,font=\small]
  \draw[step=1] (0,0) grid (3,3);
  \foreach \x/\y/\v in {.5/2.5/20,1.5/2.5/30,2.5/2.5/40,.5/1.5/50,1.5/1.5/60,2.5/1.5/70,.5/.5/80,1.5/.5/90,2.5/.5/100}
    \node at (\x,\y) {\v};
  \draw[very thick] (1,1) rectangle (2,2);
  \draw[-{Stealth}] (0,3.45)--(3,3.45) node[right] {列号$c$增加};
  \draw[-{Stealth}] (-.45,3)--(-.45,0) node[below] {行号$r$增加};
  \node[anchor=west,align=left] at (4,1.7) {若行、列从1开始计数：\\中心像素$f(2,2)=60$\\上邻居$f(1,2)=30$\\右邻居$f(2,3)=70$};
\end{tikzpicture}
\end{center}
本章写出的数值网格既可能代表\textbf{真实像素值}，也可能代表\textbf{待乘的系数}。两者外形相似，含义不同：像素网格告诉我们“图里有什么”，系数网格告诉我们“怎样计算”。

\subsection{逐像素平均：平均的是同一地点的多次观测}
对静止场景连续拍3张照片。假设已对齐，在某个像素位置观测到：
\[
  g_1=96,\qquad g_2=102,\qquad g_3=102.
\]
算术平均（arithmetic mean）就是“先相加，再除以个数”：
\[
  \bar g=\frac{96+102+102}{3}=100.
\]
对每一个位置都做同样的计算，就得到平均图。若共有$N$帧，记号为
\[
  \bar g(r,c)=\frac{1}{N}\sum_{i=1}^{N}g_i(r,c).
\]
这里$\sum$读作“求和”：把$i=1,2,\ldots,N$对应的项全部加起来。上横线只是提示这是平均值，不是另一个坐标。

为什么可能有用？真实场景若应该是100，噪声使一次观测偏小、另一次偏大；平均时这些\textbf{有正有负的偏差有机会相互抵消}。但这不保证任意3张照片都刚好抵消。若3次都是102，平均仍是102。去噪效果是对许多可能的观测而言的统计规律，而不是每次都精确还原。

\textbf{两种平均不要混淆：}
\begin{itemize}
  \item 多帧平均：固定一个空间位置，把\textbf{不同时间}拍到的数平均。对齐准确、场景静止时，不必牺牲空间细节。
  \item 邻域均值滤波：固定一张图，把\textbf{附近不同位置}的数平均。把邻居借来降噪，也可能把不同物体的亮度混在一起；下一节会具体计算。
\end{itemize}

\subsection{噪声、零均值、方差和标准差分别在说什么}
\textbf{噪声（noise）}是观测中不希望出现的扰动。先用最简单的\textbf{加性噪声模型（additive noise model）}：
\[
  g_i=f+n_i.
\]
$f$是真实值，$n_i$是第$i$次误差，$g_i$是拍到的值。真实值100、观测值96，对应$n_i=-4$；观测值103，对应$n_i=3$。为看清推理，这里省略了固定不变的位置$(r,c)$。

\textbf{期望（expectation）}写成$E[\cdot]$，可以先理解为“在同样条件下重复非常多次后，数值趋向的平均”。\textbf{零均值噪声（zero-mean noise）}满足$E[n_i]=0$，意思是长期没有持续偏亮或偏暗的倾向。它\textbf{不表示每一次噪声都是0}，也不表示有限几次噪声的和一定为0。

只知道平均为0，还无法区分“小抖动”和“大抖动”。以下两种噪声长期平均都为0：
\[
  \text{小抖动：}-2,+2,-2,+2,\ldots;
  \quad\text{大抖动：}-10,+10,-10,+10,\ldots
\]
直接平均正负偏差会抵消，所以先把偏差\textbf{平方}，再求平均。这就是\textbf{方差（variance）}：
\[
  \operatorname{Var}(n)=E\bigl[(n-E[n])^2\bigr].
\]
零均值时简化为$E[n^2]$。小抖动的方差是$(4+4)/2=4$，大抖动的方差是$(100+100)/2=100$。\textbf{标准差（standard deviation）}是方差的正平方根：
\[
  \sigma=\sqrt{\operatorname{Var}(n)}.
\]
于是上面两种噪声的标准差分别为2和10。标准差与像素灰度使用同一单位，比较直观；方差使用该单位的平方，推导更方便。\textbf{高斯噪声（Gaussian noise）}指噪声误差的概率分布呈钟形：在零均值模型中，接近0的小误差较常见，较大的正负误差较少见，并且概率关于0对称。它不是说噪声只能取几个固定值，也不是说所有误差概率相同。均值与标准差分别决定钟形的中心位置与宽度。这里讨论已知噪声模型的理论方差，不能把它与“用有限样本估计方差时为何除以$N-1$”混为一件事。

\subsection{为什么平均$N$帧后，标准差只缩小$\sqrt N$倍}
平均图可以拆成“原信号+平均噪声”：
\[
  \bar g=\frac1N\sum_{i=1}^{N}(f+n_i)
        =f+\underbrace{\frac1N\sum_{i=1}^{N}n_i}_{\bar n}.
\]
若各帧都是零均值噪声，那么
\[
  E[\bar g]=f+\frac1N\sum_{i=1}^{N}E[n_i]=f.
\]
这叫\textbf{无偏（unbiased）}：长期平均对准真实值，并非单次结果必然正确。

再看方差。设每帧噪声方差相同，都是$\sigma^2$。对零均值的$\bar n$，先把平方完整展开：
\begin{align*}
 \operatorname{Var}(\bar n)
 &=E\left[\left(\frac{n_1+\cdots+n_N}{N}\right)^2\right]\\
 &=\frac1{N^2}\left(\sum_{i=1}^{N}E[n_i^2]
       +2\sum_{i<j}E[n_i n_j]\right).
\end{align*}
$\sum_{i<j}$表示把不同帧两两配对，每一对只计一次；前面的2来自展开平方时每对会出现$n_in_j$与$n_jn_i$两项。

\textbf{互不相关（uncorrelated）}在这里的关键作用，是让不同帧之间的交叉项$E[n_i n_j]$为0。\textbf{相互独立（independent）}是更强的条件，具有有限方差的独立噪声必然互不相关；本题只需要互不相关。
\[
  \operatorname{Var}(\bar n)
  =\frac{N\sigma^2}{N^2}=\frac{\sigma^2}{N},
  \qquad
  \sigma_{\bar n}=\sqrt{\frac{\sigma^2}{N}}=\frac{\sigma}{\sqrt N}.
\]
\textbf{分母为什么不是$N$？}方差除以$N$以后，标准差还要开平方；把“方差下降”和“标准差下降”说成同一倍率，是这道题最常见的错误。

\begin{center}
\begin{tabular}{rrrr}
\toprule
平均帧数$N$ & 原标准差 & 平均后方差 & 平均后标准差\\
\midrule
1  & 10 & 100 & 10\\
4  & 10 & 25 & 5\\
25 & 10 & 4 & 2\\
100 & 10 & 1 & 1\\
\bottomrule
\end{tabular}
\end{center}
从4帧增加到25帧，成本增加很多，标准差只是从5降到2。因此平均帧数的收益会逐渐变慢。

\subsection{假设不成立时：偏差、相关噪声和运动不会自动消失}
\begin{enumerate}[label=\arabic*.]
  \item \textbf{固定偏差不会被平均掉。}若每次都额外偏亮5，即$E[n_i]=5$，那么$E[\bar g]=f+5$。零均值假设决定了能否无偏。
  \item \textbf{同步波动不能靠重复拍摄抵消。}若各帧都有完全相同的误差$n_i=z$，则$\bar n=z$，方差根本不下降。一般公式中的交叉项此时不能删除。
  \item \textbf{运动会产生重影。}如果第1帧的某位置是路面，第2帧同位置是行人，平均得到的是两种物体的混合；这不是随机噪声抵消。应先\textbf{配准（registration，即对齐对应物体位置）}，或采用适合运动场景的方法。
  \item \textbf{不能先把结果截断再套原推导。}过早四舍五入、裁剪到$0\sim255$会引入额外误差。通常先用更宽的数值类型求和和平均，最后再做显示转换。
\end{enumerate}

\subsection{E3 原题与完整解答：平均去噪需要多少帧}
\label{original:E3}
\textbf{题目。}对同一静止场景连续拍摄$N$帧并逐像素取平均，若各帧噪声互不相关、零均值、方差为$\sigma^2$，求平均后图像的噪声方差与标准差；若要求噪声标准差降为原来的$1/5$，需要平均多少帧？

\textbf{第1步：明确平均对象。}同一位置的真实值$f$保持不变，各帧$g_i=f+n_i$，所以$\bar g=f+\bar n$，其中$\bar n=N^{-1}\sum_i n_i$。

\textbf{第2步：代入题设。}零均值给出$E[\bar n]=0$；互不相关使交叉项消失。因此
\[
  \boxed{\operatorname{Var}(\bar n)=\frac{\sigma^2}{N}},
  \qquad \boxed{\sigma_{\bar n}=\frac{\sigma}{\sqrt N}}.
\]
\textbf{第3步：把文字要求变成等式。}要求的是\emph{标准差}降到$1/5$，不是方差降到$1/5$：
\[
 \frac{\sigma}{\sqrt N}=\frac{\sigma}{5}
 \ \Longrightarrow\ \sqrt N=5
 \ \Longrightarrow\ \boxed{N=25\text{帧}}.
\]
若要求“至少下降到原来的$1/5$”，应写$N\ge25$，最少25帧。例如原标准差10、原方差100，平均25帧后标准差2、方差4；$2/10=1/5$而$4/100=1/25$，互相核对一致。

\subsection{图像相减：差别会出现，负数也会出现}
\textbf{图像相减（image subtraction）}逐位置计算$d=f-h$。例如一个位置前后分别为100和130，按“后减前”得到30；另一处由100变成70，差值为$-30$。正负号携带“变亮或变暗”的信息。

它可用于比较前后变化、背景差分以及后面的细节提取，但要注意三件事：
\begin{itemize}
  \item 两图应尺寸一致、位置对齐，且曝光变化要受控；否则整幅图的差异未必来自物体运动。
  \item 用可容纳负数的类型保留中间结果。无符号8位数只能存$0\sim255$，直接相减可能发生裁剪或回绕，具体取决于软件。
  \item 想显示差值，可以取绝对值$|d|$观察变化幅度，或把0显示为中灰以区分正负；这都是\textbf{显示规则}，不等于原始差值计算。
\end{itemize}
相减也会增加不相关噪声：两图噪声方差若都是$\sigma^2$，则差值噪声方差为$\sigma^2+\sigma^2=2\sigma^2$。减法强调差异，却不天然具有去噪能力。

\section{空间滤波：让邻居一起决定一个像素}
\label{sec:spatial-filter-beginner}
本节展开原讲义模块4.5。\textbf{空间域（spatial domain）}就是直接在像素网格上操作。\textbf{滤波（filtering）}不是只指“去掉噪声”，也包括平滑、边缘提取和锐化。共同点是：为一个输出位置收集附近信息，按指定规则得到一个新数值。

\subsection{邻域、窗口、核：三个名字对应什么}
\textbf{邻域（neighborhood）}是一组附近像素；\textbf{窗口（window）}规定选取它们的位置与形状。常见$3\times3$窗口含中心及周围8个像素。$3\times3$中间有唯一中心；这就是教学中常用奇数边长的原因。

\textbf{核（kernel）}也叫\textbf{掩模（mask）}或模板。在\textbf{线性滤波（linear filtering）}中，它是一张与窗口同形状的系数表。把系数表覆盖到像素表上，\textbf{对应位置相乘，再把所有乘积相加}，得到一个输出像素。这里不是线性代数里的矩阵乘法：不需要“行乘列”，只需要“同位置乘后相加”。
\begin{center}
\begin{tikzpicture}[x=.62cm,y=.62cm,font=\small]
  \foreach \x/\y/\v in {0/2/10,1/2/20,2/2/30,0/1/40,1/1/50,2/1/60,0/0/70,1/0/80,2/0/90}
   {\draw (\x,\y) rectangle ++(1,1);\node at (\x+.5,\y+.5) {\v};}
  \node at (1.5,3.65) {像素值$P$};
  \node at (4,1.5) {$\longrightarrow$};
  \foreach \x/\y/\v in {5/2/1,6/2/2,7/2/1,5/1/2,6/1/4,7/1/2,5/0/1,6/0/2,7/0/1}
   {\draw (\x,\y) rectangle ++(1,1);\node at (\x+.5,\y+.5) {\v};}
  \node[align=center] at (6.5,3.65) {对应相乘的系数\\最后整体除以16};
  \draw[-{Stealth}] (8.5,1.5)--(10,1.5);
  \draw (10.4,1) rectangle (11.4,2);
  \node at (10.9,1.5) {50};
  \node[align=center] at (10.9,-.45) {9个输入数\\得到1个输出数};
\end{tikzpicture}
\end{center}
以上计算全部展开为
\begin{align*}
 g&=\frac1{16}(1\cdot10+2\cdot20+1\cdot30
       +2\cdot40+4\cdot50+2\cdot60\\
  &\hspace{23mm}+1\cdot70+2\cdot80+1\cdot90)\\
  &=\frac{10+40+30+80+200+120+70+160+90}{16}
   =\frac{800}{16}=50.
\end{align*}
中心50得到较大的系数4，左右上下分别得到系数2，四个角得到系数1。\textbf{系数的大小表示这个邻居的影响有多大}；负系数表示从总和中减去该邻居的贡献。

\subsection{滑动一次只写一个数：输入图和输出图要分开}
算出一个中心的结果以后，把窗口移动到下一个中心，再取新的邻居重复计算。下图中粗框表示两个相邻窗口的位置，并非一次把整块$3\times3$都改成同一个数。
\begin{center}
\begin{tikzpicture}[x=.65cm,y=.65cm,font=\small]
 \draw[step=1,gray] (0,0) grid (6,4);
 \foreach \x in {0,...,5} \foreach \y in {0,...,3}
  {\pgfmathtruncatemacro{\v}{10+\x*5+\y*3}\node at (\x+.5,\y+.5) {\v};}
 \draw[very thick] (0,1) rectangle (3,4);
 \draw[thick,dashed] (1,1) rectangle (4,4);
 \draw[very thick] (1,2) rectangle (2,3);
 \draw[thick,dashed] (2,2) rectangle (3,3);
 \draw[-{Stealth},thick] (1.5,4.5)--(2.5,4.5);
 \node[align=left,anchor=west] at (7,2.1) {实线：先计算这个中心\\虚线：再向右移动一格\\每次读取原输入图\\各结果写到另一张输出图};
\end{tikzpicture}
\end{center}
如果刚算完就覆盖原图，再让后续窗口读到这个新值，后续计算便混合了原值与已修改值；从左向右算和从右向左算甚至可能得到不同结果。普通空间滤波默认\textbf{所有窗口都读取同一份未被本次运算改写的输入}。

操作步骤可以直接照着执行：
\begin{enumerate}[label=\arabic*.]
 \item 准备输出图，并选定窗口大小、核和边界规则。
 \item 把核的中心对准一个输入像素，取出所需邻域。
 \item 对应相乘，再相加；若采用中值滤波，则改为排序取中间数。
 \item 把结果写在输出图的同一中心位置，移动窗口，直到遍历完成。
 \item 最后才按要求取整、裁剪或转换为显示格式。
\end{enumerate}

\subsection{相关与卷积：到底要不要翻转系数表}
\textbf{相关（correlation）}按原样使用核。设核$w(s,t)$中，$s$表示相对中心的行偏移，$t$表示列偏移：负1是上一行或左一列，0是本行或本列，正1是下一行或右一列。对$3\times3$核，
\[
 g(r,c)=\sum_{s=-1}^{1}\sum_{t=-1}^{1}w(s,t)f(r+s,c+t).
\]
两个求和号只是“逐行、逐列把9项加起来”。

\textbf{卷积（convolution）}则在上述乘加之前把核转$180^\circ$，也就是上下翻转并左右翻转：
\[
 g(r,c)=\sum_{s=-1}^{1}\sum_{t=-1}^{1}w(s,t)f(r-s,c-t).
\]
等价做法是把$w(s,t)$换成$w(-s,-t)$，再按相关的方法计算。\textbf{只上下翻转或只左右翻转都不是一般二维卷积所需的$180^\circ$翻转。}

例如取一行像素$[10,20,40]$，核$[1,0,-1]$：
\[
 \begin{array}{ll}
 \text{相关：}&1\cdot10+0\cdot20-1\cdot40=-30,\\
 \text{卷积：}&-1\cdot10+0\cdot20+1\cdot40=30.
 \end{array}
\]
同样的边界变化，输出符号相反。对于$180^\circ$旋转对称的均值核、高斯核、拉普拉斯核，翻转后不变，因此相关与卷积给出相同结果。本章A1、A2的7个核都具有这种对称性，题内计算不受这一区别影响；以后遇到一阶差分核时必须检查。

\subsection{走到图像边缘怎么办：延拓也是算法的一部分}
\textbf{边界延拓／填充（padding）}是为图外不存在的像素约定临时取值。比如图像最左侧像素要找“左邻居”，总得规定这个值是什么。下面用一行$[10,20,30,40]$说明，两边各补2个数，竖线中间才是真实数据。
\begin{center}
\begin{tabular}{ll}
\toprule
规则 & 延拓后的数列\\
\midrule
常数填充（constant，取0） & $0,0\mid10,20,30,40\mid0,0$\\
复制边缘（replicate） & $10,10\mid10,20,30,40\mid40,40$\\
对称反射，重复端点 & $20,10\mid10,20,30,40\mid40,30$\\
反射，不重复端点 & $30,20\mid10,20,30,40\mid30,20$\\
周期延拓（circular） & $30,40\mid10,20,30,40\mid10,20$\\
\bottomrule
\end{tabular}
\end{center}
不同软件对“symmetric”和“reflect”的命名细节可能不同，所以\textbf{看到名字后还要核对是否重复端点}。二维情况按行、列同样延伸，四角也要服从同一规则。

用3点均值计算最左边那个10，可以立刻看出差异：
\[
 \begin{array}{ll}
 \text{补0：}&(0+10+20)/3=10,\\
 \text{复制：}&(10+10+20)/3=40/3\approx13.33,\\
 \text{不重复端点的反射：}&(20+10+20)/3=50/3\approx16.67,\\
 \text{周期：}&(40+10+20)/3=70/3\approx23.33.
 \end{array}
\]
核完全一样，边缘输出仍可不同。周期延拓会把最右端接到最左端，适合本来就周期重复的内容；对普通照片，它可能创造不真实的边界跳变。补0会在亮图周围引入暗边，因此“系数和为1保持常数”在图像边缘还需要\textbf{延拓本身也保持这个常数}。

另一种选择是\textbf{只计算窗口完全落在图内的位置（valid输出）}。例如$5\times5$图使用$3\times3$核、不填充时，只能得到$3\times3$个完整输出。若要输出与输入一样大，通常需要某种填充。

\subsection{先做常数图测试：系数和告诉我们什么}
把窗口内所有像素都设成同一个数$C$，如100。每个数都一样，因此无论窗口如何移动，输出都是
\[
 g=Cw_1+Cw_2+\cdots+Cw_K
   =C\underbrace{(w_1+\cdots+w_K)}_{S},
\]
其中$K$是系数个数，$S$是系数和。这个最简单的测试能回答“均匀背景会发生什么”：
\begin{itemize}
 \item $S=1$：100仍输出100，\textbf{保留常数亮度}。
 \item $S=0$：100输出0，\textbf{去除常数背景}，只对某些变化产生响应。
 \item $S=2$：100输出200，常数背景被放大2倍。
\end{itemize}
以后频域中把不随位置变化的成分叫\textbf{直流分量（DC component）}；这里无需先懂傅里叶变换，就能理解系数和是直流增益。

\textbf{重要边界：系数和不是万能分类器。}和为1不一定是平滑核，也可能是“原图+细节”的最终锐化核；和为0只证明抑制常数，不能仅凭这一条就保证所有高频都通过。还应看正负系数的布局、对图案的实际响应，必要时再研究完整频率响应。

\subsection{均值滤波：每个邻居同样重要}
\textbf{均值滤波（mean filter）}对窗口里所有数赋予相等权重。$3\times3$中有9个数，核为
\[
 B=\frac19\begin{bmatrix}1&1&1\\1&1&1\\1&1&1\end{bmatrix}.
\]
它的系数和是$9/9=1$。假设小区域本来处处是50，观测为
\[
 P=\begin{bmatrix}48&52&49\\51&50&53\\47&51&49\end{bmatrix},
 \qquad
 g=\frac{48+52+49+51+50+53+47+51+49}{9}=50.
\]
这次正负噪声刚好抵消；其他观测未必精确得到50。在局部真实亮度相同、9个噪声零均值且互不相关、方差相同的条件下，输出噪声方差为$\sigma^2/9$，标准差为$\sigma/3$。推理与多帧平均一样，但\textbf{现在假设的是空间邻居之间的噪声关系和真实亮度关系}。

为什么会模糊边缘？如果左边物体灰度0、右边物体灰度90，窗口为
\[
 P=\begin{bmatrix}0&0&90\\0&0&90\\0&0&90\end{bmatrix},
 \qquad g=(0+0+90+0+0+90+0+0+90)/9=30.
\]
中心原本是0，却被右边物体拉到30。\textbf{平均不认识物体边界，只认识窗口内的数。}窗口越大，随机噪声通常更易平滑，但越容易混合相隔更远的区域。

\subsection{高斯滤波：中心近的像素多说一点话}
\textbf{高斯滤波（Gaussian filter）}仍然做乘加，但采用中心大、远处小的权重。常用的小型近似核为
\[
 B_G=\frac1{16}\begin{bmatrix}1&2&1\\2&4&2\\1&2&1\end{bmatrix}.
\]
其权重总和$1+2+1+2+4+2+1+2+1=16$，除以16后和为1。对称布局没有特别偏向左边或右边；中心权重$4/16$比角点$1/16$大。

“高斯”这个名字来自钟形权重曲线：离中心越远，权重逐渐下降。更大的离散核常由
\[
 \widetilde w(s,t)=\exp\!\left(-\frac{s^2+t^2}{2\sigma_G^2}\right),
 \qquad
 w(s,t)=\frac{\widetilde w(s,t)}{\sum_{s,t}\widetilde w(s,t)}
\]
构造。$s^2+t^2$是到中心的距离平方；$\sigma_G>0$控制钟形有多宽；若$s,t$以像素间隔计，它也以像素间隔为单位。\textbf{它是滤波器参数，不是前面噪声的标准差$\sigma$}。第二步归一化保证有限窗口内系数和为1。

对平坦区域中互不相关、同方差的噪声，一般加权平均的噪声方差为
\[
 \operatorname{Var}\!\left(\sum_j w_jn_j\right)=\sigma^2\sum_jw_j^2.
\]
上面的$3\times3$高斯近似核给出
\[
 \sigma^2\frac{1^2+2^2+1^2+2^2+4^2+2^2+1^2+2^2+1^2}{16^2}
 =\frac{36}{256}\sigma^2=\frac9{64}\sigma^2.
\]
同样9个独立样本，等权均值的方差$\sigma^2/9$更小；高斯权重选择的是\textbf{更重视近邻的空间折中}，不能声称它在所有条件下都比均值去噪更强。两者都会模糊边缘；“高斯”也不意味着任何高斯噪声都能无损去除。

\subsection{中值滤波：先排队，再取中间那个数}
\textbf{中值／中位数（median）}不是平均数。把数从小到大排列，取最中间的一个。9个数排序后取第5个；$5\times5$窗口有25个数，就取第13个。

下面是一块真实参与计算的9数窗口，中心遭遇了很亮的脉冲：
\[
 P=\begin{bmatrix}10&12&11\\9&255&10\\11&8&10\end{bmatrix}.
\]
\textbf{均值的算法：}把9个数相加，得到
\[
 10+12+11+9+255+10+11+8+10=336,
 \qquad\text{均值}=336/9\approx37.33.
\]
\textbf{中值的算法：}先把它们抄成一列并排序：
\[
 8,\ 9,\ 10,\ 10,\ \underbrace{10}_{\text{第5个}},\ 11,\ 11,\ 12,\ 255,
 \qquad\text{中值}=10.
\]
输出只把\textbf{中心位置}改成10，并不是把整个$3\times3$都改成10。255很大，却位于排序队伍末端，未能主导中值；均值则被它明显拉高。

若同一块窗口使用高斯近似核，中心脉冲恰好得到最大权重：
\[
 \frac{10+24+11+18+1020+20+11+16+10}{16}
 =\frac{1140}{16}=71.25.
\]
这说明面对孤立极值，简单加权平均不一定是合适选择。\textbf{椒盐噪声（salt-and-pepper noise）}常表现为随机黑点与白点；更一般的\textbf{脉冲噪声（impulse noise）}指少数位置突然出现很大的异常值。中值滤波常适合处理这类异常。

\subsection{中值为什么常保边，却也会删掉细线}
比较两种$3\times3$窗口，灰度只取0与90：
\[
 P_{\text{边缘}}=\begin{bmatrix}0&0&90\\0&0&90\\0&0&90\end{bmatrix},
 \qquad
 P_{\text{细线}}=\begin{bmatrix}0&90&0\\0&90&0\\0&90&0\end{bmatrix}.
\]
第一块中6个0、3个90，中值为0，中心暗像素保持0；均值则为30，产生过渡灰。第二块也是6个0、3个90，\textbf{中值同样为0}，但此时中心原本属于真实亮线，亮线被删除了。算法仅凭窗口数值无法自动知道哪三个90是细线、哪三个90是噪声。

\begin{center}
\begin{tikzpicture}[x=.55cm,y=.55cm,font=\small]
 \foreach \x in {0,...,5} \foreach \y in {0,...,3}{
   \ifnum\x<3 \fill[white] (\x,\y) rectangle ++(1,1);
   \else \fill[gray!55] (\x,\y) rectangle ++(1,1);\fi}
 \fill[black] (1,2) rectangle (2,3);\fill[white] (4,1) rectangle (5,2);
 \draw[step=1] (0,0) grid (6,4);
 \node at (3,-.8) {带脉冲的边界};
 \draw[-{Stealth},thick] (6.7,2)--(8,2);
 \foreach \x in {9,...,14} \foreach \y in {0,...,3}{
   \ifnum\x<12 \fill[white] (\x,\y) rectangle ++(1,1);
   \else \fill[gray!55] (\x,\y) rectangle ++(1,1);\fi}
 \draw[step=1] (9,0) grid (15,4);
 \node at (12,-.8) {去除孤立异常示意};
 \node[align=center,text width=12cm] at (7.5,-2.5) {是否保留结构取决于窗口内的数量与形状\\窄线、小点、尖角都可能被中值改变};
\end{tikzpicture}
\end{center}
所以“中值保边”是相对于某些线性平均方法的常见优势，\textbf{并非保护任意边缘、细线或角点的保证}。噪声比例很大、窗口过大、目标本身太小时，结果都可能变差。

\subsection{最大值与最小值：排序后也可以选两端}
\textbf{顺序统计滤波（order-statistic filtering）}先根据数值大小排列窗口，再选择某个次序的值。中值选择中间；\textbf{最大值滤波（maximum filter）}选择最后一个，\textbf{最小值滤波（minimum filter）}选择第一个。

仍用窗口$[10,12,11,9,255,10,11,8,10]$，排序后最小值为8、中值为10、最大值为255。这三种操作得到不同答案，因为它们在回答不同的问题：窗口中“最暗是多少”“中间是多少”“最亮是多少”。最大值会让亮区域向邻近位置扩展，有机会填掉孤立暗点，却会扩大亮脉冲；最小值反过来，让暗区域扩展，可去掉孤立亮点，却会扩大暗脉冲。因此二者并不是与中值等价的通用去椒盐方法。

在灰度形态学中，采用同样的\textbf{平坦结构元素（flat structuring element，即只指定参与比较的位置、不另加高度权重）}时，局部最大值对应\textbf{膨胀（dilation）}，局部最小值对应\textbf{腐蚀（erosion）}，具体位置还需遵守所采用的反射／锚点约定。本章只需记住数值规则与亮暗区域的变化，不必先学习完整形态学。

\subsection{中值为什么不是固定系数卷积}
线性方法应当满足“先相加再滤波”等于“分别滤波再相加”。用3点中值就能构造反例：
\[
 a=[0,0,1],\quad b=[0,1,0],\quad a+b=[0,1,1].
\]
分别取中值：$\operatorname{median}(a)=0$，$\operatorname{median}(b)=0$，两者之和为0；先相加再取中值却得到1。因此
\[
 \operatorname{median}(a+b)\ne\operatorname{median}(a)+\operatorname{median}(b).
\]
这就是\textbf{非线性（nonlinear）}的具体证据。中值可以用滑动窗口规定“看哪些像素”，但不能用一张对所有输入都固定不变的系数表实现同样操作。

\subsection{先看变化量，再认识导数}
一排像素为$[0,0,0,10,10,10]$，从左往右在第3和第4个像素之间由暗变亮。相邻两数相减，得到
\[
 [0-0,\ 0-0,\ 10-0,\ 10-10,\ 10-10]=[0,0,10,0,0].
\]
这叫\textbf{一阶差分（first difference）}：问的是“下一点比这一点变了多少”。平坦部分给0，跳变位置给10。

再比较相邻的变化量，便得到\textbf{二阶差分（second difference）}。一个常用的中心写法为
\[
 L_j=f_{j-1}-2f_j+f_{j+1}.
\]
注意中心要乘$-2$，不是只把两个邻居相加。它也可以理解为
\[
 L_j=(f_{j+1}-f_j)-(f_j-f_{j-1}),
\]
即“右边的变化量减左边的变化量”。在上面的阶跃两侧，中心取第3个像素时$L=0-2\cdot0+10=10$；中心取第4个像素时$L=0-2\cdot10+10=-10$。于是出现一正一负的成对响应。
\begin{center}
\begin{tikzpicture}[x=.83cm,y=.055cm,font=\small]
 \draw[-{Stealth}] (0,-5)--(6.7,-5) node[right] {位置};
 \draw[thick] (.5,0)--(2.5,0)--(3.5,10)--(5.5,10);
 \foreach \x/\v in {.5/0,1.5/0,2.5/0,3.5/10,4.5/10,5.5/10}{\fill (\x,\v) circle (1.7pt);}
 \node[anchor=east] at (-.1,8) {原像素};
 \draw[gray] (0,-38)--(6,-38);
 \draw[thick] (3,-38)--(3,-28);\fill (3,-28) circle (1.7pt);
 \node[anchor=east] at (-.1,-32) {一阶：$+10$};
 \draw[gray] (0,-78)--(6,-78);
 \draw[thick] (2.5,-78)--(2.5,-68);\fill (2.5,-68) circle (1.7pt);
 \draw[thick] (3.5,-78)--(3.5,-88);\fill (3.5,-88) circle (1.7pt);
 \node[anchor=east] at (-.1,-77) {二阶：$+10,-10$};
 \node[align=left,anchor=west] at (7,-40) {一阶定位变化\\二阶定位变化量的改变\\正负号不能随意丢掉};
\end{tikzpicture}
\end{center}
再看匀速上升的数列$[0,10,20,30,40]$。一阶差分都是10；中间各点的二阶差分都是0，例如$10-2\cdot20+30=0$。\textbf{二阶并非“比一阶更大的同一个数”，而是在问另一件事。}

连续数学中，一阶变化率叫\textbf{导数（derivative）}；图像里沿行、列方向分别计算的一阶变化量组成\textbf{梯度（gradient）}。像素网格上的差分是导数的离散近似。本章先掌握上面的减法，就能读懂后续导数符号。

\subsection{方向要说完整：沿列变化与竖直走向的边缘}
考虑两块窗口：
\[
 P_c=\begin{bmatrix}10&10&90\\10&10&90\\10&10&90\end{bmatrix},
 \qquad
 P_r=\begin{bmatrix}10&10&10\\10&10&10\\90&90&90\end{bmatrix}.
\]
$P_c$从左到右，也就是\textbf{沿列号$c$增加的方向}变化；边界本身却是一条\textbf{竖直走向}的线。$P_r$从上到下，也就是\textbf{沿行号$r$增加的方向}变化；边界本身是\textbf{水平走向}的线。

例如核
\[
 K_c=\begin{bmatrix}-1&0&1\\-1&0&1\\-1&0&1\end{bmatrix}
\]
按相关计算，对$P_c$得到$3(90-10)=240$，对$P_r$得到0。这类核比较右列和左列，因此对\textbf{左右亮度变化}敏感，常用来发现\textbf{竖直走向边界}。只写“水平算子”或“垂直边缘算子”容易因命名习惯不同而歧义，答题时应写清楚比较哪两侧。

\subsection{拉普拉斯：把两个方向的二阶差分加起来}
\textbf{拉普拉斯算子（Laplacian operator）}把行方向、列方向的二阶差分相加。这里默认相邻采样间隔为1；若要估计真实物理坐标中的二阶导数，还应分别除以对应采样间隔的平方。先固定一种符号约定：
\begin{align*}
 Lf(r,c)={}&f(r-1,c)+f(r+1,c)+f(r,c-1)+f(r,c+1)\\
          &-4f(r,c).
\end{align*}
这一写法近似连续的$\nabla^2 f$，对应中心为\textbf{负4}的核：
\[
 L=\begin{bmatrix}0&1&0\\1&-4&1\\0&1&0\end{bmatrix}.
\]
这里只使用上、下、左、右4个邻居，叫\textbf{四邻域（4-neighborhood）}。若定义相反符号$P=-L$，核就变成
\[
 P=\begin{bmatrix}0&-1&0\\-1&4&-1\\0&-1&0\end{bmatrix}.
\]
两种写法都可以，但\textbf{公式、核与最后加减号必须配套}。

设中心70、四邻居都是50：
\[
 Lf=50+50+50+50-4\cdot70=-80,\qquad Pf=80.
\]
中心比周围亮，$Lf$为负；使用$g=f-Lf$时，中心变成$70-(-80)=150$，比原来更亮，局部差异增强。同一个结果也可以写成$g=f+Pf=70+80=150$。
\[
 \boxed{g=f-Lf=f+Pf},\qquad
 \boxed{K_{\rm sharp}=I-L=I+P=
 \begin{bmatrix}0&-1&0\\-1&5&-1\\0&-1&0\end{bmatrix}}.
\]
$I$不是把所有数都设成1的网格，而是\textbf{只复制中心}的单位脉冲核：
\[
 I=\begin{bmatrix}0&0&0\\0&1&0\\0&0&0\end{bmatrix}.
\]
纯拉普拉斯核的系数和为0，常数区消失；最终锐化核的系数和为1，常数区保留。\textbf{“提取细节图”和“输出锐化照片”是两个不同阶段。}

加入对角邻居的常见正中心版本为
\[
 P_8=\begin{bmatrix}-1&-1&-1\\-1&8&-1\\-1&-1&-1\end{bmatrix}.
\]
它对8个方向对称取样，常称八邻域拉普拉斯类核。四邻域和八邻域版本都没有选定单一朝向，但在有限方格上\textbf{不具有连续平面任意角度下的严格旋转不变性}。考试可说“近似无单一方向偏好”，不要把它理解成各方向的数值响应必然完全一致。

\subsection{一阶、二阶与噪声：不是越尖锐越好}
灰度阶跃在一阶差分中常出现单个主要响应，在二阶差分中常出现一正一负的双响应。二阶对细线、孤立点等快速局部变化很敏感，适合构造细节增强；但孤立噪声也是快速变化，因此也会被放大。

“梯度边缘更粗、二阶边缘更细”只是对某些离散算子与显示处理的常见比较，不是所有核、尺度、阈值下都成立的定理。例如上一幅理想阶跃的一阶前向差分只有一个非零样本，二阶反而有两个。答题时应把\textbf{阶跃响应、细节敏感性和噪声敏感性}说清，不要仅背“粗”和“细”。

最终锐化也可能产生\textbf{过冲／欠冲（overshoot/undershoot）}。对一行$[0,0,0,10,10,10]$使用$g=f-Lf$，阶跃附近可出现$-10$和20，超出原范围$[0,10]$。如果直接裁剪，会改变本来推导的线性结果；如果增益很大，会在边界两侧看见不自然的明暗边。实际处理中应控制锐化强度，并考虑先抑制噪声。

\subsection{非锐化掩模：先模糊，再相减，最后加回}
\textbf{非锐化掩模（unsharp masking）}这个名字容易误导：目的是让图变清晰，但需要先制作一张“不锐”的模糊图。设原图为$f$，用归一化平滑核得到$f_{\rm blur}$，则
\[
 m=f-f_{\rm blur}
\]
是\textbf{细节掩模（detail mask）}。平坦区域中两者近似相同，相减接近0；细节处不同，相减留下正负细节。接着把细节按强度$k$加回原图：
\[
 \boxed{g=f+k\,m=f+k(f-f_{\rm blur})},\qquad k\ge0.
\]
$k=0$不改变图像，$k=1$是常见的基本锐化，$k$越大，细节放大通常越明显，噪声与明暗边也可能越明显。
\begin{center}
\begin{tikzpicture}[node distance=9mm and 13mm,font=\small,
 box/.style={draw,rounded corners=1mm,minimum height=9mm,align=center,inner sep=4pt}]
 \node[box] (f) {原图$f$};
 \node[box,right=of f] (blur) {平滑\\$f_{\rm blur}$};
 \node[box,right=of blur] (sub) {相减\\$m=f-f_{\rm blur}$};
 \node[box,below=of sub] (gain) {乘强度$k$};
 \node[box,left=of gain] (add) {加回原图\\$g=f+km$};
 \draw[-{Stealth}] (f)--(blur);
 \draw[-{Stealth}] (blur)--(sub);
 \draw[-{Stealth}] (f.north)--++(0,.7)-|(sub.north);
 \draw[-{Stealth}] (sub)--(gain);
 \draw[-{Stealth}] (gain)--(add);
 \draw[-{Stealth}] (f.south)|-(add.west);
\end{tikzpicture}
\end{center}
数值例子：一行灰度为$[40,40,40,80,80,80]$，用3点均值作模糊。在暗侧紧靠边缘的40处，
\[
 f_{\rm blur}=(40+40+80)/3=53.33\ldots,\quad
 m=40-53.33\ldots=-13.33\ldots;
\]
在亮侧紧靠边缘的80处，
\[
 f_{\rm blur}=(40+80+80)/3=66.66\ldots,\quad
 m=80-66.66\ldots=13.33\ldots.
\]
取$k=1$，两个边缘像素变成约26.67和93.33，明暗差加大；离边缘较远处仍为40或80。这说明锐化改变的是\textbf{局部对比}，不是凭空恢复相机没拍到的真实细节。

\subsection{高提升的两种写法：必须先检查亮度是否归一化}
\textbf{高提升滤波（high-boost filtering）}的记号在教材间并不完全一致。一种常用写法是上述$g=f+k(f-f_{\rm blur})$，较大的$k$表示更强的细节提升。原讲义使用的是另一种形式：
\[
 g_A=A f-f_{\rm blur}=(A-1)f+(f-f_{\rm blur}),\qquad A>1.
\]
它们都涉及原图和模糊图，但参数\textbf{不能不经换算就当作同一个量}。用常数图$f=C$测试，归一化模糊给出$f_{\rm blur}=C$，于是
\[
 f+k(f-f_{\rm blur})=C,\qquad g_A=(A-1)C.
\]
例如均匀区$C=80$：
\[
 A=1.5\Rightarrow g_A=40;\quad
 A=2\Rightarrow g_A=80;\quad
 A=3\Rightarrow g_A=160.
\]
所以\textbf{$Af-f_{\rm blur}$并非对任意$A>1$都保持亮度}；它只有$A=2$时保持常数，除非另做归一化。在$A=1$时，得到的只是$m=f-f_{\rm blur}$，不是已把细节加回原图的最终锐化图。

若将$g_A$除以$A-1$，则
\[
 \frac{g_A}{A-1}
 =f+\frac1{A-1}(f-f_{\rm blur}),\qquad A>1.
\]
它才具有常数增益1，相当于$k=1/(A-1)$。也就是说，\textbf{按这种归一化约定，$A$越大，细节相对原图的提升反而越小}。背公式时要同时说出输出是否归一化，不能只记“参数越大越锐”。

\subsection{A1 原题：四个掩模逐一计算与判读}
\label{original:A1}
\textbf{题目【练习1原题】。}指出下列掩模作用于图像，会产生什么特性。
\[
 \begin{array}{ll@{\qquad}ll}
 \mathrm{a.}&\begin{bmatrix}2&-1&-1\\-1&2&-1\\-1&-1&2\end{bmatrix}&
 \mathrm{b.}&\begin{bmatrix}0&-1&0\\-1&4&-1\\0&-1&0\end{bmatrix}\\[1.1em]
 \mathrm{c.}&\dfrac12\begin{bmatrix}-2&1&-2\\1&6&1\\-2&1&-2\end{bmatrix}&
 \mathrm{d.}&\dfrac19\begin{bmatrix}1&1&1\\1&1&1\\1&1&1\end{bmatrix}
 \end{array}
\]
判读顺序是：\textbf{先算系数和，再看系数布局，再用简单图案验证}。外面的$1/2$或$1/9$必须乘到每个系数上。

\subsubsection{A1(a)：方向选择性的零和核，不能只说“对角线”}
系数和为$3\cdot2+6\cdot(-1)=0$，所以常数背景输出0。更特别的是：\textbf{每一行、每一列的系数和也都为0}。

如果窗口每一行内部都是相同值，例如
\[
 P=\begin{bmatrix}10&10&10\\20&20&20\\40&40&40\end{bmatrix},
\]
输出为$10(2-1-1)+20(-1+2-1)+40(-1-1+2)=0$。同理，若每一列内部都相同，也输出0。因此它会抑制只沿行方向或只沿列方向变化的这种简单图案。

再对比两条斜亮线：
\[
 P_{\backslash}=\begin{bmatrix}10&0&0\\0&10&0\\0&0&10\end{bmatrix},
 \qquad
 P_{/}=\begin{bmatrix}0&0&10\\0&10&0\\10&0&0\end{bmatrix}.
\]
\[
 \begin{array}{ll}
 K_a\text{作用于}P_{\backslash}:&2\cdot10+2\cdot10+2\cdot10=60,\\
 K_a\text{作用于}P_{/}:&-1\cdot10+2\cdot10-1\cdot10=0.
 \end{array}
\]
这说明同一窗口中心处，左上到右下亮线和右上到左下亮线的响应不同。\textbf{这只是给定图案、给定对齐位置的响应，不能推出另一方向的整条线在所有输出位置都为0。}

\textbf{结论：}这是有方向选择性的零和差分／高通类核，突出某些斜向细节，抑制常数及上述单方向变化。它的$180^\circ$对称布局不同于反对称的一阶差分；若用平滑函数的差分近似分析，其低阶主项属于混合二阶变化。零基础答题时，用“对哪些明确排列响应强、对哪些排列为0”比空泛写“增强对角线”更准确。

\subsubsection{A1(b)：四邻域正中心拉普拉斯核}
系数和$4-1-1-1-1=0$。它计算
\[
 g=4f_{\rm 中}-f_{\rm 上}-f_{\rm 下}-f_{\rm 左}-f_{\rm 右}.
\]
常数50输出$4\cdot50-4\cdot50=0$。若中心70、四邻居50，则输出$4\cdot70-4\cdot50=80$。若中心30、四邻居50，则输出$120-200=-80$。正负响应对应中心相对邻居的亮暗差异。

\textbf{结论：}它是$P=-L$，用于提取二阶细节、细线、孤立点及边界变化，均匀区归零。没有单一选定朝向，但离散方向响应并非严格各向相同。若要得到保留背景的锐化图，应再加原图$g=f+Pf$，不能把这个零和输出直接叫作保亮度的最终照片。

\subsubsection{A1(c)：和为1的细节强调核，不是均值平滑}
先把所有整数相加：4个角贡献$-8$，上下左右贡献4，中心贡献6，总和2。再乘外部$1/2$：
\[
 S=\frac12(-8+4+6)=1.
\]
所以常数50仍输出50。中心实际系数为3，四角为$-1$，轴向邻居为$1/2$，具有明显的减法与中心强调；不能因为和为1就判为低通。

对中心70、其余8个邻居50的窗口：
\[
 g=\frac12\left[6\cdot70+(50+50+50+50)
                     -2(50+50+50+50)\right]
  =\frac{420+200-400}{2}=110.
\]
中心从70变成110，平坦背景50保持50。也可以把它分成$I+(K_c-I)$，其中后一项系数和为0，表示“原图+某种细节”。

\textbf{结论：}这是保留常数背景的锐化／细节强调类核，不是纯高通，也不是等权平均。其轴向邻居为正、对角邻居为负，具体频率与方向响应由完整布局决定；不能说“每一种高频都同样增强”，也不能不加验证地说它就是某个非负模糊核的标准高提升形式。

\subsubsection{A1(d)：$3\times3$均值低通核}
系数和$9\cdot(1/9)=1$，所有系数相等且非负。每个输出都是邻域9个数的平均。常数50仍为50；前文8个近10的数加一个255的窗口，则输出$336/9\approx37.33$。

\textbf{结论：}这是均值平滑／低通滤波器，可减弱局部随机起伏，也会模糊边缘和小结构。窗口跨过0到90的边界时会产生30、60等中间灰度；对强脉冲噪声不如合适的中值处理稳健。

\subsection{A2 原题：三种核的性质与方向}
\label{original:A2}
\textbf{题目。}判断下列三个掩模各属于低通还是高通、平滑还是锐化、是否有方向性。
\[
 \text{① }\begin{bmatrix}0&-1&0\\-1&5&-1\\0&-1&0\end{bmatrix}
 \qquad
 \text{② }\begin{bmatrix}-1&-1&-1\\-1&8&-1\\-1&-1&-1\end{bmatrix}
 \qquad
 \text{③ }\begin{bmatrix}-1&-1&-1\\2&2&2\\-1&-1&-1\end{bmatrix}.
\]

\subsubsection{A2①：原图加四邻域细节，最终锐化核}
$S=5-4=1$，故常数50输出50。它恰好等于$I+P$，所以是\textbf{保留常数背景的高频强调／锐化核}，不是去掉背景的纯高通输出。

中心70、上下左右50时：
\[
 g=5\cdot70-50-50-50-50=350-200=150.
\]
没有选定唯一方向，使用四邻域的近似方向对称结构。但如果题目要求二选一“低通或高通”，应说明它属于\textbf{锐化的高频强调类，同时保留直流}，而不是因为系数和1就选低通。

\subsubsection{A2②：八邻域细节提取核}
$S=8-8=0$，常数区输出0。中心70、其余8个邻居50时，
\[
 g=8\cdot70-8\cdot50=560-400=160.
\]
中心越偏离周围，绝对响应越大。\textbf{结论：}八邻域拉普拉斯类纯高通／二阶细节提取核，无单一选定朝向，包含对角比较；它本身不保留背景。需要最终锐化时还应规定如何与原图相加及使用多大系数。

\subsubsection{A2③：沿行方向二阶差分，同时沿列方向求和}
该核可写成竖直系数列与水平系数行的外积：
\[
 \begin{bmatrix}-1\\2\\-1\end{bmatrix}
 \begin{bmatrix}1&1&1\end{bmatrix}
 =\begin{bmatrix}-1&-1&-1\\2&2&2\\-1&-1&-1\end{bmatrix}.
\]
即先把上、中、下三行各自的3个像素相加，分别记为$R_{\rm 上},R_{\rm 中},R_{\rm 下}$，再算
\[
 g=-R_{\rm 上}+2R_{\rm 中}-R_{\rm 下}.
\]
沿列方向的3点求和是一种\textbf{未除以3的平滑式汇集}；跨行的$[-1,2,-1]$则是负二阶差分。

\textbf{试验1：匀速渐变。}对三行分别全是10、20、30的窗口，
\[
 g=-3\cdot10+2(3\cdot20)-3\cdot30=-30+120-90=0.
\]
一阶变化存在，二阶变化却为0，已经说明它不是一般一阶Prewitt核。

\textbf{试验2：上下阶跃。}三行分别全是10、10、40时，
\[
 g=-30+60-120=-90;
\]
窗口下移，使三行变成10、40、40时，
\[
 g=-30+240-120=90.
\]
跨过同一个水平走向边界，产生一负一正的成对响应，符合二阶差分特征。

\textbf{试验3：左右阶跃。}若每一行都是$[10,10,40]$，则三行之和都为60：
\[
 g=-60+2\cdot60-60=0.
\]
\textbf{结论：}它是有方向性的零和二阶差分／高通类核，对\textbf{从上到下的灰度变化}敏感，常响应\textbf{水平走向边界}；对上述仅左右变化的理想图案为0。原讲义把它叫作“一阶差分／类Prewitt”不正确，应以系数$[-1,2,-1]$和数值试验为准。

\subsection{D1 原题与完整解答：中值滤波的五个问题}
\label{original:D1}
\textbf{题目【练习1原题】。}对中值滤波回答：①如何操作？②作用是什么？③是否可以用模板操作？④影响算法效率的因素有哪些？⑤试估计其频域特性。

\textbf{①如何操作。}选定窗口、锚点和边界规则；取出当前窗口数值；从小到大排序；奇数个样本取中间那个值，写到输出图对应位置；再移动窗口。9数窗口$[10,12,11,9,255,10,11,8,10]$排序后为$[8,9,10,10,10,11,11,12,255]$，第5个为10。符号写法为
\[
 g(r,c)=\operatorname{median}\{f(r+s,c+t):(s,t)\in W\},
\]
其中$W$就是窗口允许的偏移集合。若用偶数个样本，必须另行规定中间两个值如何处理，因此课堂通常用奇数尺寸。

\textbf{②作用。}非线性平滑，常用于抑制孤立脉冲和椒盐噪声。在合适条件下比同尺寸均值更少混合阶跃两侧的亮度；但会改变细线、尖角和小目标，也可能在高噪声比例时失效。

\textbf{③能否用模板。}能用\emph{窗口模板}选邻居；不能用\emph{固定线性系数模板}以卷积乘加实现所有输入的中值。可用$[0,0,1]$与$[0,1,0]$的加法反例证明其非线性，而不只背一句“因为要排序”。

\textbf{④效率因素。}设图像有$M\times N$个输出位置，窗口含$K$个数。朴素方法每个窗口完整排序，大致需要$O(K\log K)$次量级的比较，总量约$O(MNK\log K)$。这里$O(\cdot)$只表示随规模增长的数量级，不是精确运行秒数。窗口大小、图像尺寸、排序或选择算法、滑动时能否复用上一个窗口的信息、灰度是否为有限整数、内存访问与并行硬件都会影响速度。例如$3\times3$窗口需处理9个数，$7\times7$需处理49个数，不能认为只多做一两次运算。利用直方图更新等方法可明显改变复杂度。

\textbf{⑤频域特性。}中值是非线性的，不能给它一个对所有输入通用的固定传递函数$H$，再一律写“输出频谱=输入频谱乘$H$”。可以定性说它常抑制孤立的快速起伏，呈现某些平滑／低通效果；但其作用依赖信号形状、幅度、窗口和噪声。用正弦条纹测试时，非线性还可能产生新频率，所得只能是\textbf{该测试条件下的等效响应}，不是唯一线性频率响应。

\subsection{D2 原题与完整解答：均值和中值如何选择}
\label{original:D2}
\textbf{题目。}比较均值滤波与中值滤波在噪声抑制上的差异，分别适用哪种噪声？

\textbf{答题的出发点是计算规则。}均值把所有样本幅值都加入总和；一个很大的错误值会拉动平均。中值看排序位置，少数极端值通常落在两端，未必改变中间值。因此前面的含255窗口，均值为37.33，中值为10，差异不是算法“更聪明”，而是对异常值的敏感性不同。
\begin{center}
\begin{tabularx}{\linewidth}{@{}p{.16\linewidth}XX@{}}
\toprule
比较项目 & 均值滤波 & 中值滤波\\
\midrule
计算 & 求和后除以样本数；线性 & 排序取中间值；非线性\\
典型适用 & 局部平坦区域中的零均值、近似独立加性噪声，如高斯噪声 & 少数位置的椒盐／脉冲噪声\\
边缘代价 & 混合边缘两侧数值，形成过渡灰 & 常较少产生过渡灰，但会改动形状、删除小目标\\
极值影响 & 极值会拉偏均值，并随窗口移动扩散影响 & 少量极值通常不成为中值；比例过高仍会失败\\
\bottomrule
\end{tabularx}
\end{center}
\textbf{结论：}考试的基本匹配是“近似独立加性高斯噪声可考虑均值／高斯平滑，椒盐噪声优先考虑中值”。这不是不看图像结构的绝对排名。若要保留非常细的线，二者都可能有明显损失；若噪声强相关，独立噪声的方差下降公式也不能直接套用。

\subsection{D3 原题与完整解答：为何系数和是1或0}
\label{original:D3}
\textbf{题目。}为什么平滑（低通）模板的系数和等于1，而纯锐化（高通）模板的系数和等于0？

\textbf{先修正题目用语。}这里“系数和0”指\textbf{纯细节提取／零直流高通核}，不是所有最终锐化核。很多最终锐化核恰好和为1。

设常数输入为$C=100$。核系数为$w_1,\ldots,w_K$，输出为
\[
 g=\sum_jw_j\cdot100=100\sum_jw_j.
\]
希望平坦亮度仍为100，就令$\sum_jw_j=1$；希望完全没有变化的区域不产生细节，就令$\sum_jw_j=0$。举例：
\[
 \begin{array}{ll}
 \text{均值核：}&100\times(9/9)=100,\\
 \text{拉普拉斯细节核：}&100\times(4-4)=0,\\
 \text{原图+拉普拉斯：}&100\times(5-4)=100.
 \end{array}
\]
\textbf{补充条件：}图像边界采用保持常数的延拓规则，且未发生裁剪；“和为1”本身只保证常数响应，不保证整个核一定低通。对有限图像，要讨论全图平均值是否严格保持，还必须考虑边界实现，不能把局部常数测试扩大成任意边界下的全局平均定理。

\subsection{D4 原题与完整解答：拉普拉斯公式及一、二阶差异}
\label{original:D4}
\textbf{题目。}写出拉普拉斯锐化的公式与典型掩模，并说明一阶导数与二阶导数在锐化中的区别。

\textbf{第一部分：符号配套写全。}
\[
 \nabla^2f\ \approx\ Lf
 =f_{\rm 上}+f_{\rm 下}+f_{\rm 左}+f_{\rm 右}-4f_{\rm 中},
 \quad L=\begin{bmatrix}0&1&0\\1&-4&1\\0&1&0\end{bmatrix},
 \quad g=f-Lf.
\]
若使用中心为正4的$P=-L$，则写成
\[
 P=\begin{bmatrix}0&-1&0\\-1&4&-1\\0&-1&0\end{bmatrix},
 \qquad g=f+Pf.
\]
两套符号给出同一个最终核$[0,-1,0;-1,5,-1;0,-1,0]$。原讲义把“邻居和减4倍中心”的公式与“中心正4”的核并排对应，符号不一致；不能照抄后再任意选加或减。

\textbf{数值验算。}中心70、邻居50：$Lf=200-280=-80$，于是$g=70-(-80)=150$；或$Pf=80$，$g=70+80=150$。假如错误地对正中心核采用$70-80=-10$，局部亮峰被反向压低，说明符号配错。

\textbf{第二部分：一阶与二阶。}一阶问相邻像素变化多少；二阶问变化量怎样变化。前述0到10阶跃的一阶有一个$+10$主要响应，二阶在阶跃两侧出现$+10,-10$；匀速斜坡一阶非零、二阶为0。二阶对细线、孤立点及噪声敏感，常用于细节增强；梯度常用于检测轮廓与方向。实际边缘宽度还取决于核、尺度和阈值，不能把“粗／细”作为绝对定义。

\subsection{D5 原题与完整解答：非锐化掩模和高提升的关系}
\label{original:D5}
\textbf{题目。}说明非锐化掩模与高提升滤波的关系。

\textbf{先定义中间结果：}$m=f-f_{\rm blur}$是原图减模糊图得到的细节掩模。它在常数区为0，单独显示时不再保留原图背景；显示正负细节还需要选择灰度映射。

\textbf{再定义最终输出：}常见非锐化掩模锐化为$g=f+km$，$k=1$时$g=2f-f_{\rm blur}$。增加$k$可加强细节，但也会放大噪声和明暗边。

\textbf{与原题高提升写法对照：}
\[
 g_A=Af-f_{\rm blur}=(A-1)f+m.
\]
$A=1$时是纯掩模$m$；$A=2$时等于$k=1$的保常数锐化。一般$A>1$的常数增益为$A-1$，不能无条件说“保持原图亮度”。若要归一化，应写
\[
 g_{\rm norm}=\frac{Af-f_{\rm blur}}{A-1}
 =f+\frac{m}{A-1}.
\]
\textbf{数值检查：}原图80、模糊图80，在任何$k$下$f+km=80$；而$A=3$时$Af-f_{\rm blur}=240-80=160$。原图100、模糊图90时，掩模10；$k=1$输出110，$A=2$也输出110；$A=3$输出210，若除以2归一化则为105。先检查参数约定，再讨论锐化程度。

\subsection{把计算结果变成能交卷的解释}
遇到陌生核，可以按下面的次序写答案，而不必靠外形猜名称：
\begin{enumerate}[label=\arabic*.]
 \item \textbf{计算系数和。}说明常数输入$C$会变成什么；这一步区分是否保留背景。
 \item \textbf{检查正负和位置。}全非负且归一化的局部平均具有平滑意义；有正负项时，说明是在比较哪些位置。
 \item \textbf{选择一个小图案手算。}常数、孤立亮点、上下阶跃、左右阶跃或斜线都可以。一次只比较一个因素。
 \item \textbf{说清输出是哪一层。}是噪声较少的图、纯细节图，还是已经把原图加回的锐化图？
 \item \textbf{补上限制。}边界规则、取整裁剪、噪声放大、细线损失，以及方向性不是绝对各向同性等。
\end{enumerate}
本章最值得记住的不是核的外号，而是三个可检验的问题：\textbf{每个邻居怎样参与计算？一块常数会输出什么？一个明确的变化图案会输出什么？}能回答这三个问题，就能把大多数滤波判断题变成清楚的算术。

\section{模块五：频域图像增强，从一行像素开始}

\noindent\textbf{这一模块要解决什么？}\quad 一幅图里可能同时有大范围的明暗变化、物体边界、草地纹理和重复条纹。我们希望把其中某些变化保留下来，把另一些变化减弱。\textbf{频域（frequency domain）}把问题改写成：“变化有多快？沿哪个方向变化？这一类变化有多强？”先理解这三个问题，再学公式，滤波器就不再只是需要背诵的字母。

本模块沿用原讲义模块五及原题 B4、D6、D7、E6 的编号。条纹小例子、四点变换、边界补零示例等是为理解原题补充的必要背景；它们不是原讲义中额外承诺的考试题目。原稿中不严谨的结论会在相应位置明确订正。

\subsection{先看一行亮度：什么叫变化快，什么叫变化慢}

把灰度图像的一行像素依次抄下来，就得到一串数。数值大表示亮，数值小表示暗。例如下面每行都有八个像素：
\[
\begin{array}{c|rrrrrrrr}
\text{位置 }x&0&1&2&3&4&5&6&7\\\hline
\text{几乎不变}&100&100&100&100&100&100&100&100\\
\text{缓慢变亮}&40&50&60&70&80&90&100&110\\
\text{反复明暗}&40&110&40&110&40&110&40&110
\end{array}
\]
第一行没有随位置发生的明暗变化；第二行要走过较长距离才有明显变化；第三行只走一个像素就从暗变亮或从亮变暗。频域中的“高”和“低”，首先描述这种\textbf{空间变化的快慢}，不描述像素灰度数值的高低。一整块灰度为 240 的白墙仍可以主要是低频；在 20 和 30 之间密集交替的暗条纹也可以含很高的空间频率。

\textbf{空间频率（spatial frequency）}是每单位距离内完成的周期数。一个完整周期可以是“亮到暗，再回到亮”。若相邻两条亮纹之间相隔 $T$ 个像素，则
\[
\nu=\frac{1}{T}.
\]
这里 $T$ 是\textbf{周期（period）}，单位为像素；$\nu$ 是空间频率，单位为周期/像素。这个公式用于把可见的条纹间距转换成变化快慢：$T=8$ 时 $\nu=1/8=0.125$ 周期/像素；$T=4$ 时 $\nu=1/4=0.25$ 周期/像素。后一组条纹更密，所以频率更高。若距离单位换成毫米，频率单位就相应成为周期/毫米。

\begin{center}
\begin{tikzpicture}[x=0.68cm,y=0.46cm,>=Latex,font=\small]
\draw[->] (0,0)--(16.8,0) node[right]{$x$};
\draw[->] (0,-.2)--(0,3.5) node[above]{灰度};
\draw[dashed,gray] (0,1.6)--(16,1.6);
\draw[thick,domain=0:16,samples=150] plot(\x,{1.6+1.1*cos(45*\x)});
\draw[<->] (0,3)--(8,3) node[midway,fill=white]{一个周期：8 像素};
\node[anchor=west] (a) at (0,-1){慢变化：16 像素中出现 2 个周期};
\begin{scope}[yshift=-4.4cm]
\draw[->] (0,0)--(16.8,0) node[right]{$x$};
\draw[->] (0,-.2)--(0,3.5) node[above]{灰度};
\draw[dashed,gray] (0,1.6)--(16,1.6);
\draw[thick,domain=0:16,samples=180] plot(\x,{1.6+1.1*cos(90*\x)});
\draw[<->] (0,3)--(4,3) node[midway,fill=white]{4 像素};
\node[anchor=west] at (0,-1){快变化：16 像素中出现 4 个周期};
\end{scope}
\end{tikzpicture}
\end{center}

\textbf{读图检查：}两条曲线的平均高度和上下摆动幅度相同，改变的是每段距离中摆动的次数。这说明“频率高”既不等于“平均更亮”，也不等于“对比度更大”。有限图像的边缘和裁切也会影响频谱；“缓慢渐变主要对应低频”是对主体变化的描述，不是保证其他频率完全为零。

\subsection{用波拼出亮度：幅度、相位和直流}

\subsubsection{一条最简单的波要用哪些量描述}
先补一个几何背景：在半径为 1 的圆上，从最右侧出发逆时针转过角度 $\theta$，到达点的横坐标叫 $\cos\theta$（余弦），纵坐标叫 $\sin\theta$（正弦）。转一整圈回到起点，所以这两个数会周期性重复，适合描述明暗反复出现的变化。\textbf{弧度（radian）}是角度的另一种单位：一整圈 $360^\circ=2\pi$ 弧度，四分之一圈 $90^\circ=\pi/2$ 弧度。公式里的角度默认按弧度读。

\begin{center}
\begin{tabular}{@{}lllll@{}}
\toprule
转过的角度&弧度&到达位置&$\cos\theta$&$\sin\theta$\\
\midrule
$0^\circ$&$0$&最右侧&$1$&$0$\\
$90^\circ$&$\pi/2$&最上方&$0$&$1$\\
$180^\circ$&$\pi$&最左侧&$-1$&$0$\\
$270^\circ$&$3\pi/2$&最下方&$0$&$-1$\\
$360^\circ$&$2\pi$&回到最右侧&$1$&$0$\\
\bottomrule
\end{tabular}
\end{center}

只看横坐标，就按 $1,0,-1,0,1$ 变化。把这个上下起伏的数乘上一个幅度，再加上基准亮度，就得到可见的明暗波。一个余弦亮度变化可以写为
\[
f(x)=b+A\cos\!\left(2\pi\nu x+\phi\right).
\]
这里 $f(x)$ 是位置 $x$ 的灰度，$b$ 是基准亮度，$A\geq0$ 是\textbf{振幅（amplitude）}，控制上下摆动的大小；$\nu$ 是频率；$\phi$ 是\textbf{相位（phase）}，控制波峰从哪里开始；$\pi$ 是圆周率。$2\pi$ 表示一个完整周期的角度，因而 $\nu x$ 先给出走过了多少周期，再乘 $2\pi$ 转成余弦函数使用的角度。

若把这条波直接用作不截断的 8 位非负灰度，需要 $b-A\geq0$ 且 $b+A\leq255$，保证最暗和最亮处都在可显示范围内。下面例子满足这一条件。

例如取 $b=100$、$A=20$、$\nu=1/4$、$\phi=0$，得到
\[
f(x)=100+20\cos(\pi x/2).
\]
逐点代入 $x=0,1,2,3$：余弦值分别为 $1,0,-1,0$，灰度就是 $120,100,80,100$。到 $x=4$ 又变成 120，因此每四个像素重复一次。

如果只把 $A$ 改为 40，结果为 $140,100,60,100$，明暗差更大，周期仍是四个像素。如果只把 $\phi$ 改为 $\pi/2$，结果为 $100,80,100,120$，只是波峰移动了位置。\textbf{频率决定重复速度，幅度决定这一成分有多强，相位决定它如何对齐。}

\subsubsection{复杂图像为什么可以谈“有哪些频率”}
一条亮度曲线通常不是单独一条波，但可以由多条不同频率、幅度和相位的波相加得到。\textbf{傅里叶分析（Fourier analysis）}就是找出这些波的“配方”。在有限离散数据上，离散傅里叶变换提供一套可逆的配方：不丢掉系数时，能够用逆变换重建原数据。

例如
\[
f(x)=100+20\cos(2\pi x/8)+5\cos(2\pi x/2).
\]
第一项 100 是处处一样的基准；第二项每八个像素一个周期，振幅 20；第三项每两个像素一个周期，振幅 5。以 $x=0$ 为例，$f(0)=100+20+5=125$；以 $x=2$ 为例，$f(2)=100+20\cos(\pi/2)+5\cos(2\pi)=105$。这里高频成分更弱，但确实存在。滤去第三项，就去掉了细密起伏，而保留较缓慢的明暗变化。

“把图像分解成波”不是说真实场景由许多条纹实体组成，而是换一种数学方式描述同一组像素。类似一个位置可以用经纬度表示，也可以换一套坐标表示；改变表示方式不等于改变图像。

\subsubsection{直流为什么与平均亮度有关}
\textbf{直流分量（direct-current component，DC）}指零频率成分。零频率的波不再随位置起伏，所以对应一个常数。采用后文的正变换不归一化、逆变换除以像素数的约定时，
\[
F(0,0)=\sum_{x=0}^{M-1}\sum_{y=0}^{N-1} f(x,y),
\qquad
\bar f=\frac{F(0,0)}{MN}.
\]
$M,N$ 分别是图像在两个坐标方向的像素数，$MN$ 是总像素数，$\bar f$ 是平均灰度。因此\textbf{直流系数等于总和，除以总像素数才是平均值}。例如 $2\times2$ 图像的四个像素全为 10，总和为 40，故 $F(0,0)=40$，平均值为 $40/4=10$。原稿把直流直接说成“平均灰度”，是在说它代表的含义；数值上须看变换的归一化约定。

\subsection{认识二维频谱：一个点到底代表什么}

\subsubsection{先约定坐标，避免把横竖记反}
本模块统一把 $x$ 看作竖直方向的行坐标，把 $y$ 看作水平方向的列坐标；$u$ 与 $x$ 配对，$v$ 与 $y$ 配对。图像尺寸为 $M$ 行、$N$ 列。一组二维条纹可写成
\[
f(x,y)=b+A\cos\!\left[2\pi\left(\frac{k_x x}{M}+\frac{k_y y}{N}\right)+\phi\right].
\]
$k_x$ 表示沿整幅图的竖直跨度走过多少周期，$k_y$ 表示沿水平跨度走过多少周期。$k_x/M$ 和 $k_y/N$ 是两个方向的空间频率。若 $k_x=0,k_y=4$，则只沿水平方向变化；沿竖直方向保持不变，所以看到的是\textbf{竖条纹}。对应频率点出现在频谱中心的水平方向上。

\begin{center}
\begin{tikzpicture}[font=\small,>=Latex]
\begin{scope}
\fill[gray!15] (0,0) rectangle(4,3);
\foreach \a in {0,...,79}{\pgfmathsetmacro{\shadeval}{50+40*cos(18*\a)}\fill[black!\shadeval] ({0.05*\a},0) rectangle({0.05*\a+0.051},3);}
\draw (0,0) rectangle(4,3);
\draw[->] (0,-.35)--(4,-.35) node[right]{$y$};
\draw[->] (-.35,3)--(-.35,0) node[below]{$x$};
\node[align=center] at (2,-1){竖条纹：向右移动时明暗交替\\向下移动时不变};
\end{scope}
\draw[->,thick] (4.6,1.5)--(5.8,1.5) node[midway,above]{变换};
\begin{scope}[xshift=7.8cm,yshift=1.5cm]
\draw (-2,-1.5) rectangle(2,1.5);
\draw[->] (-1.8,0)--(2.35,0) node[right]{$\Delta v$};
\draw[->] (0,-1.3)--(0,1.95) node[above]{$\Delta u$};
\fill (0,0) circle(2.5pt);
\fill (-1.2,0) circle(3pt);\fill (1.2,0) circle(3pt);
\node[above] at (-1.2,0){$-k_y$};\node[above] at (1.2,0){$+k_y$};
\node[below right] at (0,0){DC};
\node[align=center] at (0,-2.5){频率点沿水平方向分开\\两点连线垂直于条纹走向};
\end{scope}
\end{tikzpicture}
\end{center}

本图为方便读坐标把频谱竖轴正向画向上；软件把矩阵行号画向下时，竖轴方向可能相反，但“哪个方向发生变化”“连线垂直于条纹”不变。横条纹沿竖直方向发生变化，亮点沿频谱竖直方向分开。斜条纹也遵循同一个规则：沿条纹走，亮度变化少；垂直穿过条纹，亮暗交替最多。

\subsubsection{频谱位置、亮度和相位不能混为一谈}
对每个频率位置 $(u,v)$，变换得到一个通常为复数的系数 $F(u,v)$。复数可以写成
\[
F(u,v)=a(u,v)+j\,b(u,v),\qquad j^2=-1.
\]
$a$ 和 $b$ 是这个系数的实部、虚部；这里的 $b(u,v)$ 与前面例子中的基准常数 $b$ 不同。二者是记录波的强弱与对齐位置的一种方便写法。我们最常画出的量是
\[
|F(u,v)|=\sqrt{a(u,v)^2+b(u,v)^2},
\qquad
\theta(u,v)=\operatorname{atan2}(b(u,v),a(u,v)).
\]
$|F|$ 是\textbf{幅度谱（magnitude spectrum）}；$\theta$ 是\textbf{相位谱（phase spectrum）}；$\operatorname{atan2}$ 根据两部分的符号确定正确角度。例如 $F=3+4j$ 时幅度为 $\sqrt{9+16}=5$，相位约为 $53.13^\circ$。若 $F$ 改为 $-3-4j$，幅度仍为 5，但相位不同，重建出的波峰位置也不同。

因此，频谱图中“某处很亮”的意思是\textbf{该频率的成分很强}，不是原图同一位置有个很亮的像素。空间位置与频谱位置是两套坐标；一个频率系数通常影响整幅图像，不能把右上角的频谱点直接解释成原图右上角的物体。

幅度的数值范围常很大，直接显示时强直流会掩盖弱分量。常用显示式为
\[
s(u,v)=c_{\mathrm{disp}}\ln\bigl(1+|F(u,v)|\bigr).
\]
$s$ 是显示用的亮度，$c_{\mathrm{disp}}>0$ 是把结果放进屏幕灰度范围的缩放系数；加 1 避免对零取对数。比如幅度 0 和 99 分别变为 $\ln 1=0$ 和 $\ln100\approx4.605$，大幅度被压缩。\textbf{这只是为了看清频谱，真正滤波仍使用原来的复系数 $F$，不能拿显示图 $s$ 代替它。}

\subsubsection{为什么实图像会出现成对频率点}
一幅普通灰度图的像素是实数。它的频谱具有\textbf{共轭对称性（conjugate symmetry）}：
\[
F((-u)\bmod M,(-v)\bmod N)=\overline{F(u,v)}.
\]
$\overline{F}$ 表示把复数虚部的符号反过来；“$\bmod$”表示把索引绕回合法范围。例如 $3+4j$ 的共轭为 $3-4j$，两者幅度相同。所以中心化后的幅度谱通常关于中心成对对称。

不必先掌握复数技巧也能理解原因：一条实余弦波可由正、负两个相反频率的复波配对表示。两项的虚部互相抵消，留下实数。对非零且非特殊边界频率，单一理想余弦会产生一对亮点。DC 和某些奈奎斯特频率点可与自己配对，因此“所有亮点都另有一个不同点”并非普遍成立。

图像平移主要改变各个波的相位。在\textbf{循环平移、没有新裁切}的离散模型下，幅度谱不变；现实中移动后又裁掉一块内容，幅度也可能改变。只保存幅度而丢弃相位，一般不能唯一恢复原图。

图像整体旋转时，频谱也随之旋转相同角度：原来水平的变化方向旋转后，频率点连线也转向新方向。这是连续理想模型的性质；离散图像做插值旋转或重新裁切时，幅度可能还受到插值与边界的影响。

\subsection{DFT 与 IDFT：把“配方”写成可计算的公式}

\subsubsection{正变换是在问：每一类波用了多少}
还需要认识复指数这个简写。\textbf{欧拉公式（Euler's formula）}为
\[
e^{j\theta}=\cos\theta+j\sin\theta,\qquad
e^{-j\theta}=\cos\theta-j\sin\theta.
\]
$j$ 是满足 $j^2=-1$ 的虚数单位；$\theta$ 是弧度角。这个写法把一对余弦、正弦合在一个数里，方便同时记录波的强弱和对齐位置，并不意味着像素本身必须是复数。利用前面的四分之一圈表，$\theta=0,\pi/2,\pi,3\pi/2$ 时，$e^{-j\theta}$ 依次等于 $1,-j,-1,j$。例如 $\theta=\pi/2$ 时，$\cos\theta=0,\sin\theta=1$，所以 $e^{-j\pi/2}=0-j\times1=-j$。

\textbf{离散傅里叶变换（discrete Fourier transform，DFT）}把像素数组变成频率系数数组。本模块采用
\[
F(u,v)=\sum_{x=0}^{M-1}\sum_{y=0}^{N-1}
 f(x,y)\exp\!\left[-j2\pi\left(\frac{ux}{M}+\frac{vy}{N}\right)\right].
\]
逐一读符号：$f(x,y)$ 是原图的像素；$x=0,\ldots,M-1$ 是行号，$y=0,\ldots,N-1$ 是列号；$u=0,\ldots,M-1$ 和 $v=0,\ldots,N-1$ 是两方向的频率索引。$\sum$ 表示把所有像素的贡献加起来；$\exp(t)=e^t$ 是指数函数；$j$ 是虚数单位。指数中的 $ux/M+vy/N$ 负责指定一组波，负号约定它是正变换。

这条公式并不要求在做滤波题时逐像素手算。它告诉我们三件事：一个频率系数由整幅图决定；系数含幅度与相位；当 $u=v=0$ 时指数为 0，$e^0=1$，因而该系数就是像素总和。

\textbf{快速傅里叶变换（fast Fourier transform，FFT）}是快速计算 DFT 的算法，不是一种不同的变换。类似用更快的乘法算法算出同一个乘积，FFT 与直接求和在数学上得到同样的 DFT，浮点舍入误差除外。

\subsubsection{逆变换是在按配方把所有波加回去}
\textbf{逆离散傅里叶变换（inverse DFT，IDFT）}为
\[
f(x,y)=\frac{1}{MN}\sum_{u=0}^{M-1}\sum_{v=0}^{N-1}
F(u,v)\exp\!\left[j2\pi\left(\frac{ux}{M}+\frac{vy}{N}\right)\right].
\]
变量含义不变，但现在遍历所有频率，把各个成分加到位置 $(x,y)$。与正变换相比，指数符号变成正号，并除以总像素数 $MN$。正、逆两式必须成套使用；别的教材可能把归一化系数放在正变换或拆到两边，只要前后一致，结论相同。

用前面的 $2\times2$ 常数图验证：只有 $F(0,0)=40$ 非零，任意位置的逆变换为 $(1/4)\times40\times1=10$，正好重建原灰度。这就是“直流代表整体亮度”的一次实际运算。

\subsubsection{四个像素的完整小例子}
为看清正、逆变换在做什么，先退回一维，取四个像素
\[
f=[12,10,8,10].
\]
其一维 DFT 是 $F(k)=\sum_{x=0}^{3}f(x)e^{-j2\pi kx/4}$，其中 $k=0,1,2,3$ 是四个频率索引。$k=0$ 时四个乘数都是 1，所以 $F(0)=12+10+8+10=40$。

当 $k=1$ 时，依次使用乘数 $1,-j,-1,j$，于是
\[
F(1)=12-10j-8+10j=4.
\]
$k=2$ 时乘数是 $1,-1,1,-1$，故 $F(2)=12-10+8-10=0$。$k=3$ 时与 $k=1$ 共轭，得到 $F(3)=4$。所以
\[
F=[40,4,0,4].
\]
这个配方对应 $f(x)=10+2\cos(2\pi x/4)$：常数为 10，一条振幅为 2 的余弦；它在频谱中表现为成对的两个系数 4。\textbf{DFT 系数的幅度并不直接等于原余弦的振幅}，还受样本数和正负频率配对影响；本例每个系数为 $4\times2/2=4$。

逆变换在 $x=0$ 处为 $(40+4+0+4)/4=12$；在 $x=2$ 处为 $(40-4+0-4)/4=8$。这样就能看到“频率配方”怎样重新变成像素，而不是把 DFT 只当作一个神秘缩写。

\subsection{滤波的核心：给每种频率乘一个权重}

\textbf{频率响应或传递函数（frequency response/transfer function）}$H(u,v)$ 给每一个频率位置指定一个权重。频域滤波的核心仅是
\[
G(u,v)=H(u,v)F(u,v).
\]
$F$ 是输入频谱，$H$ 是滤波器，$G$ 是输出频谱。若某位置 $H=1$，该分量不变；$H=0$，该分量完全去掉；$H=0.2$，只剩原幅度的五分之一；$H=2$，幅度加倍。对本模块常用的实数、非负 $H$，相位保持不变；一般复数 $H$ 也可以改变相位。

例如某频率系数为 $F=3+4j$，原幅度为 5。若 $H=0.4$，则 $G=1.2+1.6j$，幅度为 $\sqrt{1.2^2+1.6^2}=2$，正好是原来的 0.4 倍。

在四像素例子中，取 $H=[1,0.5,0,0.5]$，得到
\[
G=[40,2,0,2],\qquad g=[11,10,9,10].
\]
平均亮度仍为 10，原来上下摆动 2，现在上下摆动 1。若改成 $H=[1,0,0,0]$，则输出处处为 10，只剩平均亮度。若反过来只去掉直流，$H=[0,1,1,1]$，输出为 $[2,0,-2,0]$。出现负值是正常的：这是亮度\textbf{变化量}，不能直接当成已经放进显示范围的图像。

\subsubsection{与空间卷积的关系及适用条件}
\textbf{卷积定理（convolution theorem）}说，在匹配的变换和边界约定下，空间卷积对应频谱逐点相乘：
\[
g=f*h\quad\Longleftrightarrow\quad G=FH.
\]
这里 $*$ 表示卷积，$h$ 是空间域滤波核，$H$ 是 $h$ 的傅里叶变换。以一维为例，$g(x)=\sum_t h(t)f(x-t)$；$t$ 枚举核的位置，每一项是权重乘邻近像素，再求和。

这个对应关系适用于固定核的\textbf{线性、平移不变}滤波。中值滤波需要先排序再选中间数，通常不能写成一个固定 $H$ 乘所有输入的频谱；同态滤波在取对数后的域里可用线性滤波，但整个原灰度输入到输出过程一般不是线性的。不能把“任何图像增强”都说成与某个固定频域乘法严格等价。

\subsection{为什么频谱要中心化：搬动索引，不是创造细节}

未经移动的 DFT 把 DC 放在数组左上角 $(0,0)$。低频不仅出现在左上角附近，也会绕到其他角落，肉眼较难看出“中心到外围”的规律。\textbf{频谱中心化（spectrum centering）}把零频率搬到图中央，便于显示和画圆形滤波器。

有两条常用、必须保持前后一致的路线：
\begin{enumerate}[label=\arabic*.]
\item 先算 DFT，再对频谱执行 \texttt{fftshift} 重排，把零频率搬到中央；在中心化频谱上乘 $H$ 后，用 \texttt{ifftshift} 搬回原顺序，再做 IDFT。
\item 当变换网格的两个尺寸都是偶数时，在正变换前把图像乘 $(-1)^{x+y}$；在中心化频谱上乘 $H$，做 IDFT 后再乘一次 $(-1)^{x+y}$，还原空间图像。
\end{enumerate}

$(-1)^{x+y}$ 是一张正负交替的棋盘：当 $x+y$ 为偶数时为 $+1$，为奇数时为 $-1$。
\[
\begin{bmatrix}
1&-1&1&-1\\-1&1&-1&1\\1&-1&1&-1\\-1&1&-1&1
\end{bmatrix}.
\]
它不是把图像整体变暗，而是利用傅里叶变换的调制性质，使频率索引沿两个方向各平移半个网格。对于偶数 $P\times Q$ 的变换网格，DC 移到 $(P/2,Q/2)$。再次乘棋盘能还原，因为 $[(-1)^{x+y}]^2=1$。

\textbf{为什么强调偶数尺寸？}当尺寸为奇数时，“半个网格”不是整数索引，这个棋盘调制不能简单等同于软件的整数格 \texttt{fftshift}。一般实现可采用成套的 \texttt{fftshift}/\texttt{ifftshift}，并按实际中心索引建立 $H$。中心化是方便，不是必需条件：也可以始终用未中心化坐标设计 $H$，只要滤波器与频谱的索引位置一致。

\textbf{最常见的错误：}对未中心化的 $F$，直接乘一个画在数组中央的低通圆盘。这样保留的会是错误位置的频率。另一个错误是只移动一次，逆变换前没有移回，结果带上交替符号。记住：\textbf{显示顺序可以改，配对关系不能错。}

\subsection{补零与边界：为什么右边会“绕回”左边}

\subsubsection{DFT 看待边界的方式}
有限长 DFT 把数组理解为一个周期的内容，周期结束后重新从开头接上。二维情况可以想象把同一张图片向四周反复平铺。因此在原尺寸网格上直接进行 DFT 相乘再 IDFT，得到的是\textbf{循环卷积（circular convolution）}：右边可能绕回左边，下边也可能绕回上边。

真实照片的左边和右边往往不连续，周期接缝会产生人为突变。滤波时这些接缝可能污染边界。因此“边界处理”既有数学上的卷积长度问题，也有图像外部应该取什么值的建模问题。

\subsubsection{用四则运算看懂循环卷积与线性卷积}
令一维图像 $f=[1,2,3]$，滤波核 $h=[1,1]$。把数组外面看成零，\textbf{线性卷积（linear convolution）}逐项计算为
\[
\begin{aligned}
g(0)&=1\times1=1,\\
g(1)&=2\times1+1\times1=3,\\
g(2)&=3\times1+2\times1=5,\\
g(3)&=3\times1=3.
\end{aligned}
\]
得到 $g=[1,3,5,3]$，长度为 $3+2-1=4$。如果只用长度 3 的 DFT，最后的结果没有位置存放，$g(3)=3$ 会绕回并加到 $g(0)$，得到循环卷积 $[4,3,5]$。这不是数值误差，而是计算了另一种边界模型。

正确的补零做法是把两者都放在至少长度 4 的共同网格：$f_p=[1,2,3,0]$，$h_p=[1,1,0,0]$。用四点 DFT 相乘再逆变换，就可以完整容纳 $[1,3,5,3]$，不发生首尾折叠。下标 $p$ 表示经过填充的数组。

\subsubsection{二维尺寸应怎样确定}
若图像有 $M\times N$ 个像素，有限滤波核有 $K\times L$ 个系数，那么完整线性卷积的尺寸为
\[
(M+K-1)\times(N+L-1).
\]
为了用 DFT 计算这个线性卷积，选共同变换尺寸 $P\times Q$ 满足
\[
P\geq M+K-1,\qquad Q\geq N+L-1,
\]
并把\textbf{图像和核都}填到这个共同尺寸。$P,Q$ 可以更大，以方便快速算法；不是必须取二的整数次幂。例如 $256\times256$ 图像与 $3\times3$ 核，至少需要 $258\times258$ 的共同网格；只把核补到 $256\times256$ 仍不足以容纳完整线性卷积。

此外，核的原点必须安排正确。以核左上角为索引零直接做完整卷积，随后应按核中心偏移选取与输入对齐的部分；若把核中心作为原点，则要按软件约定将正、负位移正确放入网格。\textbf{补零解决重叠，原点与裁剪解决对齐，两者不是一件事。}

\subsubsection{直接画频域滤波器，也不能跳过边界问题}
原稿称“若 $H$ 直接在频域构造，如 ILPF，则无此问题”，这不成立。只要在有限 DFT 网格上做乘法，逆变换仍对应这个网格上的循环卷积。理想低通的空间响应延伸很远，更不能因未显式写出 $h$ 就忽略边界。

若没有一个指定的有限核，便不能机械套用某个有限 $K,L$ 获得严格的全局等价。常见做法是扩展图像、在更大的网格上滤波，再裁取关心的区域；必要时结合镜像延拓等边界模型。补零可减少绕回污染，但零与边界像素之间也可能形成突变。应说明采用的是零延拓、镜像、复制边界还是周期延拓。

还有一种不同用途：为了\textbf{让频谱显示的采样更密}而给图像补零。它使频谱曲线看起来更平滑，却没有增加原图的细节信息，也没有凭空提高对两个邻近真实频率的分辨能力。必须分清“卷积不折叠”和“频谱插值更密”这两种目的。

\subsection{用半径选频率：三种低通滤波器}

\subsubsection{先把频谱位置换成到中心的距离}
对已中心化的偶数 $P\times Q$ 变换网格，定义
\[
\Delta u=u-P/2,\qquad \Delta v=v-Q/2,
\qquad D(u,v)=\sqrt{(\Delta u)^2+(\Delta v)^2}.
\]
$\Delta u,\Delta v$ 是相对中心的有符号偏移；$D$ 是到中心的欧氏距离；$D_0>0$ 是滤波器的尺度或约定截止半径。原题 E6 没有额外补零，因此网格记成 $M=N=256$；本模块用 $P,Q$ 区分原图尺寸和变换尺寸。

例如偏移 $(3,4)$ 的点，$D=\sqrt{3^2+4^2}=5$。以中心画半径 5 的圆，圆上所有方向都受到相同处理，这叫\textbf{径向对称}。上述 $D$ 的单位是“频率格”，并非直接的周期/毫米。如果图像两个方向尺寸或像素物理间距不同，要用归一化的物理频率计算圆形距离，不能把不同单位直接相加。

\textbf{低通滤波（low-pass filtering）}让慢变化通过，压低快变化。结果通常更平滑，但也可能模糊边缘、纹理和小目标。它不知道“哪种细节重要”，只按频率做选择。

\subsubsection{理想低通 ILPF：一刀切}
\textbf{理想低通滤波器（ideal low-pass filter，ILPF）}为
\[
H_{\mathrm I}(u,v)=
\begin{cases}
1,&D(u,v)\leq D_0,\\
0,&D(u,v)>D_0.
\end{cases}
\]
$H_{\mathrm I}$ 是理想低通的权重；$D_0$ 内全部保留，外面全部丢弃。本书约定边界 $D=D_0$ 取 1，采用别的边界约定时须说明。例如 $D_0=10$，$D=8$ 时 $H=1$，$D=12$ 时 $H=0$。优点是选频范围清楚；代价是过于突然的频率截断会引起后面解释的振铃。

原稿的“物理不可实现”需要限定：连续、无限范围的精确理想频率响应，不能由有限长度的空间核精确实现；但\textbf{有限 DFT 网格上的零一圆盘当然可以计算}。不能把它理解为计算机无法使用理想低通。

\subsubsection{巴特沃斯低通 BLPF：用阶数控制坡度}
\textbf{巴特沃斯低通滤波器（Butterworth low-pass filter，BLPF）}在本讲义中定义为
\[
H_{\mathrm B}(u,v)=\frac{1}{1+\left[D(u,v)/D_0\right]^{2n}},\qquad n\geq1.
\]
$n$ 是正整数阶数，控制过渡带有多陡；$D/D_0$ 是无单位的相对距离；指数是 $2n$。例如 $D_0=10$、$n=2$、$D=5$ 时，
\[
H_{\mathrm B}=\frac{1}{1+(5/10)^4}
=\frac{1}{1+1/16}=\frac{16}{17}\approx0.9412.
\]
所以这个分量保留约 $94.12\%$ 的振幅。若 $D=20$，则 $H=1/(1+2^4)=1/17\approx0.05882$，仅保留约 $5.88\%$。当 $D=D_0$ 时比值为 1，任何阶数都得到 $H=1/(1+1)=0.5$。

$n$ 越大，$D_0$ 内更接近 1，外部更接近 0，形状越接近理想低通。原稿建议的 $n=2$ 至 $3$ 可以作为尝试的起点，而非通用最优值。低阶通常较平缓，但不能无条件承诺“所有低阶实现绝无振铃”；滤波器参数、采样和边界处理都会影响实际结果。

\subsubsection{高斯低通 GLPF：平缓下降}
\textbf{高斯低通滤波器（Gaussian low-pass filter，GLPF）}为
\[
H_{\mathrm G}(u,v)=\exp\!\left[-\frac{D(u,v)^2}{2D_0^2}\right].
\]
$D_0$ 在这个写法中是频域的尺度参数。$D=0$ 时指数为 0，$H=1$；$D=D_0$ 时 $H=e^{-1/2}\approx0.60653$，不是 0.5；$D=2D_0$ 时 $H=e^{-2}\approx0.13534$。因此三种滤波器使用同一个 $D_0$，并不表示它们在该位置有相同的衰减。

\begin{center}
\begin{tikzpicture}[x=3.5cm,y=3.2cm,>=Latex,font=\small]
\draw[->] (0,0)--(2.75,0) node[right]{$D/D_0$};
\draw[->] (0,0)--(0,1.22) node[above]{$H$};
\draw[dotted] (1,0)--(1,1.08);
\draw[dotted] (0,.5)--(2.6,.5);
\draw[very thick] (0,1)--(1,1)--(1,0)--(2.6,0);
\draw[thick,dashed,domain=0:2.6,samples=100] plot(\x,{1/(1+\x^4)});
\draw[thick,dash dot,domain=0:2.6,samples=100] plot(\x,{exp(-\x^2/2)});
\foreach \x in {1,2}{\draw (\x,0)--(\x,-.025) node[below]{$\x$};}
\node[left] at (0,1){1};\node[left] at (0,.5){0.5};
\node[anchor=west,fill=white] at (1.4,1.08){实线：理想};
\node[anchor=west,fill=white] at (1.4,.94){虚线：巴特沃斯 $n=2$};
\node[anchor=west,fill=white] at (1.4,.80){点划线：高斯};
\end{tikzpicture}
\end{center}

这张图只画了从频谱中心往外走的一条半径，不是原图的灰度剖面。横轴往右表示变化更快；纵轴越低表示这种变化保留得越少。$D_0$ 变小，更多频率被压低，通常更模糊；$D_0$ 变大，更多细节得以保留，去噪也可能变弱。

\subsubsection{怎样选择半径：能量比例只是一个参考}
幅度的平方 $|F(u,v)|^2$ 常用来度量该频率的能量贡献。若希望了解半径 $D_0$ 内包含多少频谱能量，可计算
\[
\eta(D_0)=\frac{\displaystyle\sum_{D(u,v)\leq D_0}|F(u,v)|^2}
{\displaystyle\sum_{u,v}|F(u,v)|^2}.
\]
$\eta$ 为圆内能量占总能量的比例，分子只累加圆内频率，分母累加全部频率；分母为零的全零图像需单独处理。例如总能量为 1000，圆内为 920，则 $\eta=920/1000=0.92=92\%$。

原稿提出以包含约 $90\%$ 至 $95\%$ 能量的圆作选择参考，这可作为尝试参数的一种办法，\textbf{不是普遍最优标准}。若 DC 极强，它可能单独占绝大多数能量，极小半径也满足比例，却把重要边缘全部抹掉。实际还要结合目的、噪声分布与视觉结果；有时需单独分析去掉 DC 后的能量分布，并清楚说明采用的定义。

\subsubsection{纠错：0.5 到底是负 3 dB 还是负 6 dB}
\textbf{分贝（decibel，dB）}表达比值时，必须先说清是振幅比还是功率比。对振幅传递函数 $H$，
\[
A_{\mathrm{dB}}=20\log_{10}|H|.
\]
$A_{\mathrm{dB}}$ 是振幅增益的分贝数，$\log_{10}$ 是以 10 为底的对数。将 $H=0.5$ 代入，
\[
20\log_{10}(0.5)\approx-6.0206\ \mathrm{dB}.
\]
对应功率比为 $|H|^2=0.25$，用功率公式 $10\log_{10}(0.25)$ 仍得到 $-6.0206$ dB。\textbf{功率比为 0.5} 才对应 $10\log_{10}(0.5)\approx-3.0103$ dB，其振幅比为 $\sqrt{0.5}\approx0.7071$。

有些工程教材定义巴特沃斯\textbf{振幅}为 $1/\sqrt{1+(D/D_0)^{2n}}$，那时 $D_0$ 处确实约为 $-3$ dB。但原题使用的是\textbf{没有平方根}的 $1/[1+(D/D_0)^{2n}]$。不能保留这一公式，却搬用另一种定义的截止结论。

\subsection{为什么边界会出现明暗晕圈：振铃不是“还有噪声”}

\textbf{振铃（ringing）}是边缘附近出现交替的过亮、过暗波纹。原图即使完全没有噪声，经过陡峭截止的滤波器，也可能出现振铃。因此不能看见波纹就断言去噪没有去干净。

频域与空间域通过傅里叶变换联系：频域截断越突然，对应的空间核越容易长距离振荡。对一维理想低通，空间核是 sinc 型；二维\textbf{圆形}理想低通的空间核严格说是带贝塞尔函数的 jinc 型径向函数。两者都具有旁瓣且包含负值。把二维圆盘直接称为“二维 sinc”只是一种不精确的类比。

不必先计算 sinc，也可以用一个有负系数的小核看清后果。考虑从 0 突然跳到 1 的阶跃：$f(x)=0$（$x<0$），$f(x)=1$（$x\geq0$）。取仅用于示意的对称核
\[
h(-2)=-0.1,\quad h(-1)=0.2,\quad h(0)=0.8,
\quad h(1)=0.2,\quad h(2)=-0.1.
\]
系数总和为 1，能保持远离边缘处的常数；但负系数可能把输出拉到原范围之外。按 $g(x)=\sum_{t=-2}^{2}h(t)f(x-t)$，得到
\[
\begin{array}{c|rrrrrrr}
x&-3&-2&-1&0&1&2&3\\\hline
f(x)&0&0&0&1&1&1&1\\
g(x)&0&-0.1&0.1&0.9&1.1&1&1
\end{array}
\]
例如 $x=-2$ 时仅 $t=-2$ 那项看到亮侧的 1，故 $g(-2)=-0.1$，出现欠冲；$x=1$ 时前四个系数都落在亮侧，故 $g(1)=-0.1+0.2+0.8+0.2=1.1$，出现过冲。真实理想低通具有更多旁瓣，所以可能反复交替形成多道波纹。

\begin{center}
\begin{tikzpicture}[x=.95cm,y=2.1cm,>=Latex,font=\small]
\draw[->] (-3.5,0)--(3.8,0) node[right]{$x$};
\draw[->] (-3.3,-.28)--(-3.3,1.4) node[above]{灰度};
\draw[gray,dashed] (-3,1)--(3.3,1);
\draw[thick,dashed] (-3,0)--(0,0)--(0,1)--(3,1);
\draw[very thick] plot coordinates{(-3,0)(-2,-.1)(-1,.1)(0,.9)(1,1.1)(2,1)(3,1)};
\foreach \x/\y in {-3/0,-2/-.1,-1/.1,0/.9,1/1.1,2/1,3/1}{\fill (\x,\y) circle(1.5pt);}
\draw[->] (-1.6,-.35)--(-2,-.12) node[pos=0,below]{低于暗侧};
\draw[->] (1.8,1.36)--(1,1.12) node[pos=0,above]{高于亮侧};
\node[align=left] at (4.4,.55){虚线：原阶跃\\实线：示意输出};
\end{tikzpicture}
\end{center}

作为对照，若使用非负、总和为 1 的核
\[
h_{\mathrm{pos}}=\frac1{16}[1,4,6,4,1],
\]
那么每个输出是邻域像素的加权平均，不会超出邻域最小值和最大值。对同一阶跃，过渡值依次为 $1/16,5/16,11/16,15/16,1$，单调地从 0 过渡到 1，没有过冲、欠冲。高斯空间核同样非负，这解释了常说的“高斯低通无振铃”。

这一结论指理想高斯核及恰当实现本身不因负旁瓣产生振铃；边界延拓、核截断、其他锐化步骤仍可能引入伪影。高斯也会让边缘过渡变宽，“没有晕圈”不等于“完全不模糊”。

\subsection{高通、带通、带阻和陷波：先判断想保留什么}

\subsubsection{高通不是自动得到一张好看的锐化图}
\textbf{高通滤波（high-pass filtering）}保留较快变化，压低较慢变化。本模块常用低通的补函数构造：
\[
H_{\mathrm{HP}}(u,v)=1-H_{\mathrm{LP}}(u,v).
\]
$H_{\mathrm{LP}}$ 是选定低通响应，$H_{\mathrm{HP}}$ 是对应高通响应。分别代入理想、巴特沃斯和高斯低通，就得到 IHPF、BHPF 和 GHPF，即三种对应的高通滤波器。例如某点低通为 0.8，高通就是 0.2；某点低通为 0.1，高通就是 0.9。低通在 DC 为 1，故这种高通在 DC 为 0，去掉整幅图的均值。

纯高通输出常包含正负值，主要表达边缘和细节，失去原来的整体亮度。若希望保留背景并增强细节，可以采用
\[
g=f+k\,g_{\mathrm{HP}},\qquad
H_{\mathrm{sharp}}=1+kH_{\mathrm{HP}},\qquad k>0.
\]
$g_{\mathrm{HP}}$ 是高通输出，$k$ 是加入细节的强度。以四像素例子为例，原图 $[12,10,8,10]$，纯高通为 $[2,0,-2,0]$，取 $k=0.5$，得到 $[13,10,7,10]$。平均仍为 10，局部起伏增大。

原稿说“高通放大噪声”，更准确的理解是：高频噪声会相对低频背景更突出；当使用 $0\leq H_{\mathrm{HP}}\leq1$ 的纯高通时，不是每个噪声频率的绝对振幅都增大。高提升响应 $1+kH_{\mathrm{HP}}$ 可大于 1，此时相应频率的信号和噪声都确实被放大。

\subsubsection{带阻与带通：按一整圈频率选择}
\textbf{带阻滤波（band-reject filtering）}压低某个频率范围，两侧通过。以径向理想带阻为例，
\[
H_{\mathrm{BR}}(D)=
\begin{cases}
0,&D_0-W/2\leq D\leq D_0+W/2,\\
1,&\text{其他}.
\end{cases}
\]
$D_0$ 是被抑制环带的中心半径，$W>0$ 是带宽，通常选 $D_0>W/2$ 以使内半径为正。例如 $D_0=20,W=6$，被抑制的是 $17\leq D\leq23$，不是“所有 $D>20$”。\textbf{带通滤波（band-pass filtering）}反过来只保留这个频带：$H_{\mathrm{BP}}=1-H_{\mathrm{BR}}$。

原稿给出的巴特沃斯带阻为
\[
H_{\mathrm{BR,B}}(D)=
\frac{1}{1+\left[\dfrac{DW}{D^2-D_0^2}\right]^{2n}}.
\]
$D,D_0,W$ 的含义同上，正整数阶数 $n$ 控制过渡陡峭程度。$D=D_0$ 处不能直接除以零，取极限规定 $H=0$。例如 $D_0=10,W=4,n=1,D=5$，括号内为 $20/(25-100)=-4/15$，平方后为 $16/225$，因此 $H=225/241\approx0.9336$。该点仍可通过约 $93.36\%$ 的振幅。

\begin{center}
\begin{tikzpicture}[x=1cm,y=1.5cm,font=\small,>=Latex]
\begin{scope}
\draw[->] (0,0)--(4.5,0) node[right]{$D$};
\draw[->] (0,0)--(0,1.3) node[above]{$H$};
\draw[thick] (0,1)--(1.5,1)--(1.5,0)--(2.8,0)--(2.8,1)--(4.2,1);
\node at (2,-.55){带阻：中间这段不通过};
\node[left] at (0,1){1};
\draw[<->] (1.5,.45)--(2.8,.45) node[midway,above]{$W$};
\end{scope}
\begin{scope}[xshift=7cm]
\draw[->] (0,0)--(4.5,0) node[right]{$D$};
\draw[->] (0,0)--(0,1.3) node[above]{$H$};
\draw[thick] (0,0)--(1.5,0)--(1.5,1)--(2.8,1)--(2.8,0)--(4.2,0);
\node at (2,-.55){带通：只让中间这段通过};
\node[left] at (0,1){1};
\end{scope}
\end{tikzpicture}
\end{center}

原稿还提到高斯型带阻：它也以平滑过渡替代突然的零一切换。理想、巴特沃斯、高斯这些名字描述的是过渡形状；带通、带阻描述的是保留哪个频段，两组分类不要混淆。

带阻可能消除某种周期干扰，但同时会删去同一半径上各个方向的真实纹理。若噪声只表现为两个局部亮点，抑制整圈通常过宽。

\subsubsection{陷波：只处理被确认的几个频率位置}
\textbf{陷波滤波（notch filtering）}是在特定频率位置附近开小孔。令中心化频率坐标为 $\boldsymbol q=(\Delta u,\Delta v)$，识别出的干扰位置为 $\boldsymbol q_0=(a,b)$。理想成对陷波可写成
\[
H_{\mathrm{notch}}(\boldsymbol q)=
\begin{cases}
0,&\|\boldsymbol q-\boldsymbol q_0\|\leq r
\ \text{或}\ \|\boldsymbol q+\boldsymbol q_0\|\leq r,\\
1,&\text{其他}.
\end{cases}
\]
$\|\cdot\|$ 表示欧氏距离，$r$ 是每个小孔的半径。例如干扰位于偏移 $(0,20)$ 与 $(0,-20)$，可以在两点附近各压低半径 2 的区域，而不把半径 20 的整圈都删掉。实际常用平滑陷波边缘，减少突截带来的振铃。

\begin{center}
\begin{tikzpicture}[font=\small]
\begin{scope}
\fill[gray!10] (-2,-2) rectangle(2,2);
\fill[gray!65,even odd rule] (0,0) circle(1.5) (0,0) circle(1.05);
\draw (-2,-2) rectangle(2,2);\draw[gray](-2,0)--(2,0);\draw[gray](0,-2)--(0,2);
\fill (0,0) circle(1.7pt);
\node[align=center] at (0,-2.55){环形带阻\\灰色环带被抑制};
\end{scope}
\begin{scope}[xshift=6.5cm]
\fill[gray!10] (-2,-2) rectangle(2,2);
\fill[gray!65] (-1.3,0) circle(.25);\fill[gray!65] (1.3,0) circle(.25);
\draw (-2,-2) rectangle(2,2);\draw[gray](-2,0)--(2,0);\draw[gray](0,-2)--(0,2);
\fill (0,0) circle(1.7pt);
\node[align=center] at (0,-2.55){成对陷波\\仅局部小孔被抑制};
\end{scope}
\end{tikzpicture}
\end{center}

对于实输入，要保持输出为实数，应让 $H$ 保持相应的共轭对称性。实数、偶对称的成对小孔是最方便的做法；只去掉一边通常破坏配对，使 IDFT 出现不能仅归因于舍入误差的虚部。DC 是自身的共轭伙伴，只令中心 $H=0$ 是一种特殊的直流陷波，作用是去均值，\textbf{不是通用的周期条纹去除方法}。

\subsubsection{噪声与高频并非同义词}
边缘和草地纹理是有用的高频信息；条纹干扰可能落在低频、中频或高频；白噪声通常在很宽的频率范围都有成分。图像上的缓慢阴影可能主要影响低频，也可能根本不是噪声。选择滤波器前应先问：干扰的频率在哪里？它与有用信号是否重叠？

若二者在同一频率位置重叠，简单频域乘法不能知道哪个是噪声、哪个是物体。抑制噪声也会损失信号。因此“外圈全部是噪声”“所有对称亮点都该删掉”“低通一定去噪且保边”都不成立。

\subsection{同态滤波：先解决“光照乘在表面上”这个问题}

\subsubsection{照明与反射为什么是相乘}
\textbf{同态滤波（homomorphic filtering）}从一个简化成像模型出发：
\[
f(x,y)=i(x,y)\,r(x,y).
\]
$f$ 是记录到的灰度，$i>0$ 是照明强度，$r$ 是表面反射比例；为使后面取对数有明确含义，$f$ 和 $i$ 均理解为选定参考尺度下的无量纲正数，$r$ 本身就是无量纲比例。在简化模型中常取 $0<r\leq1$。同一表面得到两倍的照明，记录亮度也大致变为两倍，这就是乘法的直观来源。例如 $i=100,r=0.3$ 时 $f=30$；照明升为 200，反射仍为 0.3，则 $f=60$。

常用假设是：照明在空间上变化较缓慢，反射在物体边界和纹理处变化较快。因此希望压低缓慢的大范围亮暗差，同时加强细节。\textbf{这只是近似：}阴影边界也可能很锐利，反射率也可能缓慢变化；两者的频率范围并不一定严格分开。

\subsubsection{为什么先取对数，而不是直接做高通}
直接对 $f=i\,r$ 变换，乘积对应频域中的卷积，照明与反射不会简单地成为两个可各自加权的频谱。对数有一个关键性质：$\ln(ab)=\ln a+\ln b$。所以令
\[
z(x,y)=\ln f(x,y)=\ln i(x,y)+\ln r(x,y).
\]
$z$ 是对数灰度。对它做 DFT，得到
\[
Z(u,v)=I_{\log}(u,v)+R_{\log}(u,v).
\]
$Z$ 是 $z$ 的频谱，$I_{\log}$ 是 $\ln i$ 的频谱，$R_{\log}$ 是 $\ln r$ 的频谱。现在两部分在数学上相加，频域滤波可写成
\[
S=HZ=H I_{\log}+H R_{\log}.
\]
$S$ 是滤波后的对数域频谱。同一个 $H$ 根据频率给予不同权重，在“照明偏低频、反射细节偏高频”的假设下起作用。\textbf{取对数并没有自动、精确地识别出照明和反射}；两者频谱若重叠，仍会相互影响。

数值上，若 $i=100,r=0.25$，则 $f=25$；$\ln100\approx4.60517$、$\ln0.25\approx-1.38629$，二者之和为 $3.21888\approx\ln25$。反射率的对数可能为负，这完全正常，不需要为此删掉负值。

\subsubsection{常用滤波器的每个参数是什么意思}
原稿采用
\[
H(D)=(\gamma_H-\gamma_L)
\left[1-\exp\!\left(-cD^2/D_0^2\right)\right]+\gamma_L.
\]
$D$ 是到频谱中心的距离；$D_0>0$ 是过渡尺度；$c>0$ 调节过渡速度；$\gamma_L$ 是低频端的增益；$\gamma_H$ 是高频端的增益。常见增强设置是 $0<\gamma_L<1<\gamma_H$。

先检查两端：当 $D=0$ 时，指数为 0，方括号为 $1-1=0$，故 $H(0)=\gamma_L$。当 $D$ 很大时，指数趋向很大的负数，指数函数趋近零，方括号趋向 1，故 $H(D)\to\gamma_H$。这正是“低频压低，高频提高”的来源。

例如取 $\gamma_L=0.5,\gamma_H=1.5,c=1$，则 $H(0)=0.5$；在 $D=D_0$ 处，
\[
H(D_0)=(1.5-0.5)(1-e^{-1})+0.5
\approx1.13212;
\]
远离中心时趋近 1.5。这里“低频减半”指\textbf{对数域分量的振幅减半}，不代表把原图每个暗区域的像素直接除以二。

\begin{center}
\begin{tikzpicture}[x=3.3cm,y=2cm,>=Latex,font=\small]
\draw[->] (0,0)--(2.8,0) node[right]{$D/D_0$};
\draw[->] (0,0)--(0,1.8) node[above]{$H$};
\draw[dotted] (0,.5)--(2.6,.5);\draw[dashed] (0,1.5)--(2.6,1.5);
\draw[thick,domain=0:2.6,samples=90] plot(\x,{1-exp(-\x^2)+.5});
\node[left] at (0,.5){$\gamma_L=0.5$};\node[left] at (0,1.5){$\gamma_H=1.5$};
\node[below] at (1,0){1};\draw[dotted] (1,0)--(1,1.1321);
\node[align=center,fill=white,inner sep=1pt] at (1.6,.38){压低缓慢变化\\提高快速变化};
\end{tikzpicture}
\end{center}

\subsubsection{完整流程与“动态范围压缩”的含义}
完整流程为：准备正值图像，取自然对数，按需要扩展边界与中心化，计算 DFT，乘 $H$，逆变换并正确移回，得到 $s(x,y)$，最后指数还原：
\[
g(x,y)=\exp\bigl(s(x,y)\bigr).
\]
$g$ 是输出图像，$s=\operatorname{IDFT}(S)$ 是滤波后的对数图。空间尺寸裁取和显示范围处理应与其他频域滤波一样明确说明。

如果某个正分量在对数域整体乘以 $\alpha$，指数还原后变成 $\exp(\alpha\ln i)=i^\alpha$。例如照明从 1 到 4，亮暗比为 4；取 $\alpha=0.5$ 后变成从 1 到 2，亮暗比降为 2，这就是乘法范围被压缩的直观含义。它不保证所有像素绝对变暗：对于 $0<i<1$，平方根反而比原值大。因此应说“压缩相对变化范围”，不能仅凭 $\gamma_L<1$ 就断言整幅图平均灰度必定下降。

\subsubsection{重要订正：两个增益相等，不等于普通线性滤波}
若 $\gamma_L=\gamma_H=\gamma$，前式第一项为零，$H(D)=\gamma$ 对所有频率都是常数。于是
\[
s=\gamma\ln f,\qquad g=e^{\gamma\ln f}=f^\gamma.
\]
当 $\gamma=1$ 时，输出等于输入；一般 $\gamma\neq1$ 时得到\textbf{幂律变换}，对原始灰度仍然非线性。例如 $\gamma=2$ 时，输入 2 和 3 分别输出 4 和 9，而输入和 5 输出 25，不等于 $4+9=13$，所以不满足线性的可加性。

原稿写成“$\gamma_L=\gamma_H$ 时退化为普通线性滤波”，需要改成：“在\textbf{对数域}退化为常数增益；在\textbf{原灰度域}通常是幂律非线性变换。”即使增益不同，整个“对数、线性滤波、指数”的组合一般也不是原图上的线性系统。

\subsubsection{零像素、噪声和模型限制}
$\ln0$ 没有有限值，不能直接对含零图像照搬公式。常见实现是先选尺度并使用 $z=\ln(\max(f,\varepsilon))$，其中 $\varepsilon>0$ 为小的正下限；这样明确改变了零附近的数据。另一种实现是 $\ln(1+f)$，但它不再严格满足 $\ln(ir)=\ln i+\ln r$ 的原模型，逆变换后要配套使用 $e^s-1$。

不能把 $\varepsilon$ 看成随便写的魔法数：它相对于灰度尺度多大，会影响暗部。接近零的值对微小噪声很敏感，高频提升也可能强化纹理噪声；过强参数可产生晕圈或不自然对比度。若图像严重饱和、照明与反射频谱高度重叠，简单同态模型不能保证恢复真实表面。

\subsection{原题 B4：频谱图判读，逐问建立推理链}
\label{original:B4}

\noindent\textbf{原题 B4（Enhancement3）}\quad
频谱图判读。(1) 中心亮点代表什么？(2) 远离中心的暗区代表什么？(3) 平坦的天空与密集的草地分别落在哪一部分？(4) 过中心出现一对对称亮斑，最可能是什么原因，用什么方法去除？

\paragraph{先确认题目的图是什么。}
本题通常指经过中心化、可能经过对数压缩显示的\textbf{幅度谱}。若图未中心化，DC 不在中央；若图显示的是相位，则“亮就是强”的读法也不成立。以下按中心化幅度谱作答。

\paragraph{(1) 中心亮点。}
中心是零频率，即直流分量，代表整体平均亮度。按照本模块 DFT 约定，中心系数的数值为全部像素之和，除以 $MN$ 得到平均值。中心\textbf{附近}的较亮区域还包括低频缓慢变化，不能把一大片中心区域都称作一个 DC 点。

\paragraph{(2) 远离中心的暗区。}
“远离中心”说明这里是较高的空间频率；“暗”说明这些频率的幅度较小。二者须分开解释。不能答成“原图这里很暗”，也不能只说“这里是高频，所以必然是噪声”。若幅度严格为零，则该离散频率分量缺失；显示为黑也可能仅是动态范围压缩后太弱而看不出来。

\paragraph{(3) 天空与草地。}
平坦天空在空间上变化慢，主要贡献 DC 和中心附近的低频；草地具有密集明暗变化，通常具有更多中高频成分。因此草地可能使外围某些区域增强。但草地也有平均亮度与大范围明暗，仍有低频；天空与地面的清晰分界也会贡献高频。不要说“一个物体只属于某一个频带”。

\paragraph{(4) 对称亮斑的原因与处理。}
如果原图同时可见不希望保留的等间距条纹，这对亮斑很可能来自周期干扰。实余弦以正、负频率配对，因此出现关于 DC 对称的亮点；两点连线垂直于条纹走向；离中心越远，条纹周期越短。

确认其为干扰后，优先在亮斑处使用\textbf{成对、局部、适度平滑的陷波}；若干扰能量形成完整环带，才考虑带阻。对称亮点也可能来自合法的栅格、布料或建筑纹理，所以仅凭频谱形状不能证明是噪声。频谱选点应结合原图和处理目标。

\paragraph{可整理成的短答案。}
中心为 DC，代表平均亮度；外围暗区为较弱的高频分量。平坦天空主要贡献低频，细密草地含较多中高频。对称亮点提示周期成分；若确认为条纹干扰，用成对陷波，必要时用带阻，并留意真实纹理损失。

\subsection{原题 D6：完整流程，重点是顺序与边界}
\label{original:D6}

\noindent\textbf{原题 D6}\quad
写出频域滤波的完整流程；并说明：(1) 为什么变换前要乘 $(-1)^{x+y}$？(2) 为什么补零到 $M\times N$？(3) 什么时候该用频域而不是 $3\times3$ 掩模？

\paragraph{先订正题目中尺寸的歧义。}
若题目的 $M\times N$ 指\textbf{最终选定的共同变换网格}，补到它是为让图像与核尺寸一致，但这个网格还必须足够大。若 $M\times N$ 指\textbf{原图尺寸}，则“核补到这个尺寸就一定得到线性卷积”不成立。下面统一用 $M\times N$ 表示原图、$P\times Q$ 表示变换网格。

\paragraph{完整流程：选择一条自洽的中心化路线。}
\begin{enumerate}[label=\arabic*.]
\item \textbf{明确目标与边界。}决定是平滑、提取细节还是去周期条纹；说明用零延拓、镜像还是周期边界。
\item \textbf{确定尺寸并准备数据。}若以 DFT 实现有限 $K\times L$ 核的完整线性卷积，选 $P\geq M+K-1,Q\geq N+L-1$，把图像和核填到共同尺寸，并处理核原点。
\item \textbf{正变换。}计算填充图像的 DFT，再用 \texttt{fftshift} 中心化；或者在偶数网格上先乘棋盘再做 DFT。二选一，不要把两条路线重复使用。
\item \textbf{建立 $H$ 并相乘。}让 $H$ 的中心、尺寸、频率单位与 $F$ 完全一致，逐点计算 $G=HF$。例如 $H=0.3$ 的频率位置保留原振幅的 $30\%$。
\item \textbf{逆变换并撤销移动。}使用频谱重排路线时先 \texttt{ifftshift} 再 IDFT；使用棋盘路线时先 IDFT 再乘一次棋盘。
\item \textbf{检查与截取。}对实输入和共轭对称滤波器，虚部应只剩微小舍入量；若虚部很大，先检查配对和索引。按核中心与输出需求裁取空间区域。
\item \textbf{单独处理显示范围。}若输出必须显示为 $0$ 到 $L_{\mathrm g}-1$ 的灰度，再选择缩放或截断；$L_{\mathrm g}$ 是可用灰度级数。保留负值的高通结果不宜未经说明就截断。
\end{enumerate}

这里“裁取到原图大小”是\textbf{空间尺寸}操作；“限制到灰度 $[0,L_{\mathrm g}-1]$”是\textbf{数值范围}操作。原稿把它们压在“裁剪”二字里，初学者需要明确区分。截断和归一化还会改变数值，不属于卷积定理本身。

\paragraph{(1) 乘棋盘的原因。}
在偶数网格中，乘 $(-1)^{x+y}$ 把 DC 平移到中心，方便用到中心的距离构造滤波器。它不是去噪，也不是增加高频。它只是可选的一种坐标安排，逆过程要乘回去；使用正确的 \texttt{fftshift}/\texttt{ifftshift} 也可以。

\paragraph{(2) 补零的原因与充分条件。}
DFT 乘法自然给出循环卷积。为求零延拓下的完整线性卷积，需要让两数组共同放入能容纳完整输出的网格。例：$256\times256$ 图像与 $3\times3$ 核至少用 $258\times258$。仅“尺寸一样”是可相乘的必要条件，\textbf{不等于}“尺寸足够大”这个避免折叠的条件。直接构造频域 $H$ 也仍需考虑周期边界。

\paragraph{(3) 什么时候频域更合适。}
$3\times3$ 固定核只需要每个像素约九次加权累加，空间域实现通常简单且开销小。大核卷积、需要精确选频的带阻与陷波、或者同态流程中的滤波，往往更适合考虑频域。

若图像为 $M\times N$、方核边长为 $k$，直接空间卷积粗略计算量为 $O(MNk^2)$；对 $P\times Q$ 网格做 FFT 的典型数量级为 $O(PQ\log(PQ))$。符号 $O$ 描述数据规模增大时的增长阶，不是实际耗时秒数。填充、缓存、硬件、可分离核和是否复用频谱都会影响速度，所以不存在不看实现便保证的统一分界。

\subsection{原题 D7：把同态滤波说清楚，避免四个跳步}
\label{original:D7}

\noindent\textbf{原题 D7}\quad
同态滤波。(1) 写出模型，说明照明与反射各对应低频还是高频；(2) 为什么必须先取对数？(3) 写出流程与 $H(u,v)$，并说明 $\gamma_L$、$\gamma_H$ 的作用。

\paragraph{(1) 模型与假设。}
写 $f(x,y)=i(x,y)r(x,y)$，再说明 $i$ 为照明、$r$ 为反射率。常假设照明变化慢，主要影响低频；反射的边界与纹理变化快，含较多高频。不能省略“常假设”，也不能声称反射率只有高频、照明绝无高频。

\paragraph{(2) 为什么取对数。}
原模型是相乘，而想利用的是频域中可加成分的线性加权。先取 $z=\ln f=\ln i+\ln r$，再变换成 $Z=I_{\log}+R_{\log}$，才能在这个加性表示中用同一 $H$ 对不同频率加不同增益。对数不保证两部分完全分离，只是让这种建模方式可用。含零像素要先选定正值预处理方法。

\paragraph{(3) 流程、公式和参数。}
流程写为
\[
f\ \xrightarrow{\ln}\ z\ \xrightarrow{\mathrm{DFT}}\ Z
\ \xrightarrow{\times H}\ S\ \xrightarrow{\mathrm{IDFT}}\ s
\ \xrightarrow{\exp}\ g.
\]
图中 $z$ 是对数图，$Z$ 是其频谱，$S$ 是滤波后的频谱，$s$ 是滤波后的对数图，$g$ 是输出；边界扩展、频谱移位、移回与裁取均应按前述规则配套执行。滤波器为
\[
H(u,v)=(\gamma_H-\gamma_L)
\left[1-e^{-cD(u,v)^2/D_0^2}\right]+\gamma_L.
\]
$D$ 为中心距离，$D_0,c$ 控制过渡；$\gamma_L<1$ 抑制对数域低频的相对变化，常用于减弱大范围照明差；$\gamma_H>1$ 提升对数域高频，常用于增强细节。由 $H(0)=\gamma_L$ 和远处 $H\to\gamma_H$ 验证，而不是只背两个不等式。

\paragraph{补充检验。}
令 $\gamma_L=\gamma_H=0.5$，得到 $g=f^{0.5}$。例如灰度 4 和 9 分别变成 2 和 3，说明这时是开平方的非线性变换，而不是普通线性低通。若两者同时为 1，才是这一设置下的恒等变换。

\subsection{原题 E6：逐步计算距离、增益与振铃}
\label{original:E6}

\noindent\textbf{原题 E6}\quad
频率域计算 + 振铃。$M=N=256$，频谱已中心化，$D_0=32$，某点偏移 $(\Delta u,\Delta v)=(24,32)$。(1) 求 $D$；(2) 求 ILPF、BLPF（$n=2$）、GLPF 在该点的 $H$；(3) 用 ILPF 平滑含边界的图像，为什么边界两侧出现一圈圈明暗环？哪种滤波器不振铃？

\paragraph{先读题：哪些数有用。}
$256\times256$ 告诉我们网格中心在 $(128,128)$。题目已经给出\textbf{相对中心的偏移} $(24,32)$，所以求距离时直接使用 24 和 32，不能再各减 128。若题目给的是绝对坐标 $(152,160)$，才先计算 $152-128=24$、$160-128=32$。

\paragraph{(1) 求半径。}
\[
\begin{aligned}
D&=\sqrt{(\Delta u)^2+(\Delta v)^2}\\
 &=\sqrt{24^2+32^2}=\sqrt{576+1024}
 =\sqrt{1600}=40.
\end{aligned}
\]
先平方、相加、开方，不能把距离算成 $24+32=56$。$D=40$ 与 $D_0=32$ 单位相同，均为频率格；其比值为 $D/D_0=40/32=5/4=1.25$。

\paragraph{(2a) 理想低通。}
它只需要比较内外：$40>32$，该点在通带圆盘外，因此
\[
\boxed{H_{\mathrm I}=0}.
\]
这意味着此频率系数被完全删除，不意味着原图某个像素被设为零。

\paragraph{(2b) 巴特沃斯低通，阶数为 2。}
先确定指数：$2n=2\times2=4$。然后代入，尽量到最后才舍入：
\[
\begin{aligned}
H_{\mathrm B}
&=\frac{1}{1+(D/D_0)^{2n}}
=\frac{1}{1+(40/32)^4}\\
&=\frac{1}{1+(5/4)^4}
=\frac{1}{1+625/256}
=\frac{256}{881}\\
&\approx\boxed{0.2905788876}.
\end{aligned}
\]
该频率保留约 $29.06\%$ 的\textbf{振幅}。若输入该点的幅度为 100，输出约为 $100\times0.2905788876=29.0579$；功率保留比例是这个数的平方，约 $0.08444$，不能把振幅百分比直接当成功率百分比。

\paragraph{(2c) 高斯低通。}
\[
\begin{aligned}
H_{\mathrm G}
&=\exp\!\left[-\frac{D^2}{2D_0^2}\right]
=\exp\!\left[-\frac{40^2}{2\times32^2}\right]\\
&=\exp\!\left[-\frac{1600}{2048}\right]
=\exp(-0.78125)\\
&\approx\boxed{0.4578333618}.
\end{aligned}
\]
同一频率保留约 $45.78\%$ 振幅。原稿的三位小数结果 $0.291$、$0.458$ 可以作为四舍五入答案；完整计算应保留上面的精确分数与指数，避免中间过早截断。

\begin{center}
\begin{tabular}{@{}llll@{}}
\toprule
滤波器&代入结果&约保留振幅&本点的解释\\
\midrule
ILPF&$0$&$0\%$&圆盘外完全去掉\\
BLPF，$n=2$&$256/881$&$29.06\%$&逐渐衰减\\
GLPF&$e^{-0.78125}$&$45.78\%$&平缓衰减\\
\bottomrule
\end{tabular}
\end{center}

本题不能据此断言“高斯在所有位置都比巴特沃斯保留更多”。比较只针对当前半径和参数；两条曲线在其他位置可能有不同的相对关系。

\paragraph{(3) 为什么产生明暗环，哪一种避免这种振铃。}
ILPF 在 $D_0$ 处突然从 1 变为 0；对应空间核具有振荡旁瓣，且包含负值。与阶跃边缘卷积时，一部分旁瓣跨过边界，把亮侧与暗侧按正负权重混合，产生过冲与欠冲。二维圆形低通可能在边界附近出现环状或平行于边界的明暗晕纹，这就是振铃。

在所比较的标准滤波器中，\textbf{高斯低通的理想空间核仍为非负高斯，没有负旁瓣，因此不会产生这种旁瓣振铃}；它仍会模糊边缘。巴特沃斯过渡较缓时通常振铃较轻，提高阶数趋近理想截止，振铃风险相应增大。实际实现还要检查边界处理。

\paragraph{本题的截止点纠错。}
$D=D_0$ 时，按本题公式，BLPF 始终为 $0.5$，对应振幅增益约 $-6.02$ dB；GLPF 为 $e^{-1/2}\approx0.60653$。原稿把 BLPF 的 $0.5$ 标成 $-3$ dB，是混用了振幅比、功率比或另一种巴特沃斯定义。

\subsection{做题前的检查表与自测}

\begin{enumerate}[label=\arabic*.]
\item \textbf{先问坐标。}看到的是像素图、幅度谱还是相位谱？是否中心化？横竖轴分别对应哪一方向？
\item \textbf{再问频率。}位置远近说明变化快慢，亮暗说明幅度强弱；不要把两者混为一谈。
\item \textbf{再问滤波目标。}平滑用低通；提取变化用高通；指定尺度用带通；环形干扰考虑带阻；局部成对周期分量考虑陷波。
\item \textbf{代入前看定义。}低通的 $D_0$ 不是跨所有公式通用的半功率点；BLPF 的幂是 $2n$，GLPF 的分母是 $2D_0^2$。
\item \textbf{逆变换前后看配套。}中心化要移回；图像和核要共同够大；核原点要对齐；虚部异常要检查，不能直接掩盖。
\item \textbf{最后看副作用。}平滑可能损失边缘；高提升可能放大噪声；硬截止可能振铃；对数模型需要正值且不保证照明与反射可完全分开。
\end{enumerate}

\paragraph{自测 1：}一幅图有间隔 8 像素的竖条纹，另一幅图间隔 4 像素。哪幅的亮点更远离中心？亮点连线大致朝哪里？\quad\textbf{答案：}4 像素周期的频率更高，亮点离中心更远；竖条纹沿水平方向发生变化，连线大致水平。比较需使用相同尺寸与采样单位。

\paragraph{自测 2：}某点在中心附近，但频谱显示很暗，能说“它是高频”吗？\quad\textbf{答案：}不能。它的位置说明频率低，暗说明该低频分量幅度小。

\paragraph{自测 3：}$100\times80$ 图像与 $5\times7$ 核求完整线性卷积，共同网格至少多大？\quad\textbf{答案：}$(100+5-1)\times(80+7-1)=104\times86$，图像和核都须放入该网格。若之后裁回原图大小，还要说明核中心与裁取位置。

\paragraph{自测 4：}同态滤波中 $\gamma_L=\gamma_H=2$，输出是什么？\quad\textbf{答案：}在正值、无额外偏移与显示缩放的模型下，$g=f^2$，不是原灰度域的线性滤波。若同时为 1，输出才保持不变。

\paragraph{这一模块的一句话主线。}\quad
先把图像理解为不同快慢、不同方向的亮度变化；DFT 给出这些变化的配方；$H$ 决定各成分保留多少；IDFT 把新配方还原成图像。频谱中心化负责坐标，补零与延拓负责边界，滤波器形状负责取舍，而相位、真实纹理与噪声重叠等限制决定了哪些信息不能凭一个简单滤波器恢复。

\section{原题H：用简短题目检查概念是否建立}
本章保留原稿H类题目，并把选项背后的理由展开。选择题的答案往往依赖题目的简化条件；理解这些条件，比只记选项字母更可靠。

\subsection{H1：椒盐噪声为什么优先使用中值滤波}
\label{original:H1}
\textbf{原题：}下列哪个模板可用于去除椒盐噪声？A. $3\times3$均值；B.中值滤波；C.拉普拉斯；D.高提升。

在本题的典型孤立脉冲噪声情形下，选B。椒盐噪声表现为个别异常亮点或暗点。中值滤波先把邻域的数按大小排序，再取中间的数；只要异常值没有占据足够多的位置，最大的255或最小的0就不会被选为中间值。例如$[98,99,100,101,255]$的中值是100，而均值为130.6。均值会受到极大值的拉动。

拉普拉斯和高提升主要突出变化，可能进一步放大噪声。中值也不是对任意密度的椒盐噪声都有效：异常点太多或目标太细时，仍可能失效或删除真实细节。

\subsection{H2：点处理与邻域处理的分界}
\label{original:H2}
\textbf{原题：}下列属于空间域点处理的是？A.均值滤波；B.中值滤波；C.伽马变换；D.拉普拉斯锐化。

选C。判断只需看“计算某个输出像素时，需要读哪些输入”。伽马变换先把该像素灰度归一化为$u\in[0,1]$，再计算$v=u^\gamma$并缩放回目标灰度范围；这一过程只读取该像素本身的数值，因此属于点处理。均值要读取周围多个像素，中值要对邻域排序，拉普拉斯要比较中心与邻居，都属于邻域处理。

\subsection{H3：二阶差分的双响应}
\label{original:H3}
\textbf{原题：}对灰度阶跃最能产生“双响应”的是？A.一阶微分；B.二阶微分；C.均值滤波；D.中值滤波。

选B。考虑从0突然变到10的序列$[0,0,0,10,10,10]$。一阶差分比较相邻两项，在跳变处产生一个大的变化量。再对这一阶变化作差分，一侧由0升到较大值，另一侧由较大值降回0，于是得到一正一负的一对响应。二阶差分强调的是“变化量本身怎样变化”，而不只是灰度值是否高。

\subsection{H4：保边平滑的条件}
\label{original:H4}
\textbf{原题：}能同时“平滑噪声且较少模糊边缘”的是？A.均值；B.高斯低通；C.中值；D.理想低通。

按这组入门选项和脉冲噪声背景，选C。边缘两侧是不同的真实灰度，均值会把两侧混合为中间灰度，使边界变缓；中值选取邻域已有灰度，往往能在去掉孤立异常点的同时保留阶跃。不过，这并不是所有噪声类型下的统一排名。中值会影响细线和小目标；高斯滤波也有自己的适用场景。

\subsection{H5：8比特为什么是256级}
\label{original:H5}
\textbf{原题：}$8\,\mathrm{bit}$灰度图像最多能表示多少级灰度？

每个比特只有0和1两种状态。两个比特有$2\times2=4$种组合，八个比特有$2^8=256$种组合。因此最多表示256级灰度，常用整数编码为0至255。最大编码是255，而级数是256，因为0也占一个级别。

\subsection{H6：均衡化为什么使用累计分布}
\label{original:H6}
\textbf{原题：}直方图均衡的变换函数是什么？

基本构造是累积分布函数（Cumulative Distribution Function, CDF），再乘目标最大灰度并按规定取整。某个灰度的CDF表示“灰度不超过它的像素占了多大比例”。因此它把输入灰度的位置，改写为其在整幅图中的累计比例位置。连续理想模型中，这能把灰度概率重新分配；离散图像中多个输入等级可能合并，输出并不保证每个灰度都有相同像素数。

\subsection{H7：相同灰度为何看起来不同}
\label{original:H7}
\textbf{原题：}“同时对比”说明感知亮度依赖于什么？

依赖周围的背景亮度。同一块灰色放在黑背景中常显得较亮，放在白背景中常显得较暗。两块中心区域的数字可以完全相同，变化来自视觉系统的相对判断。因此图像中的一个灰度数值和人眼感到的亮度不能简单画等号。

\subsection{H8：混合邻接限制哪一种对角连接}
\label{original:H8}
\textbf{原题：}$m$邻接主要用于解决什么问题？

混合邻接（Mixed Adjacency）用更严格的条件处理对角像素之间的连接，避免8邻接引入的某些局部多路连接歧义。两个候选像素若是对角邻居，只有它们的共同4邻域中没有属于指定前景值集合的像素时，才允许这条对角连接。

这里检查的是\textbf{共同4邻域}，不是搜索整幅图中是否存在任意一条4连通路径；也不能把规则概括为“整张图的路径从此唯一”。

\subsection{H9：系数和为零会去掉什么}
\label{original:H9}
\textbf{原题：}纯锐化（高通）掩模的系数和是多少？

对题目所指的去直流导数型高通核，系数和为0。若邻域全是常数$c$，输出为$c\sum w$；当$\sum w=0$时，均匀背景输出0，只对变化产生响应。

但“最终锐化图像”通常还需要加回原图。加回之后的有效核可以具有系数和1，以保留均匀区域的灰度。必须区分\textbf{提取细节的高通核}与\textbf{输出增强图像的完整锐化核}。

\subsection{H10：CLAHE限制的是哪一个步骤}
\label{original:H10}
\textbf{原题：}CLAHE相比AHE的关键改进是什么？

限制局部直方图的峰高，即直方图裁剪（Histogram Clipping），并在局部映射之间进行插值以减轻块边界。插值也可以用于AHE，不能把它当成CLAHE独有的区别；名称中的“限制对比度”主要对应裁剪机制。过高的峰会使局部累计映射变化过快，从而把很小的灰度差异放得过大。限制峰高可以减轻噪声过度增强，但不能保证噪声完全不被放大。

\subsection{H11：频谱中心化是重排坐标}
\label{original:H11}
\textbf{原题：}频域滤波前把图像乘$(-1)^{x+y}$有什么作用？

在常见的偶数尺寸离散图像模型中，它把频谱的直流分量移到中心，使低通和高通滤波器便于按照“到中心的距离”描述。这里$(-1)^{x+y}$按棋盘格在$+1$与$-1$之间交替。它不是消除噪声，也不是改变截止频率本身。奇数尺寸与具体软件索引需要额外留意；不能把这一乘法与任意尺寸的\texttt{fftshift}无条件等同。

\subsection{H12：突变截止为何容易振铃}
\label{original:H12}
\textbf{原题：}下列最易产生振铃的是？A.高斯低通；B.巴特沃斯低通（$n=2$）；C.理想低通；D.均值滤波。

在本题的典型滤波器比较中，选C。理想低通在截止处从1突然跳到0，对应的空间响应有交替的旁瓣。旁瓣与图像阶跃边缘叠加，容易产生边缘附近的明暗振荡。高斯低通平滑衰减，其理想空间核没有这种正负旁瓣；巴特沃斯的阶数越高，通常越接近突变截止，振铃风险也随之增加。

\subsection{H13：周期条纹与成对频谱峰}
\label{original:H13}
\textbf{原题：}去除图像上周期性条纹干扰宜用什么方法？

当干扰集中在可识别的频率位置时，可以使用频域陷波滤波（Notch Filtering），或适当的带阻滤波（Band-reject Filtering）。实值图像的周期性正弦条纹通常对应关于频谱中心成对的峰，因此陷波位置应成对考虑。

选择抑制区域时必须兼顾真实图像细节。若真实栅栏、织物等结构恰好占据相近频率，滤波也可能把它们削弱。频谱中的亮点不是自动等于噪声，必须结合原图判断。

\section{原稿与详解版之间的对应及更正}
本书的内容范围来自提供的黑白复习讲义，原题编号用于对应，不代表所有题目都已由教师确认为考试原题。原稿未提供的题图不会被伪造为原图；相应图形仅作为明确标注的教学示意。

主要更正包括：视场与物方像素大小的数值一致性；拉普拉斯公式和模板的正负号；A2中二阶差分核的性质；中值滤波的非线性与保边条件；高提升滤波的亮度变化；离散直方图均衡与BBHE的适用条件；DFT线性卷积所需的补零尺寸；振幅0.5对应约$-6.02$ dB而非$-3$ dB；同态滤波的常数增益在原图域通常对应幂变换。

这些更正已经融入相应知识解释和解题过程，不再要求同时记忆相互冲突的两套答案。对于工程案例，保留题目提供的假设，但把“题设条件”“理想计算”和“实际设备是否满足”分开说明。

\clearpage
\section{原题索引}
本索引按原稿题号排列。页码指本书页脚编号，可直接点击跳转；正文把题目放在所需知识之后，以便从概念连续学到应用。
\begin{longtable}{p{13mm}p{119mm}p{17mm}}
\toprule 题号 & 主题 & 页码\\\midrule
\endhead
A1 & 四个掩模的逐项判读 & \pageref{original:A1}\\
A2 & 三个滤波核的性质和方向 & \pageref{original:A2}\\
B1 & 四类直方图的明暗与对比度 & \pageref{original:B1}\\
B2 & 夜间道路双峰直方图 & \pageref{original:B2}\\
B3 & 离散均衡为何不完全均匀 & \pageref{original:B3}\\
B4 & 频谱图判读 & \pageref{original:B4}\\
C1 & 暗人物与亮屏幕的直方图 & \pageref{original:C1}\\
C2 & 灰度变换曲线 & \pageref{original:C2}\\
D1 & 中值滤波五个问题 & \pageref{original:D1}\\
D2 & 均值与中值的选择 & \pageref{original:D2}\\
D3 & 滤波核系数和 & \pageref{original:D3}\\
D4 & 拉普拉斯与一、二阶差分 & \pageref{original:D4}\\
D5 & 非锐化掩模与高提升 & \pageref{original:D5}\\
D6 & 完整频域滤波过程 & \pageref{original:D6}\\
D7 & 同态滤波 & \pageref{original:D7}\\
E1 & 直方图均衡完整计算 & \pageref{original:E1}\\
E2 & 相机视场与物方分辨率 & \pageref{original:E2}\\
E3 & 平均去噪的方差与帧数 & \pageref{original:E3}\\
E4 & 韦伯比 & \pageref{original:E4}\\
E5 & 采样周期与采样间隔 & \pageref{original:E5}\\
E6 & 频域滤波数值及振铃 & \pageref{original:E6}\\
F1 & 增强与复原 & \pageref{original:F1}\\
F2 & 低级与高级处理 & \pageref{original:F2}\\
F3 & 视锥与视杆 & \pageref{original:F3}\\
F4 & 采样与量化 & \pageref{original:F4}\\
F5 & 空间与灰度分辨率 & \pageref{original:F5}\\
F6 & 邻域与混合邻接 & \pageref{original:F6}\\
F7 & 算术与逻辑运算 & \pageref{original:F7}\\
G1 & 报靶系统设计 & \pageref{original:G1}\\
G2 & 隧道裂缝检测 & \pageref{original:G2}\\
G3 & 车牌增强方案 & \pageref{original:G3}\\
H1 & 椒盐噪声与中值 & \pageref{original:H1}\\
H2 & 点处理 & \pageref{original:H2}\\
H3 & 二阶差分双响应 & \pageref{original:H3}\\
H4 & 保边平滑 & \pageref{original:H4}\\
H5 & 八比特与256级 & \pageref{original:H5}\\
H6 & 均衡变换与CDF & \pageref{original:H6}\\
H7 & 同时对比 & \pageref{original:H7}\\
H8 & 混合邻接 & \pageref{original:H8}\\
H9 & 高通核系数和 & \pageref{original:H9}\\
H10 & CLAHE裁剪 & \pageref{original:H10}\\
H11 & 频谱中心化 & \pageref{original:H11}\\
H12 & 理想低通振铃 & \pageref{original:H12}\\
H13 & 周期干扰的陷波 & \pageref{original:H13}\\
\bottomrule
\end{longtable}

\end{document}
